\documentclass[11pt,a4paper]{article}
\usepackage[T1]{fontenc}
\usepackage[utf8]{inputenc}
\usepackage{lmodern}
\usepackage[italian]{babel}
\usepackage[margin=23mm]{geometry}
\usepackage{microtype,graphicx,textcomp}
\usepackage{booktabs,xltabular,array,enumitem,caption,needspace}
\usepackage{xcolor,hyperref,listings,pdflscape,tikz}
\usepackage[italian,nameinlink,noabbrev]{cleveref}
\usetikzlibrary{arrows.meta,positioning,fit,calc,shapes.geometric}
\definecolor{ink}{HTML}{173946}
\definecolor{accent}{HTML}{9F432B}
\hypersetup{colorlinks=true,linkcolor=ink,urlcolor=accent,pdftitle={UninaFoodLab: progettazione e realizzazione},pdfauthor={Daniele De Blasio}}
\setlength{\parindent}{0pt}
\setlength{\parskip}{6pt plus 1pt minus 1pt}
\captionsetup{font=small,labelfont=bf,justification=raggedright,singlelinecheck=false,hypcap=false}
\setcounter{tocdepth}{2}
\newcommand{\sqlsource}{\path{src/main/resources/db/migration/V001__create_initial_schema.sql}}
\newcommand{\sqlcorrection}{\path{src/main/resources/db/migration/V002__protect_past_session_schedules.sql}}
\setlist{nosep}
\lstset{basicstyle=\ttfamily\footnotesize,breaklines=true,columns=fullflexible,keepspaces=true,showstringspaces=false,upquote=true,tabsize=4,frame=single,rulecolor=\color{black!15},numbers=left,numberstyle=\tiny\color{black!50},numbersep=5pt,aboveskip=10pt,belowskip=10pt,breakatwhitespace=true,captionpos=b}
\newcommand{\code}[1]{\texttt{\detokenize{#1}}}
\lstset{literate={à}{{\`a}}1 {è}{{\`e}}1 {é}{{\'e}}1 {ì}{{\`i}}1 {ù}{{\`u}}1 {È}{{\`E}}1 {°}{{\textdegree}}1 {·}{{\textperiodcentered}}1 {–}{{\textendash}}1 {—}{{\textemdash}}1}
\setlength{\emergencystretch}{3em}
\renewcommand{\lstlistingname}{Codice}
\crefname{lstlisting}{codice}{codici}
\Crefname{lstlisting}{Codice}{Codici}
\tikzset{
modelclass/.style={draw=ink,fill=white,text width=4.6cm,minimum height=1.8cm,inner sep=2.5mm,align=left,font=\small},
modelphysical/.style={draw=ink,fill=white,text width=4.9cm,minimum height=1.8cm,inner sep=2.5mm,align=left,font=\footnotesize},
modelentity/.style={draw=ink,fill=white,text width=3.6cm,minimum height=1.6cm,inner sep=2mm,align=left,font=\footnotesize},
modelrelation/.style={diamond,draw=accent,fill=white,aspect=2,text width=1.8cm,inner sep=1mm,align=center,font=\footnotesize},
modelassoc/.style={draw=ink,line width=.45pt}
}

\begin{document}
\crefname{listing}{codice}{codici}
\Crefname{listing}{Codice}{Codici}
\crefname{lstlisting}{codice}{codici}
\Crefname{lstlisting}{Codice}{Codici}

\hypersetup{pageanchor=false}\begin{titlepage}\centering\vspace*{26mm}
{\Huge\bfseries\color{ink}UninaFoodLab\par}\vspace{9mm}
{\Large Progettazione e realizzazione del sistema\par}\vspace{5mm}
{\large Relazione di progetto\par}\vfill
{\large Daniele De Blasio\par}\vspace{5mm}
\end{titlepage}
\pagenumbering{roman}\hypersetup{pageanchor=true}\tableofcontents\clearpage\pagenumbering{arabic}
\section{Introduzione}
UninaFoodLab organizza corsi di cucina, sessioni e ingredienti necessari. La relazione descrive la progettazione del sistema e della base di dati, motivando le scelte nei diversi passaggi. Il codice è disponibile nel repository \url{https://github.com/Daniele04-creator/UninaFoodLab}.

\section{Analisi dei requisiti e scelte di dominio}

Lo Chef organizza corsi con sessioni online o pratiche, associa ricette alle pratiche e consulta fabbisogni e report mensili. Il Partecipante si iscrive al corso e aderisce alle singole pratiche: questa distinzione permette di pianificare gli ingredienti in base alle conferme. Le comunicazioni sono consultabili nell'applicazione.

{\small \renewcommand{\arraystretch}{1.2}

\begin{xltabular}{\linewidth}{>{\hsize=0.240\hsize\linewidth=\hsize\raggedright\arraybackslash}X>{\hsize=0.780\hsize\linewidth=\hsize\raggedright\arraybackslash}X>{\hsize=1.980\hsize\linewidth=\hsize\raggedright\arraybackslash}X}
\caption{Analisi dei requisiti e scelte di dominio}\label{tab:dictionary-1}\\
\toprule

\textbf{ID} & \textbf{Requisito} & \textbf{Interpretazione operativa}\\\midrule\endfirsthead\toprule\textbf{ID} & \textbf{Requisito} & \textbf{Interpretazione operativa}\\\midrule\endhead
\midrule\multicolumn{3}{r}{\footnotesize Continua nella pagina successiva}\\\endfoot\bottomrule\endlastfoot


{\ttfamily\small R01} & Autenticazione e ruoli & Username/password; accessi alle operazioni autorizzati dal token.\\\addlinespace[3pt]

{\ttfamily\small R02} & Creazione dei corsi & Creazione e pubblicazione atomiche; inizialmente 2--30 sessioni.\\\addlinespace[3pt]

{\ttfamily\small R03} & Consultazione & Filtri per categoria e stato, ricerca e paginazione.\\\addlinespace[3pt]

{\ttfamily\small R04} & Programma delle sessioni & Modalità, durata, sede o URL HTTPS, capienza e scadenza delle adesioni.\\\addlinespace[3pt]

{\ttfamily\small R05} & Menu delle pratiche & Più ricette per sessione e copie delle dosi associate.\\\addlinespace[3pt]

{\ttfamily\small R06} & Iscrizioni e adesioni & Iscrizione al corso distinta dalla conferma della singola pratica.\\\addlinespace[3pt]

{\ttfamily\small R07} & Fabbisogni & Dosi della copia moltiplicate per conferme valide; lista della spesa stampabile.\\\addlinespace[3pt]

{\ttfamily\small R08} & Variazioni e comunicazioni & Programma correggibile e comunicazioni destinate agli iscritti.\\\addlinespace[3pt]

{\ttfamily\small R09} & Report mensile & Corsi distinti, online/pratiche, media/minimo/massimo delle ricette; grafici.\\\addlinespace[3pt]

{\ttfamily\small R10} & Persistenza e integrità & Vincoli, trigger e funzioni SQL; migrazioni versionate.\\*



\end{xltabular}}

\subsection{Assunzioni e regole di dominio}

Una sessione non annullata è considerata svolta alla fine prevista, nel fuso Europe/Rome; le ricette del menu sono assunte realizzate. Il report conta corsi distinti e sessioni svolte nel mese dell'appuntamento. Una pratica senza ricette contribuisce con zero; in assenza di pratiche le statistiche sulle ricette sono assenti.

Uno spostamento dell'appuntamento richiede una nuova adesione; la sola correzione di sede, link o scadenza conserva le conferme. Una reiscrizione non riattiva risposte del periodo precedente. I fabbisogni usano le copie del menu e le conferme correnti, senza margine aggiuntivo.

Il calendario consente modifiche alle sole sessioni future e programmate di corsi pubblicati. Una sessione già iniziata o conclusa conserva il proprio appuntamento, anche quando sede e link non vengono modificati.

Le dosi sono riferite a una persona e ammettono al massimo tre decimali senza arrotondamento implicito. I limiti di 2--30 sessioni iniziali e frequenza di 1--31 giorni sono protezioni operative del form di creazione.

\clearpage\begin{landscape}
\section{Modello del dominio}
Il modello rappresenta concetti e attributi del dominio. Le molteplicità sono ai due estremi: ciascun valore conta le istanze della classe vicina collegate a una singola istanza opposta. Le classi ripetute mantengono la stessa identità; / indica un valore derivato.

Account e profili sono oggetti distinti, collegati da associazioni. La registrazione crea il profilo richiesto dal ruolo; trigger e routine impediscono profili incompatibili.


\subsection{UML concettuale: identità, corsi e adesioni}
\begin{center}\begin{minipage}{\linewidth}
\centering
\begin{tikzpicture}[x=1cm,y=1cm]
\node[modelclass] (Account) at (0,0) {\textbf{Account}\\\rule{\linewidth}{0.3pt}\\username, credenziale\\ruolo, attivo};
\node[modelclass] (Chef) at (8,0) {\textbf{Chef}\\\rule{\linewidth}{0.3pt}\\nome visualizzato};
\node[modelclass] (Categoria) at (16,0) {\textbf{Categoria}\\\rule{\linewidth}{0.3pt}\\nome, descrizione};
\node[modelclass] (Partecipante) at (0,-3.3) {\textbf{Partecipante}\\\rule{\linewidth}{0.3pt}\\nome visualizzato};
\node[modelclass] (Corso) at (8,-3.3) {\textbf{Corso}\\\rule{\linewidth}{0.3pt}\\titolo, descrizione\\inizio, frequenza, n. sessioni\\stato, margine ingredienti};
\node[modelclass] (Sessione) at (16,-3.3) {\textbf{Sessione}\\\rule{\linewidth}{0.3pt}\\numero, inizio, durata, modalità\\sede oppure link, capienza\\chiusura adesioni, stato\\/svolgimento};
\node[modelclass] (Iscrizione) at (0,-6.6) {\textbf{Iscrizione}\\\rule{\linewidth}{0.3pt}\\stato, data iscrizione\\data di ritiro opzionale};
\node[modelclass] (Adesione) at (16,-6.6) {\textbf{Adesione}\\\rule{\linewidth}{0.3pt}\\stato della risposta\\data della risposta};
\draw[modelassoc] (Account.east) -- node[pos=.15,fill=white,inner sep=1pt,font=\footnotesize]{1} node[pos=.85,fill=white,inner sep=1pt,font=\footnotesize]{0..1} (Chef.west);
\draw[modelassoc] (Account.south) -- node[pos=.15,fill=white,inner sep=1pt,font=\footnotesize]{1} node[pos=.85,fill=white,inner sep=1pt,font=\footnotesize]{0..1} (Partecipante.north);
\draw[modelassoc] (Chef.south) -- node[pos=.15,fill=white,inner sep=1pt,font=\footnotesize]{1} node[pos=.85,fill=white,inner sep=1pt,font=\footnotesize]{0..*} (Corso.north);
\draw[modelassoc] (Categoria.south west) -- node[pos=.15,fill=white,inner sep=1pt,font=\footnotesize]{1} node[pos=.85,fill=white,inner sep=1pt,font=\footnotesize]{0..*} (Corso.north east);
\draw[modelassoc] (Corso.east) -- node[pos=.15,fill=white,inner sep=1pt,font=\footnotesize]{1} node[pos=.85,fill=white,inner sep=1pt,font=\footnotesize]{0..*} (Sessione.west);
\draw[modelassoc] (Partecipante.south) -- node[pos=.15,fill=white,inner sep=1pt,font=\footnotesize]{1} node[pos=.85,fill=white,inner sep=1pt,font=\footnotesize]{0..*} (Iscrizione.north);
\draw[modelassoc] (Corso.south west) -- node[pos=.15,fill=white,inner sep=1pt,font=\footnotesize]{1} node[pos=.85,fill=white,inner sep=1pt,font=\footnotesize]{0..*} (Iscrizione.north east);
\draw[modelassoc] (Iscrizione.east) -- node[pos=.15,fill=white,inner sep=1pt,font=\footnotesize]{1} node[pos=.85,fill=white,inner sep=1pt,font=\footnotesize]{0..*} (Adesione.west);
\draw[modelassoc] (Sessione.south) -- node[pos=.15,fill=white,inner sep=1pt,font=\footnotesize]{1} node[pos=.85,fill=white,inner sep=1pt,font=\footnotesize]{0..*} (Adesione.north);
\end{tikzpicture}
\captionof{figure}{UML concettuale: identità, corsi e adesioni}\label{fig:diagram-1}
\end{minipage}\end{center}
{\small Iscrizione collega Partecipante e Corso. Adesione collega quella Iscrizione a una Sessione dello stesso corso, distinguendo le risposte dei diversi periodi. La modalità condiziona gli attributi della Sessione; limiti di creazione e controlli tra righe sono nella \cref{sec:constraints}.}
\end{landscape}

\clearpage\begin{landscape}
\subsection{UML concettuale: ricettario, menu e copie}
\begin{center}\begin{minipage}{\linewidth}
\centering
\begin{tikzpicture}[x=1cm,y=1cm]
\node[modelclass] (Categoria) at (0,0) {\textbf{Categoria}\\\rule{\linewidth}{0.3pt}\\nome, descrizione};
\node[modelclass] (Ricetta) at (8,0) {\textbf{Ricetta}\\\rule{\linewidth}{0.3pt}\\nome, descrizione\\dosi per una persona};
\node[modelclass] (Ingrediente) at (16,0) {\textbf{Ingrediente}\\\rule{\linewidth}{0.3pt}\\nome, unità};
\node[modelclass] (Sessione) at (0,-4.3) {\textbf{Sessione}\\\rule{\linewidth}{0.3pt}\\numero, inizio, durata, modalità\\sede oppure link, capienza\\chiusura adesioni, stato\\/svolgimento};
\node[modelclass] (RicettaSessione) at (8,-4.3) {\textbf{RicettaSessione}\\\rule{\linewidth}{0.3pt}\\porzioni per partecipante = 1};
\node[modelclass] (DoseRicetta) at (16,-4.3) {\textbf{DoseRicetta}\\\rule{\linewidth}{0.3pt}\\quantità per una persona};
\node[modelclass] (CopiaMenu) at (8,-8.6) {\textbf{CopiaMenu}\\\rule{\linewidth}{0.3pt}\\nome, descrizione e dosi copiate\\acquisizione, provenienza};
\draw[modelassoc] (Categoria.east) -- node[pos=.15,fill=white,inner sep=1pt,font=\footnotesize]{1} node[pos=.85,fill=white,inner sep=1pt,font=\footnotesize]{0..*} (Ricetta.west);
\draw[modelassoc] (Ricetta.south east) -- node[pos=.15,fill=white,inner sep=1pt,font=\footnotesize]{1} node[pos=.85,fill=white,inner sep=1pt,font=\footnotesize]{0..*} (DoseRicetta.north west);
\draw[modelassoc] (Ingrediente.south) -- node[pos=.15,fill=white,inner sep=1pt,font=\footnotesize]{1} node[pos=.85,fill=white,inner sep=1pt,font=\footnotesize]{0..*} (DoseRicetta.north);
\draw[modelassoc] (Ricetta.south) -- node[pos=.15,fill=white,inner sep=1pt,font=\footnotesize]{1} node[pos=.85,fill=white,inner sep=1pt,font=\footnotesize]{0..*} (RicettaSessione.north);
\draw[modelassoc] (Sessione.east) -- node[pos=.15,fill=white,inner sep=1pt,font=\footnotesize]{1} node[pos=.85,fill=white,inner sep=1pt,font=\footnotesize]{0..*} (RicettaSessione.west);
\draw[modelassoc] (RicettaSessione.south) -- node[pos=.15,fill=white,inner sep=1pt,font=\footnotesize]{1} node[pos=.85,fill=white,inner sep=1pt,font=\footnotesize]{0..1} (CopiaMenu.north);
\end{tikzpicture}
\captionof{figure}{UML concettuale: ricettario, menu e copie}\label{fig:diagram-2}
\end{minipage}\end{center}
{\small DoseRicetta rende esplicita la quantità del legame tra Ricetta e Ingrediente. RicettaSessione permette più ricette nella stessa pratica; CopiaMenu conserva le dosi e il contesto acquisiti, così le modifiche al ricettario non cambiano una preparazione già associata. Il ricettario è condiviso e non registra un proprietario Chef. La molteplicità 0..1 della copia esprime la struttura; l'associazione ordinaria la crea automaticamente.}
\end{landscape}

\clearpage\begin{landscape}
\subsection{UML concettuale: comunicazioni e lettura}
\begin{center}\begin{minipage}{\linewidth}
\centering
\begin{tikzpicture}[x=1cm,y=1cm]
\node[modelclass] (Chef) at (0,0) {\textbf{Chef}\\\rule{\linewidth}{0.3pt}\\nome visualizzato};
\node[modelclass] (Comunicazione) at (8,0) {\textbf{Comunicazione}\\\rule{\linewidth}{0.3pt}\\oggetto, testo, motivo\\data di creazione};
\node[modelclass] (Corso) at (16,0) {\textbf{Corso}\\\rule{\linewidth}{0.3pt}\\titolo, descrizione\\inizio, frequenza, n. sessioni\\stato, margine ingredienti};
\node[modelclass] (Partecipante) at (0,-4.3) {\textbf{Partecipante}\\\rule{\linewidth}{0.3pt}\\nome visualizzato};
\node[modelclass] (Destinatario) at (8,-4.3) {\textbf{Destinatario}\\\rule{\linewidth}{0.3pt}\\data di lettura opzionale};
\draw[modelassoc] (Chef.east) -- node[pos=.15,fill=white,inner sep=1pt,font=\footnotesize]{1} node[pos=.85,fill=white,inner sep=1pt,font=\footnotesize]{0..*} (Comunicazione.west);
\draw[modelassoc] (Corso.west) -- node[pos=.15,fill=white,inner sep=1pt,font=\footnotesize]{0..1} node[pos=.85,fill=white,inner sep=1pt,font=\footnotesize]{0..*} (Comunicazione.east);
\draw[modelassoc] (Comunicazione.south) -- node[pos=.15,fill=white,inner sep=1pt,font=\footnotesize]{1} node[pos=.85,fill=white,inner sep=1pt,font=\footnotesize]{0..*} (Destinatario.north);
\draw[modelassoc] (Partecipante.east) -- node[pos=.15,fill=white,inner sep=1pt,font=\footnotesize]{1} node[pos=.85,fill=white,inner sep=1pt,font=\footnotesize]{0..*} (Destinatario.west);
\end{tikzpicture}
\captionof{figure}{UML concettuale: comunicazioni e lettura}\label{fig:conceptual-messages}
\end{minipage}\end{center}
{\small Comunicazione ha sempre un autore Chef. Il Corso è opzionale perché il messaggio può riguardare tutti i corsi dello chef. Destinatario identifica ogni coppia messaggio/partecipante e conserva la lettura senza modificare il contenuto comune. La registrazione del messaggio non rappresenta un invio email o SMS.}
\subsection{Modello entità-relazione}
Il modello ER esplicita gli identificatori e gli attributi. I rettangoli contengono le entità; gli identificatori sono sottolineati. I rombi rappresentano relazioni e possono contenere attributi propri. Le cardinalità minime e massime sono poste vicino all'entità e indicano quante volte ciascuna sua istanza partecipa alla relazione: 0..N ammette anche nessun collegamento.

Iscrizione è rappresentata come entità perché possiede una propria identità ed è richiamata da Adesione. Anche RicettaSessione viene resa esplicita, identificata dalla coppia Sessione/Ricetta, perché la copia del menu dipende da quella coppia. Dose, Adesione e Destinatario rimangono relazioni con attributi; nella ristrutturazione diventano tabelle associative. Le entità ripetute collegano le tre tavole.

\end{landscape}

\clearpage\begin{landscape}
\subsubsection{Identità, corsi, iscrizioni e adesioni}
\begin{center}\begin{minipage}{\linewidth}
\centering
\begin{tikzpicture}[x=1cm,y=1cm]
\node[modelentity,text width=3.6cm] (Account) at (0,0) {\textbf{Account}\\\rule{\linewidth}{0.3pt}\\\underline{id account}\\username, credenziale\\ruolo, attivo};
\node[modelentity,text width=3.6cm] (Chef) at (9,0) {\textbf{Chef}\\\rule{\linewidth}{0.3pt}\\\underline{id chef}\\nome visualizzato};
\node[modelentity,text width=3.6cm] (Categoria) at (18,0) {\textbf{Categoria}\\\rule{\linewidth}{0.3pt}\\\underline{id categoria}\\nome, descrizione};
\node[modelentity,text width=3.6cm] (Partecipante) at (0,-5) {\textbf{Partecipante}\\\rule{\linewidth}{0.3pt}\\\underline{id partecipante}\\nome visualizzato};
\node[modelentity,text width=3.6cm] (Corso) at (9,-5) {\textbf{Corso}\\\rule{\linewidth}{0.3pt}\\\underline{id corso}\\titolo, descrizione\\inizio, frequenza, n. sessioni\\stato, margine ingredienti};
\node[modelentity,text width=3.6cm] (Sessione) at (18,-5) {\textbf{Sessione}\\\rule{\linewidth}{0.3pt}\\\underline{id sessione}\\numero, inizio, durata, modalità\\sede oppure link, capienza\\chiusura adesioni, stato\\/svolgimento};
\node[modelentity,text width=3.6cm] (Iscrizione) at (0,-10) {\textbf{Iscrizione}\\\rule{\linewidth}{0.3pt}\\\underline{id iscrizione}\\stato, data iscrizione\\data di ritiro opzionale};
\node[modelrelation] (rAC) at (4.5,0) {profilo Chef};
\node[modelrelation] (rAP) at (0,-2.5) {profilo\\Partecipante};
\node[modelrelation] (rCC) at (9,-2.5) {organizza};
\node[modelrelation] (rKC) at (13.5,-2.5) {classifica};
\node[modelrelation] (rCS) at (13.5,-5) {contiene};
\node[modelrelation] (rPI) at (0,-7.5) {effettua};
\node[modelrelation] (rCI) at (4.5,-7.5) {riguarda};
\node[modelrelation] (rIS) at (9,-10) {adesione\\{\footnotesize stato, data risposta}};
\draw[modelassoc] (Account.east) -- node[pos=.25,fill=white,inner sep=1pt,font=\footnotesize]{0..1} (rAC.west);
\draw[modelassoc] (Chef.west) -- node[pos=.25,fill=white,inner sep=1pt,font=\footnotesize]{1} (rAC.east);
\draw[modelassoc] (Account.south) -- node[pos=.25,fill=white,inner sep=1pt,font=\footnotesize]{0..1} (rAP.north);
\draw[modelassoc] (Partecipante.north) -- node[pos=.25,fill=white,inner sep=1pt,font=\footnotesize]{1} (rAP.south);
\draw[modelassoc] (Chef.south) -- node[pos=.25,fill=white,inner sep=1pt,font=\footnotesize]{0..N} (rCC.north);
\draw[modelassoc] (Corso.north) -- node[pos=.25,fill=white,inner sep=1pt,font=\footnotesize]{1} (rCC.south);
\draw[modelassoc] (Categoria.south west) -- node[pos=.25,fill=white,inner sep=1pt,font=\footnotesize]{0..N} (rKC.north east);
\draw[modelassoc] (Corso.north east) -- node[pos=.25,fill=white,inner sep=1pt,font=\footnotesize]{1} (rKC.south west);
\draw[modelassoc] (Corso.east) -- node[pos=.25,fill=white,inner sep=1pt,font=\footnotesize]{0..N} (rCS.west);
\draw[modelassoc] (Sessione.west) -- node[pos=.25,fill=white,inner sep=1pt,font=\footnotesize]{1} (rCS.east);
\draw[modelassoc] (Partecipante.south) -- node[pos=.25,fill=white,inner sep=1pt,font=\footnotesize]{0..N} (rPI.north);
\draw[modelassoc] (Iscrizione.north) -- node[pos=.25,fill=white,inner sep=1pt,font=\footnotesize]{1} (rPI.south);
\draw[modelassoc] (Corso.south west) -- node[pos=.25,fill=white,inner sep=1pt,font=\footnotesize]{0..N} (rCI.north east);
\draw[modelassoc] (Iscrizione.north east) -- node[pos=.25,fill=white,inner sep=1pt,font=\footnotesize]{1} (rCI.south west);
\draw[modelassoc] (Iscrizione.east) -- node[pos=.25,fill=white,inner sep=1pt,font=\footnotesize]{0..N} (rIS.west);
\draw[modelassoc] (Sessione.south west) -- node[pos=.25,fill=white,inner sep=1pt,font=\footnotesize]{0..N} (rIS.north east);
\end{tikzpicture}
\captionof{figure}{ER: identità, corsi, iscrizioni e adesioni}\label{fig:diagram-3}
\end{minipage}\end{center}
{\small Iscrizione conserva stato e periodo ed è univoca per Partecipante/Corso. La relazione Adesione è identificata dalla coppia Iscrizione/Sessione; una conferma è ammessa solo per una pratica dello stesso corso. I profili sono alternativi per ruolo. Gli identificatori dei singoli oggetti diventano UUID nello schema logico.}
\end{landscape}

\clearpage\begin{landscape}
\subsubsection{Ricette, ingredienti, menu e copie}
\begin{center}\begin{minipage}{\linewidth}
\centering
\begin{tikzpicture}[x=1cm,y=1cm]
\node[modelentity,text width=3.6cm] (Categoria) at (0,0) {\textbf{Categoria}\\\rule{\linewidth}{0.3pt}\\\underline{id categoria}\\nome, descrizione};
\node[modelentity,text width=3.6cm] (Ricetta) at (8,0) {\textbf{Ricetta}\\\rule{\linewidth}{0.3pt}\\\underline{id ricetta}\\nome, descrizione\\dosi per una persona};
\node[modelentity,text width=3.6cm] (Ingrediente) at (16,0) {\textbf{Ingrediente}\\\rule{\linewidth}{0.3pt}\\\underline{id ingrediente}\\nome, unità};
\node[modelentity,text width=3.6cm] (Sessione) at (0,-5) {\textbf{Sessione}\\\rule{\linewidth}{0.3pt}\\\underline{id sessione}\\numero, inizio, durata, modalità\\sede oppure link, capienza\\chiusura adesioni, stato\\/svolgimento};
\node[modelentity,text width=3.6cm] (RicettaSessione) at (8,-5) {\textbf{RicettaSessione}\\\rule{\linewidth}{0.3pt}\\\underline{(Sessione, Ricetta)}\\porzioni per partecipante = 1};
\node[modelentity,text width=3.6cm] (CopiaMenu) at (16,-5) {\textbf{CopiaMenu}\\\rule{\linewidth}{0.3pt}\\\underline{(Sessione, Ricetta)}\\nome, descrizione e dosi copiate\\acquisizione, provenienza};
\node[modelrelation] (rKR) at (4,0) {classifica};
\node[modelrelation] (rRI) at (12,0) {dose\\{\footnotesize quantità}};
\node[modelrelation] (rSR) at (4,-5) {comprende};
\node[modelrelation] (rRR) at (8,-2.5) {usa};
\node[modelrelation] (rMC) at (12,-5) {acquisisce};
\draw[modelassoc] (Categoria.east) -- node[pos=.25,fill=white,inner sep=1pt,font=\footnotesize]{0..N} (rKR.west);
\draw[modelassoc] (Ricetta.west) -- node[pos=.25,fill=white,inner sep=1pt,font=\footnotesize]{1} (rKR.east);
\draw[modelassoc] (Ricetta.east) -- node[pos=.25,fill=white,inner sep=1pt,font=\footnotesize]{0..N} (rRI.west);
\draw[modelassoc] (Ingrediente.west) -- node[pos=.25,fill=white,inner sep=1pt,font=\footnotesize]{0..N} (rRI.east);
\draw[modelassoc] (Sessione.east) -- node[pos=.25,fill=white,inner sep=1pt,font=\footnotesize]{0..N} (rSR.west);
\draw[modelassoc] (RicettaSessione.west) -- node[pos=.25,fill=white,inner sep=1pt,font=\footnotesize]{1} (rSR.east);
\draw[modelassoc] (Ricetta.south) -- node[pos=.25,fill=white,inner sep=1pt,font=\footnotesize]{0..N} (rRR.north);
\draw[modelassoc] (RicettaSessione.north) -- node[pos=.25,fill=white,inner sep=1pt,font=\footnotesize]{1} (rRR.south);
\draw[modelassoc] (RicettaSessione.east) -- node[pos=.25,fill=white,inner sep=1pt,font=\footnotesize]{0..1} (rMC.west);
\draw[modelassoc] (CopiaMenu.west) -- node[pos=.25,fill=white,inner sep=1pt,font=\footnotesize]{1} (rMC.east);
\end{tikzpicture}
\captionof{figure}{ER: ricette, ingredienti, menu e copie}\label{fig:diagram-4}
\end{minipage}\end{center}
{\small Dose è identificata dalla coppia Ricetta/Ingrediente e contiene la quantità per una persona. RicettaSessione conserva le porzioni per partecipante; la copia usa lo stesso identificatore composto. Una Sessione online non può avere un menu. Categoria è la stessa entità che classifica i corsi.}
\end{landscape}

\clearpage\begin{landscape}
\subsubsection{Comunicazioni e destinatari}
\begin{center}\begin{minipage}{\linewidth}
\centering
\begin{tikzpicture}[x=1cm,y=1cm]
\node[modelentity,text width=3.6cm] (Chef) at (0,0) {\textbf{Chef}\\\rule{\linewidth}{0.3pt}\\\underline{id chef}\\nome visualizzato};
\node[modelentity,text width=3.6cm] (Comunicazione) at (8,0) {\textbf{Comunicazione}\\\rule{\linewidth}{0.3pt}\\\underline{id comunicazione}\\oggetto, testo, motivo\\data di creazione};
\node[modelentity,text width=3.6cm] (Corso) at (16,0) {\textbf{Corso}\\\rule{\linewidth}{0.3pt}\\\underline{id corso}\\titolo, descrizione\\inizio, frequenza, n. sessioni\\stato, margine ingredienti};
\node[modelentity,text width=3.6cm] (Partecipante) at (8,-5) {\textbf{Partecipante}\\\rule{\linewidth}{0.3pt}\\\underline{id partecipante}\\nome visualizzato};
\node[modelrelation] (rCN) at (4,0) {scrive};
\node[modelrelation] (rNC) at (12,0) {riguarda};
\node[modelrelation] (rNP) at (8,-2.5) {destinatario\\{\footnotesize data lettura opzionale}};
\draw[modelassoc] (Chef.east) -- node[pos=.25,fill=white,inner sep=1pt,font=\footnotesize]{0..N} (rCN.west);
\draw[modelassoc] (Comunicazione.west) -- node[pos=.25,fill=white,inner sep=1pt,font=\footnotesize]{1} (rCN.east);
\draw[modelassoc] (Comunicazione.east) -- node[pos=.25,fill=white,inner sep=1pt,font=\footnotesize]{0..1} (rNC.west);
\draw[modelassoc] (Corso.west) -- node[pos=.25,fill=white,inner sep=1pt,font=\footnotesize]{0..N} (rNC.east);
\draw[modelassoc] (Comunicazione.south) -- node[pos=.25,fill=white,inner sep=1pt,font=\footnotesize]{0..N} (rNP.north);
\draw[modelassoc] (Partecipante.north) -- node[pos=.25,fill=white,inner sep=1pt,font=\footnotesize]{0..N} (rNP.south);
\end{tikzpicture}
\captionof{figure}{ER: comunicazioni e destinatari}\label{fig:diagram-5}
\end{minipage}\end{center}
{\small Destinatario è identificato dalla coppia Comunicazione/Partecipante e registra una data di lettura opzionale. Il Corso può essere assente in un messaggio globale; l'autore Chef è sempre presente. La materializzazione dei destinatari conserva a chi è stata indirizzata la comunicazione.}
\end{landscape}

\clearpage\begin{landscape}
\section{Ristrutturazione e schema relazionale}
Account e profili diventano tabelle distinte con FK univoca, evitando credenziali duplicate. Le relazioni molti-a-molti diventano associative: Dose, Adesione e Destinatario usano una coppia di FK come PK; Iscrizione mantiene un identificatore autonomo perché viene richiamata dalle adesioni.

Ricette, ingredienti e dosi correnti sono normalizzati. CopiaMenu conserva invece il contesto in JSONB: questa ridondanza impedisce modifiche retroattive della preparazione. La sua FK composta punta all'intera coppia Sessione/Ricetta.

Le cinque tavole comprendono 22 tabelle e 24 FK. I nodi di raccordo sono ripetuti; gli attributi completi sono nella \cref{sec:logical-schema}. PK indica la chiave primaria, FK un riferimento, U l'unicità di un attributo e ? la nullabilità. Le molteplicità descrivono la struttura; trigger e routine aggiungono le regole tra righe.


\subsection{UML ristrutturato: account e profili}
\begin{center}\begin{minipage}{\linewidth}
\centering
\begin{tikzpicture}[x=1cm,y=1cm]
\node[modelphysical] (chef_profile) at (0,0) {\textbf{chef\_\allowbreak{}profile}\\\rule{\linewidth}{0.3pt}\\id: UUID [PK]\\account\_\allowbreak{}id: UUID [FK,U]\\display\_\allowbreak{}name: VARCHAR};
\node[modelphysical] (account) at (8,0) {\textbf{account}\\\rule{\linewidth}{0.3pt}\\id: UUID [PK]\\username: VARCHAR [U]\\role: VARCHAR\\active: BOOLEAN};
\node[modelphysical] (participant_profile) at (16,0) {\textbf{participant\_\allowbreak{}profile}\\\rule{\linewidth}{0.3pt}\\id: UUID [PK]\\account\_\allowbreak{}id: UUID [FK,U]\\display\_\allowbreak{}name: VARCHAR};
\node[modelphysical] (auth_session) at (0,-3.5) {\textbf{auth\_\allowbreak{}session}\\\rule{\linewidth}{0.3pt}\\token\_\allowbreak{}hash: BYTEA [PK]\\account\_\allowbreak{}id: UUID [FK]\\expires\_\allowbreak{}at: TIMESTAMPTZ};
\node[modelphysical] (account_password_reset) at (8,-3.5) {\textbf{account\_\allowbreak{}password\_\allowbreak{}reset}\\\rule{\linewidth}{0.3pt}\\token\_\allowbreak{}hash: BYTEA [PK]\\account\_\allowbreak{}id: UUID [FK]\\expires\_\allowbreak{}at: TIMESTAMPTZ\\used\_\allowbreak{}at: TIMESTAMPTZ [?]};
\node[modelphysical] (chef_invitation) at (16,-3.5) {\textbf{chef\_\allowbreak{}invitation}\\\rule{\linewidth}{0.3pt}\\token\_\allowbreak{}hash: BYTEA [PK]\\expires\_\allowbreak{}at: TIMESTAMPTZ\\used\_\allowbreak{}at: TIMESTAMPTZ [?]};
\draw[modelassoc] (account.west) -- node[pos=.15,fill=white,inner sep=1pt,font=\footnotesize]{1} node[pos=.85,fill=white,inner sep=1pt,font=\footnotesize]{0..1} (chef_profile.east);
\draw[modelassoc] (account.east) -- node[pos=.15,fill=white,inner sep=1pt,font=\footnotesize]{1} node[pos=.85,fill=white,inner sep=1pt,font=\footnotesize]{0..1} (participant_profile.west);
\draw[modelassoc] (account.south west) -- node[pos=.15,fill=white,inner sep=1pt,font=\footnotesize]{1} node[pos=.85,fill=white,inner sep=1pt,font=\footnotesize]{0..*} (auth_session.north east);
\draw[modelassoc] (account.south) -- node[pos=.15,fill=white,inner sep=1pt,font=\footnotesize]{1} node[pos=.85,fill=white,inner sep=1pt,font=\footnotesize]{0..*} (account_password_reset.north);
\end{tikzpicture}
\captionof{figure}{UML ristrutturato: account e profili}\label{fig:restructured-accounts}
\end{minipage}\end{center}
{\small account\_id UNIQUE limita ogni tipo a un profilo per account. auth\_session supporta gli accessi correnti; account\_password\_reset e chef\_invitation conservano dati storici. chef\_invitation non ha FK. Routine e trigger controllano ruolo e creazione dei profili.}
\end{landscape}

\clearpage\begin{landscape}
\subsection{UML ristrutturato: corsi, iscrizioni e adesioni}
\begin{center}\begin{minipage}{\linewidth}
\centering
\begin{tikzpicture}[x=1cm,y=1cm]
\node[modelphysical] (participant_profile) at (0,0) {\textbf{participant\_\allowbreak{}profile}\\\rule{\linewidth}{0.3pt}\\id: UUID [PK]\\account\_\allowbreak{}id: UUID [FK,U]\\display\_\allowbreak{}name: VARCHAR};
\node[modelphysical] (chef_profile) at (8,0) {\textbf{chef\_\allowbreak{}profile}\\\rule{\linewidth}{0.3pt}\\id: UUID [PK]\\account\_\allowbreak{}id: UUID [FK,U]\\display\_\allowbreak{}name: VARCHAR};
\node[modelphysical] (category) at (16,0) {\textbf{category}\\\rule{\linewidth}{0.3pt}\\id: UUID [PK]\\name: VARCHAR [U]};
\node[modelphysical] (enrollment) at (0,-5) {\textbf{enrollment}\\\rule{\linewidth}{0.3pt}\\id: UUID [PK]\\participant\_\allowbreak{}profile\_\allowbreak{}id: UUID [FK]\\course\_\allowbreak{}id: UUID [FK]\\status: VARCHAR\\enrolled\_\allowbreak{}at: TIMESTAMPTZ};
\node[modelphysical] (course) at (8,-5) {\textbf{course}\\\rule{\linewidth}{0.3pt}\\id: UUID [PK]\\chef\_\allowbreak{}profile\_\allowbreak{}id: UUID [FK]\\category\_\allowbreak{}id: UUID [FK]\\title: VARCHAR\\start\_\allowbreak{}date: DATE\\status: VARCHAR};
\node[modelphysical] (course_session) at (16,-5) {\textbf{course\_\allowbreak{}session}\\\rule{\linewidth}{0.3pt}\\id: UUID [PK]\\course\_\allowbreak{}id: UUID [FK]\\sequence\_\allowbreak{}number: INTEGER\\scheduled\_\allowbreak{}at: TIMESTAMP\\mode: VARCHAR\\duration\_\allowbreak{}minutes: INTEGER};
\node[modelphysical] (practical_attendance) at (0,-10) {\textbf{practical\_\allowbreak{}attendance}\\\rule{\linewidth}{0.3pt}\\enrollment\_\allowbreak{}id: UUID [PK,FK]\\session\_\allowbreak{}id: UUID [PK,FK]\\status: VARCHAR\\updated\_\allowbreak{}at: TIMESTAMPTZ};
\draw[modelassoc] (participant_profile.south) -- node[pos=.15,fill=white,inner sep=1pt,font=\footnotesize]{1} node[pos=.85,fill=white,inner sep=1pt,font=\footnotesize]{0..*} (enrollment.north);
\draw[modelassoc] (chef_profile.south) -- node[pos=.15,fill=white,inner sep=1pt,font=\footnotesize]{1} node[pos=.85,fill=white,inner sep=1pt,font=\footnotesize]{0..*} (course.north);
\draw[modelassoc] (category.south west) -- node[pos=.15,fill=white,inner sep=1pt,font=\footnotesize]{1} node[pos=.85,fill=white,inner sep=1pt,font=\footnotesize]{0..*} (course.north east);
\draw[modelassoc] (course.west) -- node[pos=.15,fill=white,inner sep=1pt,font=\footnotesize]{1} node[pos=.85,fill=white,inner sep=1pt,font=\footnotesize]{0..*} (enrollment.east);
\draw[modelassoc] (course.east) -- node[pos=.15,fill=white,inner sep=1pt,font=\footnotesize]{1} node[pos=.85,fill=white,inner sep=1pt,font=\footnotesize]{0..*} (course_session.west);
\draw[modelassoc] (enrollment.south) -- node[pos=.15,fill=white,inner sep=1pt,font=\footnotesize]{1} node[pos=.85,fill=white,inner sep=1pt,font=\footnotesize]{0..*} (practical_attendance.north);
\draw[modelassoc] (course_session.south west) -- node[pos=.15,fill=white,inner sep=1pt,font=\footnotesize]{1} node[pos=.85,fill=white,inner sep=1pt,font=\footnotesize]{0..*} (practical_attendance.east);
\end{tikzpicture}
\captionof{figure}{UML ristrutturato: corsi, iscrizioni e adesioni}\label{fig:diagram-6}
\end{minipage}\end{center}
{\small La coppia partecipante/corso di enrollment è UNIQUE. Anche corso/numero della sessione è UNIQUE. La PK di practical\_attendance è la coppia enrollment\_id/session\_id; il trigger verifica corso, iscrizione attiva e capienza. Le FK da sole consentono zero sessioni, mentre la creazione applicativa ne richiede inizialmente 2--30.}
\end{landscape}

\clearpage\begin{landscape}
\subsection{UML ristrutturato: ricettario, menu e copie}
\begin{center}\begin{minipage}{\linewidth}
\centering
\begin{tikzpicture}[x=1cm,y=1cm]
\node[modelphysical] (category) at (0,0) {\textbf{category}\\\rule{\linewidth}{0.3pt}\\id: UUID [PK]\\name: VARCHAR [U]};
\node[modelphysical] (recipe) at (8,0) {\textbf{recipe}\\\rule{\linewidth}{0.3pt}\\id: UUID [PK]\\category\_\allowbreak{}id: UUID [FK]\\name: VARCHAR [U]\\base\_\allowbreak{}servings: INTEGER};
\node[modelphysical] (ingredient) at (16,0) {\textbf{ingredient}\\\rule{\linewidth}{0.3pt}\\id: UUID [PK]\\name: VARCHAR [U]\\unit: VARCHAR};
\node[modelphysical] (course_session) at (0,-5) {\textbf{course\_\allowbreak{}session}\\\rule{\linewidth}{0.3pt}\\id: UUID [PK]\\course\_\allowbreak{}id: UUID [FK]\\sequence\_\allowbreak{}number: INTEGER\\scheduled\_\allowbreak{}at: TIMESTAMP\\mode: VARCHAR\\duration\_\allowbreak{}minutes: INTEGER};
\node[modelphysical] (session_recipe) at (8,-5) {\textbf{session\_\allowbreak{}recipe}\\\rule{\linewidth}{0.3pt}\\session\_\allowbreak{}id: UUID [PK,FK]\\recipe\_\allowbreak{}id: UUID [PK,FK]\\portions\_\allowbreak{}per\_\allowbreak{}participant: NUMERIC};
\node[modelphysical] (recipe_ingredient) at (16,-5) {\textbf{recipe\_\allowbreak{}ingredient}\\\rule{\linewidth}{0.3pt}\\recipe\_\allowbreak{}id: UUID [PK,FK]\\ingredient\_\allowbreak{}id: UUID [PK,FK]\\quantity: NUMERIC};
\node[modelphysical] (session_menu_snapshot) at (8,-10) {\textbf{session\_\allowbreak{}menu\_\allowbreak{}snapshot}\\\rule{\linewidth}{0.3pt}\\session\_\allowbreak{}id: UUID [PK,FK]\\recipe\_\allowbreak{}id: UUID [PK,FK]\\recipe\_\allowbreak{}name: TEXT\\ingredients: JSONB\\provenance: TEXT};
\draw[modelassoc] (category.east) -- node[pos=.15,fill=white,inner sep=1pt,font=\footnotesize]{1} node[pos=.85,fill=white,inner sep=1pt,font=\footnotesize]{0..*} (recipe.west);
\draw[modelassoc] (recipe.south east) -- node[pos=.15,fill=white,inner sep=1pt,font=\footnotesize]{1} node[pos=.85,fill=white,inner sep=1pt,font=\footnotesize]{0..*} (recipe_ingredient.north west);
\draw[modelassoc] (ingredient.south) -- node[pos=.15,fill=white,inner sep=1pt,font=\footnotesize]{1} node[pos=.85,fill=white,inner sep=1pt,font=\footnotesize]{0..*} (recipe_ingredient.north);
\draw[modelassoc] (recipe.south) -- node[pos=.15,fill=white,inner sep=1pt,font=\footnotesize]{1} node[pos=.85,fill=white,inner sep=1pt,font=\footnotesize]{0..*} (session_recipe.north);
\draw[modelassoc] (course_session.east) -- node[pos=.15,fill=white,inner sep=1pt,font=\footnotesize]{1} node[pos=.85,fill=white,inner sep=1pt,font=\footnotesize]{0..*} (session_recipe.west);
\draw[modelassoc] (session_recipe.south) -- node[pos=.15,fill=white,inner sep=1pt,font=\footnotesize]{1} node[pos=.85,fill=white,inner sep=1pt,font=\footnotesize]{0..1} (session_menu_snapshot.north);
\end{tikzpicture}
\captionof{figure}{UML ristrutturato: ricettario, menu e copie}\label{fig:diagram-7}
\end{minipage}\end{center}
{\small Le PK di recipe\_ingredient e session\_recipe sono coppie di UUID. La FK della copia usa insieme session\_id e recipe\_id. Essendo anche PK, quella coppia ammette al massimo una copia per associazione; il trigger la crea nel flusso ordinario. Le quantità NUMERIC mantengono la precisione delle dosi.}
\end{landscape}

\clearpage\begin{landscape}
\subsection{UML ristrutturato: comunicazioni e destinatari}
\begin{center}\begin{minipage}{\linewidth}
\centering
\begin{tikzpicture}[x=1cm,y=1cm]
\node[modelphysical] (chef_profile) at (0,0) {\textbf{chef\_\allowbreak{}profile}\\\rule{\linewidth}{0.3pt}\\id: UUID [PK]\\account\_\allowbreak{}id: UUID [FK,U]\\display\_\allowbreak{}name: VARCHAR};
\node[modelphysical] (notification) at (8,0) {\textbf{notification}\\\rule{\linewidth}{0.3pt}\\id: UUID [PK]\\chef\_\allowbreak{}profile\_\allowbreak{}id: UUID [FK]\\course\_\allowbreak{}id: UUID [FK,?]\\title: VARCHAR\\message: VARCHAR\\created\_\allowbreak{}at: TIMESTAMPTZ};
\node[modelphysical] (course) at (16,0) {\textbf{course}\\\rule{\linewidth}{0.3pt}\\id: UUID [PK]\\chef\_\allowbreak{}profile\_\allowbreak{}id: UUID [FK]\\category\_\allowbreak{}id: UUID [FK]\\title: VARCHAR\\start\_\allowbreak{}date: DATE\\status: VARCHAR};
\node[modelphysical] (participant_profile) at (0,-5) {\textbf{participant\_\allowbreak{}profile}\\\rule{\linewidth}{0.3pt}\\id: UUID [PK]\\account\_\allowbreak{}id: UUID [FK,U]\\display\_\allowbreak{}name: VARCHAR};
\node[modelphysical] (notification_recipient) at (8,-5) {\textbf{notification\_\allowbreak{}recipient}\\\rule{\linewidth}{0.3pt}\\notification\_\allowbreak{}id: UUID [PK,FK]\\participant\_\allowbreak{}profile\_\allowbreak{}id: UUID [PK,FK]\\read\_\allowbreak{}at: TIMESTAMPTZ [?]};
\draw[modelassoc] (chef_profile.east) -- node[pos=.15,fill=white,inner sep=1pt,font=\footnotesize]{1} node[pos=.85,fill=white,inner sep=1pt,font=\footnotesize]{0..*} (notification.west);
\draw[modelassoc] (course.west) -- node[pos=.15,fill=white,inner sep=1pt,font=\footnotesize]{0..1} node[pos=.85,fill=white,inner sep=1pt,font=\footnotesize]{0..*} (notification.east);
\draw[modelassoc] (notification.south) -- node[pos=.15,fill=white,inner sep=1pt,font=\footnotesize]{1} node[pos=.85,fill=white,inner sep=1pt,font=\footnotesize]{0..*} (notification_recipient.north);
\draw[modelassoc] (participant_profile.east) -- node[pos=.15,fill=white,inner sep=1pt,font=\footnotesize]{1} node[pos=.85,fill=white,inner sep=1pt,font=\footnotesize]{0..*} (notification_recipient.west);
\end{tikzpicture}
\captionof{figure}{UML ristrutturato: comunicazioni e destinatari}\label{fig:restructured-messages}
\end{minipage}\end{center}
{\small notification.course\_id è nullable: ogni messaggio può riguardare zero o un corso. La PK di notification\_recipient è la coppia notification\_id/participant\_profile\_id. read\_at è opzionale, così un destinatario non ancora letto rimane comunque registrato.}
\end{landscape}

\clearpage\begin{landscape}
\subsection{UML ristrutturato: storico e supporto dello schema}
\begin{center}\begin{minipage}{\linewidth}
\centering
\begin{tikzpicture}[x=1cm,y=1cm]
\node[modelphysical] (session_change_history) at (0,0) {\textbf{session\_\allowbreak{}change\_\allowbreak{}history}\\\rule{\linewidth}{0.3pt}\\id: BIGINT [PK]\\session\_\allowbreak{}id: UUID [FK]\\changed\_\allowbreak{}at: TIMESTAMPTZ\\previous\_\allowbreak{}schedule: TIMESTAMP\\next\_\allowbreak{}schedule: TIMESTAMP};
\node[modelphysical] (course_session) at (8,0) {\textbf{course\_\allowbreak{}session}\\\rule{\linewidth}{0.3pt}\\id: UUID [PK]\\course\_\allowbreak{}id: UUID [FK]\\sequence\_\allowbreak{}number: INTEGER\\scheduled\_\allowbreak{}at: TIMESTAMP\\mode: VARCHAR\\duration\_\allowbreak{}minutes: INTEGER};
\node[modelphysical] (session_outcome) at (16,0) {\textbf{session\_\allowbreak{}outcome}\\\rule{\linewidth}{0.3pt}\\session\_\allowbreak{}id: UUID [PK,FK]\\recorded\_\allowbreak{}at: TIMESTAMPTZ\\note: TEXT};
\node[modelphysical] (session_preparation_revision) at (8,-5) {\textbf{session\_\allowbreak{}preparation\_\allowbreak{}revision}\\\rule{\linewidth}{0.3pt}\\session\_\allowbreak{}id: UUID [PK,FK]\\revision: INTEGER [PK]\\preparation: JSONB\\active: BOOLEAN};
\node[modelphysical] (schema_metadata) at (0,-5) {\textbf{schema\_\allowbreak{}metadata}\\\rule{\linewidth}{0.3pt}\\component: TEXT [PK]\\installed\_\allowbreak{}at: TIMESTAMPTZ};
\draw[modelassoc] (course_session.west) -- node[pos=.15,fill=white,inner sep=1pt,font=\footnotesize]{1} node[pos=.85,fill=white,inner sep=1pt,font=\footnotesize]{0..*} (session_change_history.east);
\draw[modelassoc] (course_session.east) -- node[pos=.15,fill=white,inner sep=1pt,font=\footnotesize]{1} node[pos=.85,fill=white,inner sep=1pt,font=\footnotesize]{0..1} (session_outcome.west);
\draw[modelassoc] (course_session.south) -- node[pos=.15,fill=white,inner sep=1pt,font=\footnotesize]{1} node[pos=.85,fill=white,inner sep=1pt,font=\footnotesize]{0..*} (session_preparation_revision.north);
\end{tikzpicture}
\captionof{figure}{UML ristrutturato: storico e supporto dello schema}\label{fig:restructured-history}
\end{minipage}\end{center}
{\small session\_change\_history registra ancora spostamenti e annullamenti tramite trigger. session\_outcome e session\_preparation\_revision conservano esiti e revisioni storici, senza governare lo svolgimento o i fabbisogni correnti. La revisione è identificata da sessione/numero; un indice parziale limita a una la revisione attiva per sessione. schema\_metadata è tecnica e non ha FK. Gli array UUID e i dati JSONB non aggiungono FK implicite.}
\end{landscape}

\subsection{Dizionario delle classi concettuali}\label{sec:conceptual-dictionary}

{\small \renewcommand{\arraystretch}{1.2}

\begin{xltabular}{\linewidth}{>{\hsize=0.540\hsize\linewidth=\hsize\raggedright\arraybackslash}X>{\hsize=1.380\hsize\linewidth=\hsize\raggedright\arraybackslash}X>{\hsize=1.080\hsize\linewidth=\hsize\raggedright\arraybackslash}X}
\caption{Dizionario delle classi concettuali}\label{tab:dictionary-2}\\
\toprule

\textbf{Classe} & \textbf{Significato} & \textbf{Tabella}\\\midrule\endfirsthead\toprule\textbf{Classe} & \textbf{Significato} & \textbf{Tabella}\\\midrule\endhead
\midrule\multicolumn{3}{r}{\footnotesize Continua nella pagina successiva}\\\endfoot\bottomrule\endlastfoot


Account & Identità autenticabile con ruolo, credenziale e attivazione. & {\ttfamily\small app.\allowbreak{}account}\\\addlinespace[3pt]

Chef & Organizza corsi, menu, comunicazioni e report. & {\ttfamily\small app.\allowbreak{}chef\_\allowbreak{}profile}\\\addlinespace[3pt]

Partecipante & Si iscrive ai corsi e risponde alle pratiche. & {\ttfamily\small app.\allowbreak{}participant\_\allowbreak{}profile}\\\addlinespace[3pt]

Categoria & Classifica corsi e ricette. & {\ttfamily\small app.\allowbreak{}category}\\\addlinespace[3pt]

Corso & Percorso tematico con calendario e frequenza. & {\ttfamily\small app.\allowbreak{}course}\\\addlinespace[3pt]

Sessione & Appuntamento online o pratico con data, durata e dati specifici. & {\ttfamily\small app.\allowbreak{}course\_\allowbreak{}session}\\\addlinespace[3pt]

Ricetta & Ricetta corrente per una persona. & {\ttfamily\small app.\allowbreak{}recipe}\\\addlinespace[3pt]

Ingrediente & Ingrediente con unità canonica. & {\ttfamily\small app.\allowbreak{}ingredient}\\\addlinespace[3pt]

DoseRicetta & Quantità di un ingrediente in una ricetta. & {\ttfamily\small app.\allowbreak{}recipe\_\allowbreak{}ingredient}\\\addlinespace[3pt]

RicettaSessione & Ricetta associata al menu di una pratica. & {\ttfamily\small app.\allowbreak{}session\_\allowbreak{}recipe}\\\addlinespace[3pt]

CopiaMenu & Nome, descrizione e dosi acquisite per la sessione. & {\ttfamily\small app.\allowbreak{}session\_\allowbreak{}menu\_\allowbreak{}snapshot}\\\addlinespace[3pt]

Iscrizione & Legame partecipante/corso con stato e periodo corrente. & {\ttfamily\small app.\allowbreak{}enrollment}\\\addlinespace[3pt]

Adesione & Risposta esplicita alla pratica di un corso a cui si è iscritti. & {\ttfamily\small app.\allowbreak{}practical\_\allowbreak{}attendance}\\\addlinespace[3pt]

Comunicazione & Messaggio interno dello Chef per uno o tutti i corsi. & {\ttfamily\small app.\allowbreak{}notification}\\\addlinespace[3pt]

Destinatario & Destinatario registrato con lettura opzionale. & {\ttfamily\small app.\allowbreak{}notification\_\allowbreak{}recipient}\\*

\end{xltabular}}

\subsection{Dizionario delle associazioni}

Le cardinalità esprimono quante istanze possono essere collegate; l’ultima colonna indica come il collegamento viene conservato e perché. Nella \cref{tab:dictionary-3} sono comprese anche le relazioni non mostrate nelle singole tavole.


{\small \renewcommand{\arraystretch}{1.2}

\begin{xltabular}{\linewidth}{>{\hsize=0.720\hsize\linewidth=\hsize\raggedright\arraybackslash}X>{\hsize=0.480\hsize\linewidth=\hsize\raggedright\arraybackslash}X>{\hsize=1.800\hsize\linewidth=\hsize\raggedright\arraybackslash}X}
\caption{Dizionario delle associazioni}\label{tab:dictionary-3}\\
\toprule

\textbf{Associazione} & \textbf{Cardinalità} & \textbf{Realizzazione e motivazione}\\\midrule\endfirsthead\toprule\textbf{Associazione} & \textbf{Cardinalità} & \textbf{Realizzazione e motivazione}\\\midrule\endhead
\midrule\multicolumn{3}{r}{\footnotesize Continua nella pagina successiva}\\\endfoot\bottomrule\endlastfoot


Account / Chef & 1 : 0..1 & account\_id UNIQUE FK; trigger ruolo; profilo creato insieme ad account da routine\\\addlinespace[3pt]

Account / Partecipante & 1 : 0..1 & account\_id UNIQUE FK; trigger ruolo; profili incompatibili impediti dal ruolo\\\addlinespace[3pt]

Chef / Corso & 1 : 0..N & chef\_profile\_id FK\\\addlinespace[3pt]

Categoria / Corso & 1 : 0..N & category\_id FK\\\addlinespace[3pt]

Categoria / Ricetta & 1 : 0..N & category\_id FK\\\addlinespace[3pt]

Corso / Sessione & 1 : 0..N & course\_id FK; creazione applicativa iniziale di 2--30 sessioni; il CHECK sul conteggio ammette 1--30, ma non impone il minimo di righe\\\addlinespace[3pt]

Ricetta / DoseRicetta & 1 : 0..N & PK ricetta/ingrediente\\\addlinespace[3pt]

Ingrediente / DoseRicetta & 1 : 0..N & ingrediente in più ricette\\\addlinespace[3pt]

Sessione / RicettaSessione & 1 : 0..N & solo pratica, più ricette consentite\\\addlinespace[3pt]

Ricetta / RicettaSessione & 1 : 0..N & recipe\_id FK\\\addlinespace[3pt]

RicettaSessione / CopiaMenu & 1 : 0..1 & FK composita; trigger inserimento produce copia ordinariamente\\\addlinespace[3pt]

Partecipante / Iscrizione & 1 : 0..N & FK, UNIQUE partecipante/corso\\\addlinespace[3pt]

Corso / Iscrizione & 1 : 0..N & FK\\\addlinespace[3pt]

Iscrizione / Adesione & 1 : 0..N & FK\\\addlinespace[3pt]

Sessione / Adesione & 1 : 0..N & FK, solo pratica stesso corso\\\addlinespace[3pt]

Chef / Comunicazione & 1 : 0..N & autore FK\\\addlinespace[3pt]

Corso / Comunicazione & 0..1 : 0..N & course\_id nullable per messaggi globali\\\addlinespace[3pt]

Comunicazione / Destinatario & 1 : 0..N & FK\\\addlinespace[3pt]

Partecipante / Destinatario & 1 : 0..N & FK\\*

\end{xltabular}}

\subsection{Schema logico e dizionario degli attributi}\label{sec:logical-schema}

Sono utilizzati identificatori UUID per gli oggetti referenziati da più relazioni, così una correzione del nome non cambia i collegamenti. I vincoli UNIQUE proteggono separatamente le identità di dominio, come lo username e la coppia partecipante/corso. Le associative senza identità autonoma usano la coppia delle FK come chiave primaria. NUMERIC e BigDecimal conservano le dosi esatte: un CHECK rifiuta oltre tre decimali invece di arrotondarli. Gli appuntamenti sono date e orari locali interpretati in Europe/Rome; accessi e comunicazioni usano timestamp con fuso per rappresentare istanti.

Nella \cref{tab:dictionary-4} ogni relazione elenca i propri attributi; PK indica la chiave primaria. Le FK e le chiavi alternative sono dettagliate subito dopo, insieme ai tipi e ai vincoli.

{\small \renewcommand{\arraystretch}{1.2}

\begin{xltabular}{\linewidth}{>{\hsize=1.500\hsize\linewidth=\hsize\raggedright\arraybackslash}X>{\hsize=0.500\hsize\linewidth=\hsize\raggedright\arraybackslash}X}
\caption{Schema logico e dizionario degli attributi}\label{tab:dictionary-4}\\
\toprule

\textbf{Relazione e attributi} & \textbf{Chiave primaria}\\\midrule\endfirsthead\toprule\textbf{Relazione e attributi} & \textbf{Chiave primaria}\\\midrule\endhead
\midrule\multicolumn{2}{r}{\footnotesize Continua nella pagina successiva}\\\endfoot\bottomrule\endlastfoot


{\ttfamily\small app.\allowbreak{}account(\allowbreak{}id,\allowbreak{} username,\allowbreak{} password\_\allowbreak{}hash,\allowbreak{} role,\allowbreak{} active,\allowbreak{} created\_\allowbreak{}at)} & {\ttfamily\small PRIMARY KEY (\allowbreak{}id)}\\\addlinespace[3pt]

{\ttfamily\small app.\allowbreak{}account\_\allowbreak{}password\_\allowbreak{}reset(\allowbreak{}token\_\allowbreak{}hash,\allowbreak{} account\_\allowbreak{}id,\allowbreak{} created\_\allowbreak{}at,\allowbreak{} expires\_\allowbreak{}at,\allowbreak{} used\_\allowbreak{}at)} & {\ttfamily\small PRIMARY KEY (\allowbreak{}token\_\allowbreak{}hash)}\\\addlinespace[3pt]

{\ttfamily\small app.\allowbreak{}auth\_\allowbreak{}session(\allowbreak{}token\_\allowbreak{}hash,\allowbreak{} account\_\allowbreak{}id,\allowbreak{} created\_\allowbreak{}at,\allowbreak{} expires\_\allowbreak{}at)} & {\ttfamily\small PRIMARY KEY (\allowbreak{}token\_\allowbreak{}hash)}\\\addlinespace[3pt]

{\ttfamily\small app.\allowbreak{}category(\allowbreak{}id,\allowbreak{} name,\allowbreak{} description)} & {\ttfamily\small PRIMARY KEY (\allowbreak{}id)}\\\addlinespace[3pt]

{\ttfamily\small app.\allowbreak{}chef\_\allowbreak{}invitation(\allowbreak{}token\_\allowbreak{}hash,\allowbreak{} created\_\allowbreak{}at,\allowbreak{} expires\_\allowbreak{}at,\allowbreak{} used\_\allowbreak{}at)} & {\ttfamily\small PRIMARY KEY (\allowbreak{}token\_\allowbreak{}hash)}\\\addlinespace[3pt]

{\ttfamily\small app.\allowbreak{}chef\_\allowbreak{}profile(\allowbreak{}id,\allowbreak{} account\_\allowbreak{}id,\allowbreak{} display\_\allowbreak{}name)} & {\ttfamily\small PRIMARY KEY (\allowbreak{}id)}\\\addlinespace[3pt]

{\ttfamily\small app.\allowbreak{}course(\allowbreak{}id,\allowbreak{} chef\_\allowbreak{}profile\_\allowbreak{}id,\allowbreak{} category\_\allowbreak{}id,\allowbreak{} title,\allowbreak{} description,\allowbreak{} start\_\allowbreak{}date,\allowbreak{} start\_\allowbreak{}time,\allowbreak{} recurrence\_\allowbreak{}days,\allowbreak{} session\_\allowbreak{}count,\allowbreak{} status,\allowbreak{} created\_\allowbreak{}at,\allowbreak{} ingredient\_\allowbreak{}margin)} & {\ttfamily\small PRIMARY KEY (\allowbreak{}id)}\\\addlinespace[3pt]

{\ttfamily\small app.\allowbreak{}course\_\allowbreak{}session(\allowbreak{}id,\allowbreak{} course\_\allowbreak{}id,\allowbreak{} sequence\_\allowbreak{}number,\allowbreak{} scheduled\_\allowbreak{}at,\allowbreak{} mode,\allowbreak{} status,\allowbreak{} duration\_\allowbreak{}minutes,\allowbreak{} online\_\allowbreak{}url,\allowbreak{} venue,\allowbreak{} capacity,\allowbreak{} attendance\_\allowbreak{}closes\_\allowbreak{}at,\allowbreak{} completed\_\allowbreak{}at)} & {\ttfamily\small PRIMARY KEY (\allowbreak{}id)}\\\addlinespace[3pt]

{\ttfamily\small app.\allowbreak{}enrollment(\allowbreak{}id,\allowbreak{} participant\_\allowbreak{}profile\_\allowbreak{}id,\allowbreak{} course\_\allowbreak{}id,\allowbreak{} status,\allowbreak{} enrolled\_\allowbreak{}at,\allowbreak{} withdrawn\_\allowbreak{}at)} & {\ttfamily\small PRIMARY KEY (\allowbreak{}id)}\\\addlinespace[3pt]

{\ttfamily\small app.\allowbreak{}ingredient(\allowbreak{}id,\allowbreak{} name,\allowbreak{} unit)} & {\ttfamily\small PRIMARY KEY (\allowbreak{}id)}\\\addlinespace[3pt]

{\ttfamily\small app.\allowbreak{}notification(\allowbreak{}id,\allowbreak{} chef\_\allowbreak{}profile\_\allowbreak{}id,\allowbreak{} course\_\allowbreak{}id,\allowbreak{} title,\allowbreak{} message,\allowbreak{} reason,\allowbreak{} created\_\allowbreak{}at)} & {\ttfamily\small PRIMARY KEY (\allowbreak{}id)}\\\addlinespace[3pt]

{\ttfamily\small app.\allowbreak{}notification\_\allowbreak{}recipient(\allowbreak{}notification\_\allowbreak{}id,\allowbreak{} participant\_\allowbreak{}profile\_\allowbreak{}id,\allowbreak{} read\_\allowbreak{}at)} & {\ttfamily\small PRIMARY KEY (\allowbreak{}notification\_\allowbreak{}id,\allowbreak{} participant\_\allowbreak{}profile\_\allowbreak{}id)}\\\addlinespace[3pt]

{\ttfamily\small app.\allowbreak{}participant\_\allowbreak{}profile(\allowbreak{}id,\allowbreak{} account\_\allowbreak{}id,\allowbreak{} display\_\allowbreak{}name)} & {\ttfamily\small PRIMARY KEY (\allowbreak{}id)}\\\addlinespace[3pt]

{\ttfamily\small app.\allowbreak{}practical\_\allowbreak{}attendance(\allowbreak{}enrollment\_\allowbreak{}id,\allowbreak{} session\_\allowbreak{}id,\allowbreak{} status,\allowbreak{} updated\_\allowbreak{}at)} & {\ttfamily\small PRIMARY KEY (\allowbreak{}enrollment\_\allowbreak{}id,\allowbreak{} session\_\allowbreak{}id)}\\\addlinespace[3pt]

{\ttfamily\small app.\allowbreak{}recipe(\allowbreak{}id,\allowbreak{} category\_\allowbreak{}id,\allowbreak{} name,\allowbreak{} description,\allowbreak{} base\_\allowbreak{}servings)} & {\ttfamily\small PRIMARY KEY (\allowbreak{}id)}\\\addlinespace[3pt]

{\ttfamily\small app.\allowbreak{}recipe\_\allowbreak{}ingredient(\allowbreak{}recipe\_\allowbreak{}id,\allowbreak{} ingredient\_\allowbreak{}id,\allowbreak{} quantity)} & {\ttfamily\small PRIMARY KEY (\allowbreak{}recipe\_\allowbreak{}id,\allowbreak{} ingredient\_\allowbreak{}id)}\\\addlinespace[3pt]

{\ttfamily\small app.\allowbreak{}schema\_\allowbreak{}metadata(\allowbreak{}component,\allowbreak{} installed\_\allowbreak{}at)} & {\ttfamily\small PRIMARY KEY (\allowbreak{}component)}\\\addlinespace[3pt]

{\ttfamily\small app.\allowbreak{}session\_\allowbreak{}change\_\allowbreak{}history(\allowbreak{}id,\allowbreak{} session\_\allowbreak{}id,\allowbreak{} changed\_\allowbreak{}at,\allowbreak{} previous\_\allowbreak{}schedule,\allowbreak{} next\_\allowbreak{}schedule,\allowbreak{} previous\_\allowbreak{}status,\allowbreak{} next\_\allowbreak{}status,\allowbreak{} previous\_\allowbreak{}attendance)} & {\ttfamily\small PRIMARY KEY (\allowbreak{}id)}\\\addlinespace[3pt]

{\ttfamily\small app.\allowbreak{}session\_\allowbreak{}menu\_\allowbreak{}snapshot(\allowbreak{}session\_\allowbreak{}id,\allowbreak{} recipe\_\allowbreak{}id,\allowbreak{} recipe\_\allowbreak{}name,\allowbreak{} description,\allowbreak{} ingredients,\allowbreak{} captured\_\allowbreak{}at,\allowbreak{} provenance)} & {\ttfamily\small PRIMARY KEY (\allowbreak{}session\_\allowbreak{}id,\allowbreak{} recipe\_\allowbreak{}id)}\\\addlinespace[3pt]

{\ttfamily\small app.\allowbreak{}session\_\allowbreak{}outcome(\allowbreak{}session\_\allowbreak{}id,\allowbreak{} recorded\_\allowbreak{}at,\allowbreak{} recipe\_\allowbreak{}ids,\allowbreak{} enrollment\_\allowbreak{}ids,\allowbreak{} note)} & {\ttfamily\small PRIMARY KEY (\allowbreak{}session\_\allowbreak{}id)}\\\addlinespace[3pt]

{\ttfamily\small app.\allowbreak{}session\_\allowbreak{}preparation\_\allowbreak{}revision(\allowbreak{}session\_\allowbreak{}id,\allowbreak{} revision,\allowbreak{} created\_\allowbreak{}at,\allowbreak{} reason,\allowbreak{} preparation,\allowbreak{} confirmed\_\allowbreak{}enrollments,\allowbreak{} active)} & {\ttfamily\small PRIMARY KEY (\allowbreak{}session\_\allowbreak{}id,\allowbreak{} revision)}\\\addlinespace[3pt]

{\ttfamily\small app.\allowbreak{}session\_\allowbreak{}recipe(\allowbreak{}session\_\allowbreak{}id,\allowbreak{} recipe\_\allowbreak{}id,\allowbreak{} portions\_\allowbreak{}per\_\allowbreak{}participant)} & {\ttfamily\small PRIMARY KEY (\allowbreak{}session\_\allowbreak{}id,\allowbreak{} recipe\_\allowbreak{}id)}\\*

\end{xltabular}}

La struttura delle tabelle è definita dalla migrazione iniziale e rimane invariata nella migrazione correttiva. Comprende 22 tabelle applicative e 115 colonne; la tabella di cronologia Flyway è metadato del framework e non viene contata. Per ciascuna relazione sono indicati attributi, tipi, nullabilità e chiavi. Nella colonna NULL, sì indica un attributo facoltativo e no indica NOT NULL. Il trattino nella colonna Default indica l’assenza di un valore predefinito.

\Needspace{10\baselineskip}
\subsection{app.account}

Classificazione: operativa. Identità autenticabile con ruolo, credenziale e attivazione.

{\small \renewcommand{\arraystretch}{1.2}

\begin{xltabular}{\linewidth}{>{\hsize=1.120\hsize\linewidth=\hsize\raggedright\arraybackslash}X>{\hsize=1.200\hsize\linewidth=\hsize\raggedright\arraybackslash}X>{\hsize=0.400\hsize\linewidth=\hsize\raggedright\arraybackslash}X>{\hsize=1.280\hsize\linewidth=\hsize\raggedright\arraybackslash}X}
\caption{Attributi di app.account}\label{tab:dictionary-5}\\
\toprule

\textbf{Attributo} & \textbf{Tipo fisico} & \textbf{NULL} & \textbf{Default}\\\midrule\endfirsthead\toprule\textbf{Attributo} & \textbf{Tipo fisico} & \textbf{NULL} & \textbf{Default}\\\midrule\endhead
\midrule\multicolumn{4}{r}{\footnotesize Continua nella pagina successiva}\\\endfoot\bottomrule\endlastfoot


{\ttfamily\small id} & {\ttfamily\small uuid} & no & {\ttfamily\small gen\_\allowbreak{}random\_\allowbreak{}uuid(\allowbreak{})}\\\addlinespace[3pt]

{\ttfamily\small username} & {\ttfamily\small character varying(\allowbreak{}50)} & no & {\ttfamily\small --}\\\addlinespace[3pt]

{\ttfamily\small password\_\allowbreak{}hash} & {\ttfamily\small text} & no & {\ttfamily\small --}\\\addlinespace[3pt]

{\ttfamily\small role} & {\ttfamily\small character varying(\allowbreak{}20)} & no & {\ttfamily\small --}\\\addlinespace[3pt]

{\ttfamily\small active} & {\ttfamily\small boolean} & no & {\ttfamily\small true}\\\addlinespace[3pt]

{\ttfamily\small created\_\allowbreak{}at} & {\ttfamily\small timestamp with time zone} & no & {\ttfamily\small clock\_\allowbreak{}timestamp(\allowbreak{})}\\*

\end{xltabular}}

{\ttfamily\small PRIMARY KEY (\allowbreak{}id)}\par

{\ttfamily\small UNIQUE (\allowbreak{}username)}\par

\Needspace{10\baselineskip}
\subsection{app.account\_password\_reset}

Classificazione: storica. Archivio dei token di recupero, senza operazioni disponibili nel client.

{\small \renewcommand{\arraystretch}{1.2}

\begin{xltabular}{\linewidth}{>{\hsize=1.120\hsize\linewidth=\hsize\raggedright\arraybackslash}X>{\hsize=1.200\hsize\linewidth=\hsize\raggedright\arraybackslash}X>{\hsize=0.400\hsize\linewidth=\hsize\raggedright\arraybackslash}X>{\hsize=1.280\hsize\linewidth=\hsize\raggedright\arraybackslash}X}
\caption{Attributi di app.account\_password\_reset}\label{tab:dictionary-6}\\
\toprule

\textbf{Attributo} & \textbf{Tipo fisico} & \textbf{NULL} & \textbf{Default}\\\midrule\endfirsthead\toprule\textbf{Attributo} & \textbf{Tipo fisico} & \textbf{NULL} & \textbf{Default}\\\midrule\endhead
\midrule\multicolumn{4}{r}{\footnotesize Continua nella pagina successiva}\\\endfoot\bottomrule\endlastfoot


{\ttfamily\small token\_\allowbreak{}hash} & {\ttfamily\small bytea} & no & {\ttfamily\small --}\\\addlinespace[3pt]

{\ttfamily\small account\_\allowbreak{}id} & {\ttfamily\small uuid} & no & {\ttfamily\small --}\\\addlinespace[3pt]

{\ttfamily\small created\_\allowbreak{}at} & {\ttfamily\small timestamp with time zone} & no & {\ttfamily\small clock\_\allowbreak{}timestamp(\allowbreak{})}\\\addlinespace[3pt]

{\ttfamily\small expires\_\allowbreak{}at} & {\ttfamily\small timestamp with time zone} & no & {\ttfamily\small --}\\\addlinespace[3pt]

{\ttfamily\small used\_\allowbreak{}at} & {\ttfamily\small timestamp with time zone} & sì & {\ttfamily\small --}\\*

\end{xltabular}}

{\ttfamily\small FOREIGN KEY (\allowbreak{}account\_\allowbreak{}id) REFERENCES app.\allowbreak{}account(\allowbreak{}id)}\par

{\ttfamily\small PRIMARY KEY (\allowbreak{}token\_\allowbreak{}hash)}\par

\Needspace{12\baselineskip}
\subsection{app.auth\_session}

Classificazione: tecnica. Sessioni autenticate: hash del token, account e intervallo di validità.

{\small \renewcommand{\arraystretch}{1.2}

\begin{xltabular}{\linewidth}{>{\hsize=1.120\hsize\linewidth=\hsize\raggedright\arraybackslash}X>{\hsize=1.200\hsize\linewidth=\hsize\raggedright\arraybackslash}X>{\hsize=0.400\hsize\linewidth=\hsize\raggedright\arraybackslash}X>{\hsize=1.280\hsize\linewidth=\hsize\raggedright\arraybackslash}X}
\caption{Attributi di app.auth\_session}\label{tab:dictionary-7}\\
\toprule

\textbf{Attributo} & \textbf{Tipo fisico} & \textbf{NULL} & \textbf{Default}\\\midrule\endfirsthead\toprule\textbf{Attributo} & \textbf{Tipo fisico} & \textbf{NULL} & \textbf{Default}\\\midrule\endhead
\midrule\multicolumn{4}{r}{\footnotesize Continua nella pagina successiva}\\\endfoot\bottomrule\endlastfoot


{\ttfamily\small token\_\allowbreak{}hash} & {\ttfamily\small bytea} & no & {\ttfamily\small --}\\\addlinespace[3pt]

{\ttfamily\small account\_\allowbreak{}id} & {\ttfamily\small uuid} & no & {\ttfamily\small --}\\\addlinespace[3pt]

{\ttfamily\small created\_\allowbreak{}at} & {\ttfamily\small timestamp with time zone} & no & {\ttfamily\small clock\_\allowbreak{}timestamp(\allowbreak{})}\\\addlinespace[3pt]

{\ttfamily\small expires\_\allowbreak{}at} & {\ttfamily\small timestamp with time zone} & no & {\ttfamily\small --}\\*

\end{xltabular}}

{\ttfamily\small FOREIGN KEY (\allowbreak{}account\_\allowbreak{}id) REFERENCES app.\allowbreak{}account(\allowbreak{}id) ON DELETE CASCADE}\par

{\ttfamily\small PRIMARY KEY (\allowbreak{}token\_\allowbreak{}hash)}\par

\Needspace{10\baselineskip}
\subsection{app.category}

Classificazione: operativa. Classifica corsi e ricette.

{\small \renewcommand{\arraystretch}{1.2}

\begin{xltabular}{\linewidth}{>{\hsize=1.120\hsize\linewidth=\hsize\raggedright\arraybackslash}X>{\hsize=1.200\hsize\linewidth=\hsize\raggedright\arraybackslash}X>{\hsize=0.400\hsize\linewidth=\hsize\raggedright\arraybackslash}X>{\hsize=1.280\hsize\linewidth=\hsize\raggedright\arraybackslash}X}
\caption{Attributi di app.category}\label{tab:dictionary-8}\\
\toprule

\textbf{Attributo} & \textbf{Tipo fisico} & \textbf{NULL} & \textbf{Default}\\\midrule\endfirsthead\toprule\textbf{Attributo} & \textbf{Tipo fisico} & \textbf{NULL} & \textbf{Default}\\\midrule\endhead
\midrule\multicolumn{4}{r}{\footnotesize Continua nella pagina successiva}\\\endfoot\bottomrule\endlastfoot


{\ttfamily\small id} & {\ttfamily\small uuid} & no & {\ttfamily\small gen\_\allowbreak{}random\_\allowbreak{}uuid(\allowbreak{})}\\\addlinespace[3pt]

{\ttfamily\small name} & {\ttfamily\small character varying(\allowbreak{}80)} & no & {\ttfamily\small --}\\\addlinespace[3pt]

{\ttfamily\small description} & {\ttfamily\small text} & no & {\ttfamily\small --}\\*

\end{xltabular}}

{\ttfamily\small UNIQUE (\allowbreak{}name)}\par

{\ttfamily\small PRIMARY KEY (\allowbreak{}id)}\par

\Needspace{10\baselineskip}
\subsection{app.chef\_invitation}

Classificazione: storica. Archivio dei codici di invito, senza operazioni disponibili nel client.

{\small \renewcommand{\arraystretch}{1.2}

\begin{xltabular}{\linewidth}{>{\hsize=1.120\hsize\linewidth=\hsize\raggedright\arraybackslash}X>{\hsize=1.200\hsize\linewidth=\hsize\raggedright\arraybackslash}X>{\hsize=0.400\hsize\linewidth=\hsize\raggedright\arraybackslash}X>{\hsize=1.280\hsize\linewidth=\hsize\raggedright\arraybackslash}X}
\caption{Attributi di app.chef\_invitation}\label{tab:dictionary-9}\\
\toprule

\textbf{Attributo} & \textbf{Tipo fisico} & \textbf{NULL} & \textbf{Default}\\\midrule\endfirsthead\toprule\textbf{Attributo} & \textbf{Tipo fisico} & \textbf{NULL} & \textbf{Default}\\\midrule\endhead
\midrule\multicolumn{4}{r}{\footnotesize Continua nella pagina successiva}\\\endfoot\bottomrule\endlastfoot


{\ttfamily\small token\_\allowbreak{}hash} & {\ttfamily\small bytea} & no & {\ttfamily\small --}\\\addlinespace[3pt]

{\ttfamily\small created\_\allowbreak{}at} & {\ttfamily\small timestamp with time zone} & no & {\ttfamily\small clock\_\allowbreak{}timestamp(\allowbreak{})}\\\addlinespace[3pt]

{\ttfamily\small expires\_\allowbreak{}at} & {\ttfamily\small timestamp with time zone} & no & {\ttfamily\small --}\\\addlinespace[3pt]

{\ttfamily\small used\_\allowbreak{}at} & {\ttfamily\small timestamp with time zone} & sì & {\ttfamily\small --}\\*

\end{xltabular}}

{\ttfamily\small PRIMARY KEY (\allowbreak{}token\_\allowbreak{}hash)}\par

\Needspace{10\baselineskip}
\subsection{app.chef\_profile}

Classificazione: operativa. Organizza corsi, menu, comunicazioni e report.

{\small \renewcommand{\arraystretch}{1.2}

\begin{xltabular}{\linewidth}{>{\hsize=1.120\hsize\linewidth=\hsize\raggedright\arraybackslash}X>{\hsize=1.200\hsize\linewidth=\hsize\raggedright\arraybackslash}X>{\hsize=0.400\hsize\linewidth=\hsize\raggedright\arraybackslash}X>{\hsize=1.280\hsize\linewidth=\hsize\raggedright\arraybackslash}X}
\caption{Attributi di app.chef\_profile}\label{tab:dictionary-10}\\
\toprule

\textbf{Attributo} & \textbf{Tipo fisico} & \textbf{NULL} & \textbf{Default}\\\midrule\endfirsthead\toprule\textbf{Attributo} & \textbf{Tipo fisico} & \textbf{NULL} & \textbf{Default}\\\midrule\endhead
\midrule\multicolumn{4}{r}{\footnotesize Continua nella pagina successiva}\\\endfoot\bottomrule\endlastfoot


{\ttfamily\small id} & {\ttfamily\small uuid} & no & {\ttfamily\small gen\_\allowbreak{}random\_\allowbreak{}uuid(\allowbreak{})}\\\addlinespace[3pt]

{\ttfamily\small account\_\allowbreak{}id} & {\ttfamily\small uuid} & no & {\ttfamily\small --}\\\addlinespace[3pt]

{\ttfamily\small display\_\allowbreak{}name} & {\ttfamily\small character varying(\allowbreak{}100)} & no & {\ttfamily\small --}\\*

\end{xltabular}}

{\ttfamily\small FOREIGN KEY (\allowbreak{}account\_\allowbreak{}id) REFERENCES app.\allowbreak{}account(\allowbreak{}id)}\par

{\ttfamily\small UNIQUE (\allowbreak{}account\_\allowbreak{}id)}\par

{\ttfamily\small PRIMARY KEY (\allowbreak{}id)}\par

\Needspace{10\baselineskip}
\subsection{app.course}

Classificazione: operativa. Percorso tematico con calendario e frequenza.

{\small \renewcommand{\arraystretch}{1.2}

\begin{xltabular}{\linewidth}{>{\hsize=1.120\hsize\linewidth=\hsize\raggedright\arraybackslash}X>{\hsize=1.200\hsize\linewidth=\hsize\raggedright\arraybackslash}X>{\hsize=0.400\hsize\linewidth=\hsize\raggedright\arraybackslash}X>{\hsize=1.280\hsize\linewidth=\hsize\raggedright\arraybackslash}X}
\caption{Attributi di app.course}\label{tab:dictionary-11}\\
\toprule

\textbf{Attributo} & \textbf{Tipo fisico} & \textbf{NULL} & \textbf{Default}\\\midrule\endfirsthead\toprule\textbf{Attributo} & \textbf{Tipo fisico} & \textbf{NULL} & \textbf{Default}\\\midrule\endhead
\midrule\multicolumn{4}{r}{\footnotesize Continua nella pagina successiva}\\\endfoot\bottomrule\endlastfoot


{\ttfamily\small id} & {\ttfamily\small uuid} & no & {\ttfamily\small gen\_\allowbreak{}random\_\allowbreak{}uuid(\allowbreak{})}\\\addlinespace[3pt]

{\ttfamily\small chef\_\allowbreak{}profile\_\allowbreak{}id} & {\ttfamily\small uuid} & no & {\ttfamily\small --}\\\addlinespace[3pt]

{\ttfamily\small category\_\allowbreak{}id} & {\ttfamily\small uuid} & no & {\ttfamily\small --}\\\addlinespace[3pt]

{\ttfamily\small title} & {\ttfamily\small character varying(\allowbreak{}120)} & no & {\ttfamily\small --}\\\addlinespace[3pt]

{\ttfamily\small description} & {\ttfamily\small text} & no & {\ttfamily\small --}\\\addlinespace[3pt]

{\ttfamily\small start\_\allowbreak{}date} & {\ttfamily\small date} & no & {\ttfamily\small --}\\\addlinespace[3pt]

{\ttfamily\small start\_\allowbreak{}time} & {\ttfamily\small time without time zone} & no & {\ttfamily\small --}\\\addlinespace[3pt]

{\ttfamily\small recurrence\_\allowbreak{}days} & {\ttfamily\small integer} & no & {\ttfamily\small --}\\\addlinespace[3pt]

{\ttfamily\small session\_\allowbreak{}count} & {\ttfamily\small integer} & no & {\ttfamily\small --}\\\addlinespace[3pt]

{\ttfamily\small status} & {\ttfamily\small character varying(\allowbreak{}20)} & no & {\ttfamily\small 'DRAFT'::character varying}\\\addlinespace[3pt]

{\ttfamily\small created\_\allowbreak{}at} & {\ttfamily\small timestamp with time zone} & no & {\ttfamily\small clock\_\allowbreak{}timestamp(\allowbreak{})}\\\addlinespace[3pt]

{\ttfamily\small ingredient\_\allowbreak{}margin} & {\ttfamily\small numeric} & no & {\ttfamily\small 0}\\*

\end{xltabular}}

{\ttfamily\small FOREIGN KEY (\allowbreak{}category\_\allowbreak{}id) REFERENCES app.\allowbreak{}category(\allowbreak{}id)}\par

{\ttfamily\small FOREIGN KEY (\allowbreak{}chef\_\allowbreak{}profile\_\allowbreak{}id) REFERENCES app.\allowbreak{}chef\_\allowbreak{}profile(\allowbreak{}id)}\par

{\ttfamily\small UNIQUE (\allowbreak{}chef\_\allowbreak{}profile\_\allowbreak{}id,\allowbreak{} title,\allowbreak{} start\_\allowbreak{}date)}\par

{\ttfamily\small PRIMARY KEY (\allowbreak{}id)}\par

\Needspace{10\baselineskip}
\subsection{app.course\_session}

Classificazione: operativa. Appuntamento online o pratico con data, durata e dati specifici.

{\small \renewcommand{\arraystretch}{1.2}

\begin{xltabular}{\linewidth}{>{\hsize=1.120\hsize\linewidth=\hsize\raggedright\arraybackslash}X>{\hsize=1.200\hsize\linewidth=\hsize\raggedright\arraybackslash}X>{\hsize=0.400\hsize\linewidth=\hsize\raggedright\arraybackslash}X>{\hsize=1.280\hsize\linewidth=\hsize\raggedright\arraybackslash}X}
\caption{Attributi di app.course\_session}\label{tab:dictionary-12}\\
\toprule

\textbf{Attributo} & \textbf{Tipo fisico} & \textbf{NULL} & \textbf{Default}\\\midrule\endfirsthead\toprule\textbf{Attributo} & \textbf{Tipo fisico} & \textbf{NULL} & \textbf{Default}\\\midrule\endhead
\midrule\multicolumn{4}{r}{\footnotesize Continua nella pagina successiva}\\\endfoot\bottomrule\endlastfoot


{\ttfamily\small id} & {\ttfamily\small uuid} & no & {\ttfamily\small gen\_\allowbreak{}random\_\allowbreak{}uuid(\allowbreak{})}\\\addlinespace[3pt]

{\ttfamily\small course\_\allowbreak{}id} & {\ttfamily\small uuid} & no & {\ttfamily\small --}\\\addlinespace[3pt]

{\ttfamily\small sequence\_\allowbreak{}number} & {\ttfamily\small integer} & no & {\ttfamily\small --}\\\addlinespace[3pt]

{\ttfamily\small scheduled\_\allowbreak{}at} & {\ttfamily\small timestamp without time zone} & no & {\ttfamily\small --}\\\addlinespace[3pt]

{\ttfamily\small mode} & {\ttfamily\small character varying(\allowbreak{}20)} & no & {\ttfamily\small --}\\\addlinespace[3pt]

{\ttfamily\small status} & {\ttfamily\small character varying(\allowbreak{}20)} & no & {\ttfamily\small 'SCHEDULED'::character varying}\\\addlinespace[3pt]

{\ttfamily\small duration\_\allowbreak{}minutes} & {\ttfamily\small integer} & no & {\ttfamily\small 120}\\\addlinespace[3pt]

{\ttfamily\small online\_\allowbreak{}url} & {\ttfamily\small text} & sì & {\ttfamily\small --}\\\addlinespace[3pt]

{\ttfamily\small venue} & {\ttfamily\small character varying(\allowbreak{}180)} & sì & {\ttfamily\small --}\\\addlinespace[3pt]

{\ttfamily\small capacity} & {\ttfamily\small integer} & sì & {\ttfamily\small --}\\\addlinespace[3pt]

{\ttfamily\small attendance\_\allowbreak{}closes\_\allowbreak{}at} & {\ttfamily\small timestamp without time zone} & sì & {\ttfamily\small --}\\\addlinespace[3pt]

{\ttfamily\small completed\_\allowbreak{}at} & {\ttfamily\small timestamp with time zone} & sì & {\ttfamily\small --}\\*

\end{xltabular}}

{\ttfamily\small FOREIGN KEY (\allowbreak{}course\_\allowbreak{}id) REFERENCES app.\allowbreak{}course(\allowbreak{}id) ON DELETE CASCADE}\par

{\ttfamily\small PRIMARY KEY (\allowbreak{}id)}\par

{\ttfamily\small UNIQUE (\allowbreak{}course\_\allowbreak{}id,\allowbreak{} sequence\_\allowbreak{}number)}\par

\Needspace{10\baselineskip}
\subsection{app.enrollment}

Classificazione: operativa. Legame partecipante/corso con stato e periodo corrente.

{\small \renewcommand{\arraystretch}{1.2}

\begin{xltabular}{\linewidth}{>{\hsize=1.120\hsize\linewidth=\hsize\raggedright\arraybackslash}X>{\hsize=1.200\hsize\linewidth=\hsize\raggedright\arraybackslash}X>{\hsize=0.400\hsize\linewidth=\hsize\raggedright\arraybackslash}X>{\hsize=1.280\hsize\linewidth=\hsize\raggedright\arraybackslash}X}
\caption{Attributi di app.enrollment}\label{tab:dictionary-13}\\
\toprule

\textbf{Attributo} & \textbf{Tipo fisico} & \textbf{NULL} & \textbf{Default}\\\midrule\endfirsthead\toprule\textbf{Attributo} & \textbf{Tipo fisico} & \textbf{NULL} & \textbf{Default}\\\midrule\endhead
\midrule\multicolumn{4}{r}{\footnotesize Continua nella pagina successiva}\\\endfoot\bottomrule\endlastfoot


{\ttfamily\small id} & {\ttfamily\small uuid} & no & {\ttfamily\small gen\_\allowbreak{}random\_\allowbreak{}uuid(\allowbreak{})}\\\addlinespace[3pt]

{\ttfamily\small participant\_\allowbreak{}profile\_\allowbreak{}id} & {\ttfamily\small uuid} & no & {\ttfamily\small --}\\\addlinespace[3pt]

{\ttfamily\small course\_\allowbreak{}id} & {\ttfamily\small uuid} & no & {\ttfamily\small --}\\\addlinespace[3pt]

{\ttfamily\small status} & {\ttfamily\small character varying(\allowbreak{}20)} & no & {\ttfamily\small --}\\\addlinespace[3pt]

{\ttfamily\small enrolled\_\allowbreak{}at} & {\ttfamily\small timestamp with time zone} & no & {\ttfamily\small clock\_\allowbreak{}timestamp(\allowbreak{})}\\\addlinespace[3pt]

{\ttfamily\small withdrawn\_\allowbreak{}at} & {\ttfamily\small timestamp with time zone} & sì & {\ttfamily\small --}\\*

\end{xltabular}}

{\ttfamily\small FOREIGN KEY (\allowbreak{}course\_\allowbreak{}id) REFERENCES app.\allowbreak{}course(\allowbreak{}id)}\par

{\ttfamily\small UNIQUE (\allowbreak{}participant\_\allowbreak{}profile\_\allowbreak{}id,\allowbreak{} course\_\allowbreak{}id)}\par

{\ttfamily\small FOREIGN KEY (\allowbreak{}participant\_\allowbreak{}profile\_\allowbreak{}id) REFERENCES app.\allowbreak{}participant\_\allowbreak{}profile(\allowbreak{}id)}\par

{\ttfamily\small PRIMARY KEY (\allowbreak{}id)}\par

\Needspace{10\baselineskip}
\subsection{app.ingredient}

Classificazione: operativa. Ingrediente con unità canonica.

{\small \renewcommand{\arraystretch}{1.2}

\begin{xltabular}{\linewidth}{>{\hsize=1.120\hsize\linewidth=\hsize\raggedright\arraybackslash}X>{\hsize=1.200\hsize\linewidth=\hsize\raggedright\arraybackslash}X>{\hsize=0.400\hsize\linewidth=\hsize\raggedright\arraybackslash}X>{\hsize=1.280\hsize\linewidth=\hsize\raggedright\arraybackslash}X}
\caption{Attributi di app.ingredient}\label{tab:dictionary-14}\\
\toprule

\textbf{Attributo} & \textbf{Tipo fisico} & \textbf{NULL} & \textbf{Default}\\\midrule\endfirsthead\toprule\textbf{Attributo} & \textbf{Tipo fisico} & \textbf{NULL} & \textbf{Default}\\\midrule\endhead
\midrule\multicolumn{4}{r}{\footnotesize Continua nella pagina successiva}\\\endfoot\bottomrule\endlastfoot


{\ttfamily\small id} & {\ttfamily\small uuid} & no & {\ttfamily\small gen\_\allowbreak{}random\_\allowbreak{}uuid(\allowbreak{})}\\\addlinespace[3pt]

{\ttfamily\small name} & {\ttfamily\small character varying(\allowbreak{}100)} & no & {\ttfamily\small --}\\\addlinespace[3pt]

{\ttfamily\small unit} & {\ttfamily\small character varying(\allowbreak{}20)} & no & {\ttfamily\small --}\\*

\end{xltabular}}

{\ttfamily\small UNIQUE (\allowbreak{}name)}\par

{\ttfamily\small PRIMARY KEY (\allowbreak{}id)}\par

\Needspace{10\baselineskip}
\subsection{app.notification}

Classificazione: operativa. Messaggio interno dello Chef per uno o tutti i corsi.

{\small \renewcommand{\arraystretch}{1.2}

\begin{xltabular}{\linewidth}{>{\hsize=1.120\hsize\linewidth=\hsize\raggedright\arraybackslash}X>{\hsize=1.200\hsize\linewidth=\hsize\raggedright\arraybackslash}X>{\hsize=0.400\hsize\linewidth=\hsize\raggedright\arraybackslash}X>{\hsize=1.280\hsize\linewidth=\hsize\raggedright\arraybackslash}X}
\caption{Attributi di app.notification}\label{tab:dictionary-15}\\
\toprule

\textbf{Attributo} & \textbf{Tipo fisico} & \textbf{NULL} & \textbf{Default}\\\midrule\endfirsthead\toprule\textbf{Attributo} & \textbf{Tipo fisico} & \textbf{NULL} & \textbf{Default}\\\midrule\endhead
\midrule\multicolumn{4}{r}{\footnotesize Continua nella pagina successiva}\\\endfoot\bottomrule\endlastfoot


{\ttfamily\small id} & {\ttfamily\small uuid} & no & {\ttfamily\small gen\_\allowbreak{}random\_\allowbreak{}uuid(\allowbreak{})}\\\addlinespace[3pt]

{\ttfamily\small chef\_\allowbreak{}profile\_\allowbreak{}id} & {\ttfamily\small uuid} & no & {\ttfamily\small --}\\\addlinespace[3pt]

{\ttfamily\small course\_\allowbreak{}id} & {\ttfamily\small uuid} & sì & {\ttfamily\small --}\\\addlinespace[3pt]

{\ttfamily\small title} & {\ttfamily\small character varying(\allowbreak{}120)} & no & {\ttfamily\small --}\\\addlinespace[3pt]

{\ttfamily\small message} & {\ttfamily\small character varying(\allowbreak{}2000)} & no & {\ttfamily\small --}\\\addlinespace[3pt]

{\ttfamily\small reason} & {\ttfamily\small character varying(\allowbreak{}30)} & no & {\ttfamily\small --}\\\addlinespace[3pt]

{\ttfamily\small created\_\allowbreak{}at} & {\ttfamily\small timestamp with time zone} & no & {\ttfamily\small clock\_\allowbreak{}timestamp(\allowbreak{})}\\*

\end{xltabular}}

{\ttfamily\small FOREIGN KEY (\allowbreak{}chef\_\allowbreak{}profile\_\allowbreak{}id) REFERENCES app.\allowbreak{}chef\_\allowbreak{}profile(\allowbreak{}id)}\par

{\ttfamily\small FOREIGN KEY (\allowbreak{}course\_\allowbreak{}id) REFERENCES app.\allowbreak{}course(\allowbreak{}id)}\par

{\ttfamily\small PRIMARY KEY (\allowbreak{}id)}\par

\Needspace{10\baselineskip}
\subsection{app.notification\_recipient}

Classificazione: operativa. Destinatario registrato con lettura opzionale.

{\small \renewcommand{\arraystretch}{1.2}

\begin{xltabular}{\linewidth}{>{\hsize=1.120\hsize\linewidth=\hsize\raggedright\arraybackslash}X>{\hsize=1.200\hsize\linewidth=\hsize\raggedright\arraybackslash}X>{\hsize=0.400\hsize\linewidth=\hsize\raggedright\arraybackslash}X>{\hsize=1.280\hsize\linewidth=\hsize\raggedright\arraybackslash}X}
\caption{Attributi di app.notification\_recipient}\label{tab:dictionary-16}\\
\toprule

\textbf{Attributo} & \textbf{Tipo fisico} & \textbf{NULL} & \textbf{Default}\\\midrule\endfirsthead\toprule\textbf{Attributo} & \textbf{Tipo fisico} & \textbf{NULL} & \textbf{Default}\\\midrule\endhead
\midrule\multicolumn{4}{r}{\footnotesize Continua nella pagina successiva}\\\endfoot\bottomrule\endlastfoot


{\ttfamily\small notification\_\allowbreak{}id} & {\ttfamily\small uuid} & no & {\ttfamily\small --}\\\addlinespace[3pt]

{\ttfamily\small participant\_\allowbreak{}profile\_\allowbreak{}id} & {\ttfamily\small uuid} & no & {\ttfamily\small --}\\\addlinespace[3pt]

{\ttfamily\small read\_\allowbreak{}at} & {\ttfamily\small timestamp with time zone} & sì & {\ttfamily\small --}\\*

\end{xltabular}}

{\ttfamily\small FOREIGN KEY (\allowbreak{}notification\_\allowbreak{}id) REFERENCES app.\allowbreak{}notification(\allowbreak{}id) ON DELETE CASCADE}\par

{\ttfamily\small FOREIGN KEY (\allowbreak{}participant\_\allowbreak{}profile\_\allowbreak{}id) REFERENCES app.\allowbreak{}participant\_\allowbreak{}profile(\allowbreak{}id)}\par

{\ttfamily\small PRIMARY KEY (\allowbreak{}notification\_\allowbreak{}id,\allowbreak{} participant\_\allowbreak{}profile\_\allowbreak{}id)}\par

\Needspace{10\baselineskip}
\subsection{app.participant\_profile}

Classificazione: operativa. Si iscrive ai corsi e risponde alle pratiche.

{\small \renewcommand{\arraystretch}{1.2}

\begin{xltabular}{\linewidth}{>{\hsize=1.120\hsize\linewidth=\hsize\raggedright\arraybackslash}X>{\hsize=1.200\hsize\linewidth=\hsize\raggedright\arraybackslash}X>{\hsize=0.400\hsize\linewidth=\hsize\raggedright\arraybackslash}X>{\hsize=1.280\hsize\linewidth=\hsize\raggedright\arraybackslash}X}
\caption{Attributi di app.participant\_profile}\label{tab:dictionary-17}\\
\toprule

\textbf{Attributo} & \textbf{Tipo fisico} & \textbf{NULL} & \textbf{Default}\\\midrule\endfirsthead\toprule\textbf{Attributo} & \textbf{Tipo fisico} & \textbf{NULL} & \textbf{Default}\\\midrule\endhead
\midrule\multicolumn{4}{r}{\footnotesize Continua nella pagina successiva}\\\endfoot\bottomrule\endlastfoot


{\ttfamily\small id} & {\ttfamily\small uuid} & no & {\ttfamily\small gen\_\allowbreak{}random\_\allowbreak{}uuid(\allowbreak{})}\\\addlinespace[3pt]

{\ttfamily\small account\_\allowbreak{}id} & {\ttfamily\small uuid} & no & {\ttfamily\small --}\\\addlinespace[3pt]

{\ttfamily\small display\_\allowbreak{}name} & {\ttfamily\small character varying(\allowbreak{}100)} & no & {\ttfamily\small --}\\*

\end{xltabular}}

{\ttfamily\small FOREIGN KEY (\allowbreak{}account\_\allowbreak{}id) REFERENCES app.\allowbreak{}account(\allowbreak{}id)}\par

{\ttfamily\small UNIQUE (\allowbreak{}account\_\allowbreak{}id)}\par

{\ttfamily\small PRIMARY KEY (\allowbreak{}id)}\par

\Needspace{10\baselineskip}
\subsection{app.practical\_attendance}

Classificazione: operativa. Risposta esplicita alla pratica di un corso a cui si è iscritti.

{\small \renewcommand{\arraystretch}{1.2}

\begin{xltabular}{\linewidth}{>{\hsize=1.120\hsize\linewidth=\hsize\raggedright\arraybackslash}X>{\hsize=1.200\hsize\linewidth=\hsize\raggedright\arraybackslash}X>{\hsize=0.400\hsize\linewidth=\hsize\raggedright\arraybackslash}X>{\hsize=1.280\hsize\linewidth=\hsize\raggedright\arraybackslash}X}
\caption{Attributi di app.practical\_attendance}\label{tab:dictionary-18}\\
\toprule

\textbf{Attributo} & \textbf{Tipo fisico} & \textbf{NULL} & \textbf{Default}\\\midrule\endfirsthead\toprule\textbf{Attributo} & \textbf{Tipo fisico} & \textbf{NULL} & \textbf{Default}\\\midrule\endhead
\midrule\multicolumn{4}{r}{\footnotesize Continua nella pagina successiva}\\\endfoot\bottomrule\endlastfoot


{\ttfamily\small enrollment\_\allowbreak{}id} & {\ttfamily\small uuid} & no & {\ttfamily\small --}\\\addlinespace[3pt]

{\ttfamily\small session\_\allowbreak{}id} & {\ttfamily\small uuid} & no & {\ttfamily\small --}\\\addlinespace[3pt]

{\ttfamily\small status} & {\ttfamily\small character varying(\allowbreak{}20)} & no & {\ttfamily\small --}\\\addlinespace[3pt]

{\ttfamily\small updated\_\allowbreak{}at} & {\ttfamily\small timestamp with time zone} & no & {\ttfamily\small clock\_\allowbreak{}timestamp(\allowbreak{})}\\*

\end{xltabular}}

{\ttfamily\small FOREIGN KEY (\allowbreak{}enrollment\_\allowbreak{}id) REFERENCES app.\allowbreak{}enrollment(\allowbreak{}id)}\par

{\ttfamily\small PRIMARY KEY (\allowbreak{}enrollment\_\allowbreak{}id,\allowbreak{} session\_\allowbreak{}id)}\par

{\ttfamily\small FOREIGN KEY (\allowbreak{}session\_\allowbreak{}id) REFERENCES app.\allowbreak{}course\_\allowbreak{}session(\allowbreak{}id)}\par

\Needspace{10\baselineskip}
\subsection{app.recipe}

Classificazione: operativa. Ricetta corrente per una persona.

{\small \renewcommand{\arraystretch}{1.2}

\begin{xltabular}{\linewidth}{>{\hsize=1.120\hsize\linewidth=\hsize\raggedright\arraybackslash}X>{\hsize=1.200\hsize\linewidth=\hsize\raggedright\arraybackslash}X>{\hsize=0.400\hsize\linewidth=\hsize\raggedright\arraybackslash}X>{\hsize=1.280\hsize\linewidth=\hsize\raggedright\arraybackslash}X}
\caption{Attributi di app.recipe}\label{tab:dictionary-19}\\
\toprule

\textbf{Attributo} & \textbf{Tipo fisico} & \textbf{NULL} & \textbf{Default}\\\midrule\endfirsthead\toprule\textbf{Attributo} & \textbf{Tipo fisico} & \textbf{NULL} & \textbf{Default}\\\midrule\endhead
\midrule\multicolumn{4}{r}{\footnotesize Continua nella pagina successiva}\\\endfoot\bottomrule\endlastfoot


{\ttfamily\small id} & {\ttfamily\small uuid} & no & {\ttfamily\small gen\_\allowbreak{}random\_\allowbreak{}uuid(\allowbreak{})}\\\addlinespace[3pt]

{\ttfamily\small category\_\allowbreak{}id} & {\ttfamily\small uuid} & no & {\ttfamily\small --}\\\addlinespace[3pt]

{\ttfamily\small name} & {\ttfamily\small character varying(\allowbreak{}120)} & no & {\ttfamily\small --}\\\addlinespace[3pt]

{\ttfamily\small description} & {\ttfamily\small text} & no & {\ttfamily\small --}\\\addlinespace[3pt]

{\ttfamily\small base\_\allowbreak{}servings} & {\ttfamily\small integer} & no & {\ttfamily\small --}\\*

\end{xltabular}}

{\ttfamily\small FOREIGN KEY (\allowbreak{}category\_\allowbreak{}id) REFERENCES app.\allowbreak{}category(\allowbreak{}id)}\par

{\ttfamily\small UNIQUE (\allowbreak{}name)}\par

{\ttfamily\small PRIMARY KEY (\allowbreak{}id)}\par

\Needspace{10\baselineskip}
\subsection{app.recipe\_ingredient}

Classificazione: operativa. Quantità di un ingrediente in una ricetta.

{\small \renewcommand{\arraystretch}{1.2}

\begin{xltabular}{\linewidth}{>{\hsize=1.120\hsize\linewidth=\hsize\raggedright\arraybackslash}X>{\hsize=1.200\hsize\linewidth=\hsize\raggedright\arraybackslash}X>{\hsize=0.400\hsize\linewidth=\hsize\raggedright\arraybackslash}X>{\hsize=1.280\hsize\linewidth=\hsize\raggedright\arraybackslash}X}
\caption{Attributi di app.recipe\_ingredient}\label{tab:dictionary-20}\\
\toprule

\textbf{Attributo} & \textbf{Tipo fisico} & \textbf{NULL} & \textbf{Default}\\\midrule\endfirsthead\toprule\textbf{Attributo} & \textbf{Tipo fisico} & \textbf{NULL} & \textbf{Default}\\\midrule\endhead
\midrule\multicolumn{4}{r}{\footnotesize Continua nella pagina successiva}\\\endfoot\bottomrule\endlastfoot


{\ttfamily\small recipe\_\allowbreak{}id} & {\ttfamily\small uuid} & no & {\ttfamily\small --}\\\addlinespace[3pt]

{\ttfamily\small ingredient\_\allowbreak{}id} & {\ttfamily\small uuid} & no & {\ttfamily\small --}\\\addlinespace[3pt]

{\ttfamily\small quantity} & {\ttfamily\small numeric} & no & {\ttfamily\small --}\\*

\end{xltabular}}

{\ttfamily\small FOREIGN KEY (\allowbreak{}ingredient\_\allowbreak{}id) REFERENCES app.\allowbreak{}ingredient(\allowbreak{}id)}\par

{\ttfamily\small PRIMARY KEY (\allowbreak{}recipe\_\allowbreak{}id,\allowbreak{} ingredient\_\allowbreak{}id)}\par

{\ttfamily\small FOREIGN KEY (\allowbreak{}recipe\_\allowbreak{}id) REFERENCES app.\allowbreak{}recipe(\allowbreak{}id) ON DELETE CASCADE}\par

\Needspace{10\baselineskip}
\subsection{app.schema\_metadata}

Classificazione: tecnica. Marcatore tecnico della fondazione dello schema.

{\small \renewcommand{\arraystretch}{1.2}

\begin{xltabular}{\linewidth}{>{\hsize=1.120\hsize\linewidth=\hsize\raggedright\arraybackslash}X>{\hsize=1.200\hsize\linewidth=\hsize\raggedright\arraybackslash}X>{\hsize=0.400\hsize\linewidth=\hsize\raggedright\arraybackslash}X>{\hsize=1.280\hsize\linewidth=\hsize\raggedright\arraybackslash}X}
\caption{Attributi di app.schema\_metadata}\label{tab:dictionary-21}\\
\toprule

\textbf{Attributo} & \textbf{Tipo fisico} & \textbf{NULL} & \textbf{Default}\\\midrule\endfirsthead\toprule\textbf{Attributo} & \textbf{Tipo fisico} & \textbf{NULL} & \textbf{Default}\\\midrule\endhead
\midrule\multicolumn{4}{r}{\footnotesize Continua nella pagina successiva}\\\endfoot\bottomrule\endlastfoot


{\ttfamily\small component} & {\ttfamily\small text} & no & {\ttfamily\small --}\\\addlinespace[3pt]

{\ttfamily\small installed\_\allowbreak{}at} & {\ttfamily\small timestamp with time zone} & no & {\ttfamily\small clock\_\allowbreak{}timestamp(\allowbreak{})}\\*

\end{xltabular}}

{\ttfamily\small PRIMARY KEY (\allowbreak{}component)}\par

\Needspace{10\baselineskip}
\subsection{app.session\_change\_history}

Classificazione: storica. Cronologia del calendario ancora alimentata: il trigger preserve\_session\_change richiama la funzione omonima e registra nuovi spostamenti o annullamenti. È distinta dagli archivi di invito, recupero, revisioni ed esiti non più gestiti dai flussi correnti.

{\small \renewcommand{\arraystretch}{1.2}

\begin{xltabular}{\linewidth}{>{\hsize=1.120\hsize\linewidth=\hsize\raggedright\arraybackslash}X>{\hsize=1.200\hsize\linewidth=\hsize\raggedright\arraybackslash}X>{\hsize=0.400\hsize\linewidth=\hsize\raggedright\arraybackslash}X>{\hsize=1.280\hsize\linewidth=\hsize\raggedright\arraybackslash}X}
\caption{Attributi di app.session\_change\_history}\label{tab:dictionary-22}\\
\toprule

\textbf{Attributo} & \textbf{Tipo fisico} & \textbf{NULL} & \textbf{Default}\\\midrule\endfirsthead\toprule\textbf{Attributo} & \textbf{Tipo fisico} & \textbf{NULL} & \textbf{Default}\\\midrule\endhead
\midrule\multicolumn{4}{r}{\footnotesize Continua nella pagina successiva}\\\endfoot\bottomrule\endlastfoot


{\ttfamily\small id} & {\ttfamily\small bigint} & no & {\ttfamily\small --}\\\addlinespace[3pt]

{\ttfamily\small session\_\allowbreak{}id} & {\ttfamily\small uuid} & no & {\ttfamily\small --}\\\addlinespace[3pt]

{\ttfamily\small changed\_\allowbreak{}at} & {\ttfamily\small timestamp with time zone} & no & {\ttfamily\small clock\_\allowbreak{}timestamp(\allowbreak{})}\\\addlinespace[3pt]

{\ttfamily\small previous\_\allowbreak{}schedule} & {\ttfamily\small timestamp without time zone} & no & {\ttfamily\small --}\\\addlinespace[3pt]

{\ttfamily\small next\_\allowbreak{}schedule} & {\ttfamily\small timestamp without time zone} & no & {\ttfamily\small --}\\\addlinespace[3pt]

{\ttfamily\small previous\_\allowbreak{}status} & {\ttfamily\small text} & no & {\ttfamily\small --}\\\addlinespace[3pt]

{\ttfamily\small next\_\allowbreak{}status} & {\ttfamily\small text} & no & {\ttfamily\small --}\\\addlinespace[3pt]

{\ttfamily\small previous\_\allowbreak{}attendance} & {\ttfamily\small jsonb} & no & {\ttfamily\small --}\\*

\end{xltabular}}

{\ttfamily\small PRIMARY KEY (\allowbreak{}id)}\par

{\ttfamily\small FOREIGN KEY (\allowbreak{}session\_\allowbreak{}id) REFERENCES app.\allowbreak{}course\_\allowbreak{}session(\allowbreak{}id)}\par

\Needspace{10\baselineskip}
\subsection{app.session\_menu\_snapshot}

Classificazione: operativa. Nome, descrizione e dosi acquisite per la sessione.

{\small \renewcommand{\arraystretch}{1.2}

\begin{xltabular}{\linewidth}{>{\hsize=1.120\hsize\linewidth=\hsize\raggedright\arraybackslash}X>{\hsize=1.200\hsize\linewidth=\hsize\raggedright\arraybackslash}X>{\hsize=0.400\hsize\linewidth=\hsize\raggedright\arraybackslash}X>{\hsize=1.280\hsize\linewidth=\hsize\raggedright\arraybackslash}X}
\caption{Attributi di app.session\_menu\_snapshot}\label{tab:dictionary-23}\\
\toprule

\textbf{Attributo} & \textbf{Tipo fisico} & \textbf{NULL} & \textbf{Default}\\\midrule\endfirsthead\toprule\textbf{Attributo} & \textbf{Tipo fisico} & \textbf{NULL} & \textbf{Default}\\\midrule\endhead
\midrule\multicolumn{4}{r}{\footnotesize Continua nella pagina successiva}\\\endfoot\bottomrule\endlastfoot


{\ttfamily\small session\_\allowbreak{}id} & {\ttfamily\small uuid} & no & {\ttfamily\small --}\\\addlinespace[3pt]

{\ttfamily\small recipe\_\allowbreak{}id} & {\ttfamily\small uuid} & no & {\ttfamily\small --}\\\addlinespace[3pt]

{\ttfamily\small recipe\_\allowbreak{}name} & {\ttfamily\small text} & no & {\ttfamily\small --}\\\addlinespace[3pt]

{\ttfamily\small description} & {\ttfamily\small text} & no & {\ttfamily\small --}\\\addlinespace[3pt]

{\ttfamily\small ingredients} & {\ttfamily\small jsonb} & no & {\ttfamily\small --}\\\addlinespace[3pt]

{\ttfamily\small captured\_\allowbreak{}at} & {\ttfamily\small timestamp with time zone} & no & {\ttfamily\small clock\_\allowbreak{}timestamp(\allowbreak{})}\\\addlinespace[3pt]

{\ttfamily\small provenance} & {\ttfamily\small text} & no & {\ttfamily\small --}\\*

\end{xltabular}}

{\ttfamily\small PRIMARY KEY (\allowbreak{}session\_\allowbreak{}id,\allowbreak{} recipe\_\allowbreak{}id)}\par

{\ttfamily\small FOREIGN KEY (\allowbreak{}session\_\allowbreak{}id,\allowbreak{} recipe\_\allowbreak{}id) REFERENCES app.\allowbreak{}session\_\allowbreak{}recipe(\allowbreak{}session\_\allowbreak{}id,\allowbreak{} recipe\_\allowbreak{}id) ON DELETE CASCADE}\par

\Needspace{10\baselineskip}
\subsection{app.session\_outcome}

Classificazione: storica. Esiti manuali conservati come storico; lo stato corrente è derivato dal calendario.

{\small \renewcommand{\arraystretch}{1.2}

\begin{xltabular}{\linewidth}{>{\hsize=1.120\hsize\linewidth=\hsize\raggedright\arraybackslash}X>{\hsize=1.200\hsize\linewidth=\hsize\raggedright\arraybackslash}X>{\hsize=0.400\hsize\linewidth=\hsize\raggedright\arraybackslash}X>{\hsize=1.280\hsize\linewidth=\hsize\raggedright\arraybackslash}X}
\caption{Attributi di app.session\_outcome}\label{tab:dictionary-24}\\
\toprule

\textbf{Attributo} & \textbf{Tipo fisico} & \textbf{NULL} & \textbf{Default}\\\midrule\endfirsthead\toprule\textbf{Attributo} & \textbf{Tipo fisico} & \textbf{NULL} & \textbf{Default}\\\midrule\endhead
\midrule\multicolumn{4}{r}{\footnotesize Continua nella pagina successiva}\\\endfoot\bottomrule\endlastfoot


{\ttfamily\small session\_\allowbreak{}id} & {\ttfamily\small uuid} & no & {\ttfamily\small --}\\\addlinespace[3pt]

{\ttfamily\small recorded\_\allowbreak{}at} & {\ttfamily\small timestamp with time zone} & no & {\ttfamily\small clock\_\allowbreak{}timestamp(\allowbreak{})}\\\addlinespace[3pt]

{\ttfamily\small recipe\_\allowbreak{}ids} & {\ttfamily\small uuid[]} & no & {\ttfamily\small --}\\\addlinespace[3pt]

{\ttfamily\small enrollment\_\allowbreak{}ids} & {\ttfamily\small uuid[]} & sì & {\ttfamily\small --}\\\addlinespace[3pt]

{\ttfamily\small note} & {\ttfamily\small text} & no & {\ttfamily\small --}\\*

\end{xltabular}}

{\ttfamily\small PRIMARY KEY (\allowbreak{}session\_\allowbreak{}id)}\par

{\ttfamily\small FOREIGN KEY (\allowbreak{}session\_\allowbreak{}id) REFERENCES app.\allowbreak{}course\_\allowbreak{}session(\allowbreak{}id)}\par

\Needspace{10\baselineskip}
\subsection{app.session\_preparation\_revision}

Classificazione: storica. Revisioni della preparazione conservate come storico; il fabbisogno corrente usa menu e adesioni valide.

{\small \renewcommand{\arraystretch}{1.2}

\begin{xltabular}{\linewidth}{>{\hsize=1.120\hsize\linewidth=\hsize\raggedright\arraybackslash}X>{\hsize=1.200\hsize\linewidth=\hsize\raggedright\arraybackslash}X>{\hsize=0.400\hsize\linewidth=\hsize\raggedright\arraybackslash}X>{\hsize=1.280\hsize\linewidth=\hsize\raggedright\arraybackslash}X}
\caption{Attributi di app.session\_preparation\_revision}\label{tab:dictionary-25}\\
\toprule

\textbf{Attributo} & \textbf{Tipo fisico} & \textbf{NULL} & \textbf{Default}\\\midrule\endfirsthead\toprule\textbf{Attributo} & \textbf{Tipo fisico} & \textbf{NULL} & \textbf{Default}\\\midrule\endhead
\midrule\multicolumn{4}{r}{\footnotesize Continua nella pagina successiva}\\\endfoot\bottomrule\endlastfoot


{\ttfamily\small session\_\allowbreak{}id} & {\ttfamily\small uuid} & no & {\ttfamily\small --}\\\addlinespace[3pt]

{\ttfamily\small revision} & {\ttfamily\small integer} & no & {\ttfamily\small --}\\\addlinespace[3pt]

{\ttfamily\small created\_\allowbreak{}at} & {\ttfamily\small timestamp with time zone} & no & {\ttfamily\small clock\_\allowbreak{}timestamp(\allowbreak{})}\\\addlinespace[3pt]

{\ttfamily\small reason} & {\ttfamily\small text} & no & {\ttfamily\small --}\\\addlinespace[3pt]

{\ttfamily\small preparation} & {\ttfamily\small jsonb} & no & {\ttfamily\small --}\\\addlinespace[3pt]

{\ttfamily\small confirmed\_\allowbreak{}enrollments} & {\ttfamily\small uuid[]} & no & {\ttfamily\small --}\\\addlinespace[3pt]

{\ttfamily\small active} & {\ttfamily\small boolean} & no & {\ttfamily\small true}\\*

\end{xltabular}}

{\ttfamily\small PRIMARY KEY (\allowbreak{}session\_\allowbreak{}id,\allowbreak{} revision)}\par

{\ttfamily\small FOREIGN KEY (\allowbreak{}session\_\allowbreak{}id) REFERENCES app.\allowbreak{}course\_\allowbreak{}session(\allowbreak{}id)}\par

\Needspace{10\baselineskip}
\subsection{app.session\_recipe}

Classificazione: operativa. Ricetta associata al menu di una pratica.

{\small \renewcommand{\arraystretch}{1.2}

\begin{xltabular}{\linewidth}{>{\hsize=1.120\hsize\linewidth=\hsize\raggedright\arraybackslash}X>{\hsize=1.200\hsize\linewidth=\hsize\raggedright\arraybackslash}X>{\hsize=0.400\hsize\linewidth=\hsize\raggedright\arraybackslash}X>{\hsize=1.280\hsize\linewidth=\hsize\raggedright\arraybackslash}X}
\caption{Attributi di app.session\_recipe}\label{tab:dictionary-26}\\
\toprule

\textbf{Attributo} & \textbf{Tipo fisico} & \textbf{NULL} & \textbf{Default}\\\midrule\endfirsthead\toprule\textbf{Attributo} & \textbf{Tipo fisico} & \textbf{NULL} & \textbf{Default}\\\midrule\endhead
\midrule\multicolumn{4}{r}{\footnotesize Continua nella pagina successiva}\\\endfoot\bottomrule\endlastfoot


{\ttfamily\small session\_\allowbreak{}id} & {\ttfamily\small uuid} & no & {\ttfamily\small --}\\\addlinespace[3pt]

{\ttfamily\small recipe\_\allowbreak{}id} & {\ttfamily\small uuid} & no & {\ttfamily\small --}\\\addlinespace[3pt]

{\ttfamily\small portions\_\allowbreak{}per\_\allowbreak{}participant} & {\ttfamily\small numeric} & no & {\ttfamily\small 1}\\*

\end{xltabular}}

{\ttfamily\small PRIMARY KEY (\allowbreak{}session\_\allowbreak{}id,\allowbreak{} recipe\_\allowbreak{}id)}\par

{\ttfamily\small FOREIGN KEY (\allowbreak{}recipe\_\allowbreak{}id) REFERENCES app.\allowbreak{}recipe(\allowbreak{}id)}\par

{\ttfamily\small FOREIGN KEY (\allowbreak{}session\_\allowbreak{}id) REFERENCES app.\allowbreak{}course\_\allowbreak{}session(\allowbreak{}id) ON DELETE CASCADE}\par

\subsection{Dizionario dei vincoli}\label{sec:constraints}

Il catalogo contiene 205 vincoli: 22 chiavi primarie, 24 chiavi esterne, 9 unicità, 46 CHECK e 104 NOT NULL. Chiavi e nullabilità sono descritte insieme alle relazioni nella \cref{sec:logical-schema}. Per evitare di ripetere il codice SQL, il dizionario seguente esprime in parole le condizioni dei CHECK ed è raggruppato per tabella. I nomi e le espressioni complete sono in \sqlsource. I controlli tra più righe sono approfonditi nella \cref{sec:sql-operations}.

\Needspace{9\baselineskip}
\subsubsection{account}
{\small\renewcommand{\arraystretch}{1.25}
\begin{xltabular}{\linewidth}{>{\hsize=0.680\hsize\linewidth=\hsize\raggedright\arraybackslash}X>{\hsize=1.320\hsize\linewidth=\hsize\raggedright\arraybackslash}X}
\caption{Vincoli di account}\label{tab:dictionary-27}\\
\toprule
\textbf{Vincolo} & \textbf{Condizione e significato}\\\midrule\endfirsthead\toprule\textbf{Vincolo} & \textbf{Condizione e significato}\\\midrule\endhead
\midrule\multicolumn{2}{r}{\footnotesize Continua nella pagina successiva}\\\endfoot\bottomrule\endlastfoot

{\ttfamily\small account\_\allowbreak{}password\_\allowbreak{}hash\_\allowbreak{}format} & L'hash inizia con il prefisso \$pbkdf2-sha256\$600000\$; il CHECK verifica il prefisso del formato, senza validare il contenuto successivo.\\\addlinespace[3pt]
{\ttfamily\small account\_\allowbreak{}role\_\allowbreak{}valid} & Il ruolo ammette soltanto CHEF o PARTICIPANT.\\\addlinespace[3pt]
{\ttfamily\small account\_\allowbreak{}username\_\allowbreak{}format} & Lo username è minuscolo, lungo da 3 a 50 caratteri e contiene soltanto lettere da a a z, cifre, punto, trattino basso o trattino.\\*
\end{xltabular}}

\Needspace{9\baselineskip}
\subsubsection{account\_password\_reset}
{\small\renewcommand{\arraystretch}{1.25}
\begin{xltabular}{\linewidth}{>{\hsize=0.680\hsize\linewidth=\hsize\raggedright\arraybackslash}X>{\hsize=1.320\hsize\linewidth=\hsize\raggedright\arraybackslash}X}
\caption{Vincoli di account\_password\_reset}\label{tab:dictionary-28}\\
\toprule
\textbf{Vincolo} & \textbf{Condizione e significato}\\\midrule\endfirsthead\toprule\textbf{Vincolo} & \textbf{Condizione e significato}\\\midrule\endhead
\midrule\multicolumn{2}{r}{\footnotesize Continua nella pagina successiva}\\\endfoot\bottomrule\endlastfoot

{\ttfamily\small account\_\allowbreak{}password\_\allowbreak{}reset\_\allowbreak{}expiry\_\allowbreak{}valid} & La scadenza del token è successiva alla sua creazione.\\\addlinespace[3pt]
{\ttfamily\small account\_\allowbreak{}password\_\allowbreak{}reset\_\allowbreak{}used\_\allowbreak{}after\_\allowbreak{}creation} & La data di utilizzo può essere assente; quando presente, non precede la creazione del token.\\*
\end{xltabular}}

\Needspace{9\baselineskip}
\subsubsection{auth\_session}
{\small\renewcommand{\arraystretch}{1.25}
\begin{xltabular}{\linewidth}{>{\hsize=0.680\hsize\linewidth=\hsize\raggedright\arraybackslash}X>{\hsize=1.320\hsize\linewidth=\hsize\raggedright\arraybackslash}X}
\caption{Vincoli di auth\_session}\label{tab:dictionary-29}\\
\toprule
\textbf{Vincolo} & \textbf{Condizione e significato}\\\midrule\endfirsthead\toprule\textbf{Vincolo} & \textbf{Condizione e significato}\\\midrule\endhead
\midrule\multicolumn{2}{r}{\footnotesize Continua nella pagina successiva}\\\endfoot\bottomrule\endlastfoot

{\ttfamily\small auth\_\allowbreak{}session\_\allowbreak{}expiry\_\allowbreak{}valid} & La scadenza della sessione di autenticazione è successiva alla sua creazione.\\*
\end{xltabular}}

\Needspace{9\baselineskip}
\subsubsection{category}
{\small\renewcommand{\arraystretch}{1.25}
\begin{xltabular}{\linewidth}{>{\hsize=0.680\hsize\linewidth=\hsize\raggedright\arraybackslash}X>{\hsize=1.320\hsize\linewidth=\hsize\raggedright\arraybackslash}X}
\caption{Vincoli di category}\label{tab:dictionary-30}\\
\toprule
\textbf{Vincolo} & \textbf{Condizione e significato}\\\midrule\endfirsthead\toprule\textbf{Vincolo} & \textbf{Condizione e significato}\\\midrule\endhead
\midrule\multicolumn{2}{r}{\footnotesize Continua nella pagina successiva}\\\endfoot\bottomrule\endlastfoot

{\ttfamily\small category\_\allowbreak{}description\_\allowbreak{}length\_\allowbreak{}check} & La descrizione, esclusi gli spazi iniziali e finali, contiene al massimo 500 caratteri; è ammessa una descrizione vuota.\\\addlinespace[3pt]
{\ttfamily\small category\_\allowbreak{}name\_\allowbreak{}check} & Il nome non è vuoto dopo la rimozione degli spazi iniziali e finali.\\\addlinespace[3pt]
{\ttfamily\small category\_\allowbreak{}name\_\allowbreak{}length\_\allowbreak{}check} & Il nome, esclusi gli spazi iniziali e finali, contiene da 3 a 80 caratteri.\\*
\end{xltabular}}

\Needspace{9\baselineskip}
\subsubsection{chef\_invitation}
{\small\renewcommand{\arraystretch}{1.25}
\begin{xltabular}{\linewidth}{>{\hsize=0.680\hsize\linewidth=\hsize\raggedright\arraybackslash}X>{\hsize=1.320\hsize\linewidth=\hsize\raggedright\arraybackslash}X}
\caption{Vincoli di chef\_invitation}\label{tab:dictionary-31}\\
\toprule
\textbf{Vincolo} & \textbf{Condizione e significato}\\\midrule\endfirsthead\toprule\textbf{Vincolo} & \textbf{Condizione e significato}\\\midrule\endhead
\midrule\multicolumn{2}{r}{\footnotesize Continua nella pagina successiva}\\\endfoot\bottomrule\endlastfoot

{\ttfamily\small chef\_\allowbreak{}invitation\_\allowbreak{}expiry\_\allowbreak{}valid} & La scadenza dell'invito è successiva alla sua creazione.\\\addlinespace[3pt]
{\ttfamily\small chef\_\allowbreak{}invitation\_\allowbreak{}used\_\allowbreak{}after\_\allowbreak{}creation} & La data di utilizzo può essere assente; quando presente, non precede la creazione dell'invito.\\*
\end{xltabular}}

\Needspace{9\baselineskip}
\subsubsection{chef\_profile}
{\small\renewcommand{\arraystretch}{1.25}
\begin{xltabular}{\linewidth}{>{\hsize=0.680\hsize\linewidth=\hsize\raggedright\arraybackslash}X>{\hsize=1.320\hsize\linewidth=\hsize\raggedright\arraybackslash}X}
\caption{Vincoli di chef\_profile}\label{tab:dictionary-32}\\
\toprule
\textbf{Vincolo} & \textbf{Condizione e significato}\\\midrule\endfirsthead\toprule\textbf{Vincolo} & \textbf{Condizione e significato}\\\midrule\endhead
\midrule\multicolumn{2}{r}{\footnotesize Continua nella pagina successiva}\\\endfoot\bottomrule\endlastfoot

{\ttfamily\small chef\_\allowbreak{}profile\_\allowbreak{}display\_\allowbreak{}name\_\allowbreak{}check} & Il nome visualizzato non è vuoto dopo la rimozione degli spazi iniziali e finali.\\*
\end{xltabular}}

\Needspace{9\baselineskip}
\subsubsection{course}
{\small\renewcommand{\arraystretch}{1.25}
\begin{xltabular}{\linewidth}{>{\hsize=0.680\hsize\linewidth=\hsize\raggedright\arraybackslash}X>{\hsize=1.320\hsize\linewidth=\hsize\raggedright\arraybackslash}X}
\caption{Vincoli di course}\label{tab:dictionary-33}\\
\toprule
\textbf{Vincolo} & \textbf{Condizione e significato}\\\midrule\endfirsthead\toprule\textbf{Vincolo} & \textbf{Condizione e significato}\\\midrule\endhead
\midrule\multicolumn{2}{r}{\footnotesize Continua nella pagina successiva}\\\endfoot\bottomrule\endlastfoot

{\ttfamily\small course\_\allowbreak{}description\_\allowbreak{}valid} & La descrizione non è vuota dopo la rimozione degli spazi iniziali e finali.\\\addlinespace[3pt]
{\ttfamily\small course\_\allowbreak{}ingredient\_\allowbreak{}margin\_\allowbreak{}valid} & Il margine degli ingredienti è compreso tra 0 e 25, estremi inclusi.\\\addlinespace[3pt]
{\ttfamily\small course\_\allowbreak{}no\_\allowbreak{}ingredient\_\allowbreak{}margin} & Il margine degli ingredienti è uguale a 0; questo vincolo restringe ulteriormente l'intervallo ammesso dall'altro CHECK sul margine.\\\addlinespace[3pt]
{\ttfamily\small course\_\allowbreak{}recurrence\_\allowbreak{}valid} & L'intervallo tra le sessioni è compreso tra 1 e 31 giorni, estremi inclusi.\\\addlinespace[3pt]
{\ttfamily\small course\_\allowbreak{}session\_\allowbreak{}count\_\allowbreak{}valid} & Il numero delle sessioni è compreso tra 1 e 30, estremi inclusi.\\\addlinespace[3pt]
{\ttfamily\small course\_\allowbreak{}status\_\allowbreak{}valid} & Lo stato ammette soltanto DRAFT, PUBLISHED, COMPLETED o CANCELLED.\\\addlinespace[3pt]
{\ttfamily\small course\_\allowbreak{}title\_\allowbreak{}valid} & Il titolo contiene da 3 a 30 caratteri, senza spazi iniziali o finali. Sono esclusi ritorno carrello, nuova riga, tabulazione verticale, avanzamento pagina e separatori Unicode U+0085, U+2028 e U+2029.\\*
\end{xltabular}}

\Needspace{9\baselineskip}
\subsubsection{course\_session}
{\small\renewcommand{\arraystretch}{1.25}
\begin{xltabular}{\linewidth}{>{\hsize=0.680\hsize\linewidth=\hsize\raggedright\arraybackslash}X>{\hsize=1.320\hsize\linewidth=\hsize\raggedright\arraybackslash}X}
\caption{Vincoli di course\_session}\label{tab:dictionary-34}\\
\toprule
\textbf{Vincolo} & \textbf{Condizione e significato}\\\midrule\endfirsthead\toprule\textbf{Vincolo} & \textbf{Condizione e significato}\\\midrule\endhead
\midrule\multicolumn{2}{r}{\footnotesize Continua nella pagina successiva}\\\endfoot\bottomrule\endlastfoot

{\ttfamily\small course\_\allowbreak{}session\_\allowbreak{}attendance\_\allowbreak{}cutoff\_\allowbreak{}valid} & Per ONLINE la scadenza delle adesioni è assente; per PRACTICAL è presente e precede strettamente l'inizio programmato.\\\addlinespace[3pt]
{\ttfamily\small course\_\allowbreak{}session\_\allowbreak{}duration\_\allowbreak{}valid} & La durata è compresa tra 1 e 720 minuti, estremi inclusi.\\\addlinespace[3pt]
{\ttfamily\small course\_\allowbreak{}session\_\allowbreak{}mode\_\allowbreak{}check} & La modalità ammette soltanto ONLINE o PRACTICAL.\\\addlinespace[3pt]
{\ttfamily\small course\_\allowbreak{}session\_\allowbreak{}mode\_\allowbreak{}fields\_\allowbreak{}valid} & Per ONLINE il link è presente, mentre sede e capienza sono assenti. Il link rispetta il formato HTTPS senza distinzione tra maiuscole e minuscole: host con primo carattere alfanumerico e successivi alfanumerici, punti o trattini; porta numerica facoltativa; eventuale parte finale introdotta da /, ? o \# e senza caratteri di spaziatura. Per PRACTICAL il link è assente, la sede è presente, non vuota e senza spazi iniziali o finali, e la capienza è presente e positiva.\\\addlinespace[3pt]
{\ttfamily\small course\_\allowbreak{}session\_\allowbreak{}sequence\_\allowbreak{}number\_\allowbreak{}check} & Il numero progressivo della sessione è maggiore di 0.\\\addlinespace[3pt]
{\ttfamily\small course\_\allowbreak{}session\_\allowbreak{}status\_\allowbreak{}check} & Lo stato memorizzato ammette soltanto SCHEDULED, COMPLETED o CANCELLED.\\*
\end{xltabular}}

\Needspace{9\baselineskip}
\subsubsection{enrollment}
{\small\renewcommand{\arraystretch}{1.25}
\begin{xltabular}{\linewidth}{>{\hsize=0.680\hsize\linewidth=\hsize\raggedright\arraybackslash}X>{\hsize=1.320\hsize\linewidth=\hsize\raggedright\arraybackslash}X}
\caption{Vincoli di enrollment}\label{tab:dictionary-35}\\
\toprule
\textbf{Vincolo} & \textbf{Condizione e significato}\\\midrule\endfirsthead\toprule\textbf{Vincolo} & \textbf{Condizione e significato}\\\midrule\endhead
\midrule\multicolumn{2}{r}{\footnotesize Continua nella pagina successiva}\\\endfoot\bottomrule\endlastfoot

{\ttfamily\small enrollment\_\allowbreak{}status\_\allowbreak{}valid} & Lo stato dell'iscrizione ammette soltanto ACTIVE o WITHDRAWN.\\\addlinespace[3pt]
{\ttfamily\small enrollment\_\allowbreak{}withdrawal\_\allowbreak{}valid} & Per ACTIVE la data di ritiro è assente; per WITHDRAWN è presente.\\*
\end{xltabular}}

\Needspace{9\baselineskip}
\subsubsection{ingredient}
{\small\renewcommand{\arraystretch}{1.25}
\begin{xltabular}{\linewidth}{>{\hsize=0.680\hsize\linewidth=\hsize\raggedright\arraybackslash}X>{\hsize=1.320\hsize\linewidth=\hsize\raggedright\arraybackslash}X}
\caption{Vincoli di ingredient}\label{tab:dictionary-36}\\
\toprule
\textbf{Vincolo} & \textbf{Condizione e significato}\\\midrule\endfirsthead\toprule\textbf{Vincolo} & \textbf{Condizione e significato}\\\midrule\endhead
\midrule\multicolumn{2}{r}{\footnotesize Continua nella pagina successiva}\\\endfoot\bottomrule\endlastfoot

{\ttfamily\small ingredient\_\allowbreak{}name\_\allowbreak{}check} & Il nome non è vuoto dopo la rimozione degli spazi iniziali e finali.\\\addlinespace[3pt]
{\ttfamily\small ingredient\_\allowbreak{}unit\_\allowbreak{}valid} & L'unità ammette soltanto GRAM, MILLILITER, PIECE o TEASPOON.\\*
\end{xltabular}}

\Needspace{9\baselineskip}
\subsubsection{notification}
{\small\renewcommand{\arraystretch}{1.25}
\begin{xltabular}{\linewidth}{>{\hsize=0.680\hsize\linewidth=\hsize\raggedright\arraybackslash}X>{\hsize=1.320\hsize\linewidth=\hsize\raggedright\arraybackslash}X}
\caption{Vincoli di notification}\label{tab:dictionary-37}\\
\toprule
\textbf{Vincolo} & \textbf{Condizione e significato}\\\midrule\endfirsthead\toprule\textbf{Vincolo} & \textbf{Condizione e significato}\\\midrule\endhead
\midrule\multicolumn{2}{r}{\footnotesize Continua nella pagina successiva}\\\endfoot\bottomrule\endlastfoot

{\ttfamily\small notification\_\allowbreak{}message\_\allowbreak{}valid} & Il messaggio contiene da 3 a 2000 caratteri, senza spazi iniziali o finali.\\\addlinespace[3pt]
{\ttfamily\small notification\_\allowbreak{}reason\_\allowbreak{}valid} & Il motivo ammette soltanto SCHEDULE\_CHANGED, SESSION\_CANCELLED, SESSION\_ADDED o GENERAL.\\\addlinespace[3pt]
{\ttfamily\small notification\_\allowbreak{}title\_\allowbreak{}valid} & Il titolo contiene da 3 a 120 caratteri, senza spazi iniziali o finali.\\*
\end{xltabular}}

\Needspace{9\baselineskip}
\subsubsection{participant\_profile}
{\small\renewcommand{\arraystretch}{1.25}
\begin{xltabular}{\linewidth}{>{\hsize=0.680\hsize\linewidth=\hsize\raggedright\arraybackslash}X>{\hsize=1.320\hsize\linewidth=\hsize\raggedright\arraybackslash}X}
\caption{Vincoli di participant\_profile}\label{tab:dictionary-38}\\
\toprule
\textbf{Vincolo} & \textbf{Condizione e significato}\\\midrule\endfirsthead\toprule\textbf{Vincolo} & \textbf{Condizione e significato}\\\midrule\endhead
\midrule\multicolumn{2}{r}{\footnotesize Continua nella pagina successiva}\\\endfoot\bottomrule\endlastfoot

{\ttfamily\small participant\_\allowbreak{}profile\_\allowbreak{}display\_\allowbreak{}name\_\allowbreak{}check} & Il nome visualizzato non è vuoto dopo la rimozione degli spazi iniziali e finali.\\*
\end{xltabular}}

\Needspace{9\baselineskip}
\subsubsection{practical\_attendance}
{\small\renewcommand{\arraystretch}{1.25}
\begin{xltabular}{\linewidth}{>{\hsize=0.680\hsize\linewidth=\hsize\raggedright\arraybackslash}X>{\hsize=1.320\hsize\linewidth=\hsize\raggedright\arraybackslash}X}
\caption{Vincoli di practical\_attendance}\label{tab:dictionary-39}\\
\toprule
\textbf{Vincolo} & \textbf{Condizione e significato}\\\midrule\endfirsthead\toprule\textbf{Vincolo} & \textbf{Condizione e significato}\\\midrule\endhead
\midrule\multicolumn{2}{r}{\footnotesize Continua nella pagina successiva}\\\endfoot\bottomrule\endlastfoot

{\ttfamily\small practical\_\allowbreak{}attendance\_\allowbreak{}status\_\allowbreak{}valid} & La risposta alla sessione pratica ammette soltanto CONFIRMED o DECLINED.\\*
\end{xltabular}}

\Needspace{9\baselineskip}
\subsubsection{recipe}
{\small\renewcommand{\arraystretch}{1.25}
\begin{xltabular}{\linewidth}{>{\hsize=0.680\hsize\linewidth=\hsize\raggedright\arraybackslash}X>{\hsize=1.320\hsize\linewidth=\hsize\raggedright\arraybackslash}X}
\caption{Vincoli di recipe}\label{tab:dictionary-40}\\
\toprule
\textbf{Vincolo} & \textbf{Condizione e significato}\\\midrule\endfirsthead\toprule\textbf{Vincolo} & \textbf{Condizione e significato}\\\midrule\endhead
\midrule\multicolumn{2}{r}{\footnotesize Continua nella pagina successiva}\\\endfoot\bottomrule\endlastfoot

{\ttfamily\small recipe\_\allowbreak{}base\_\allowbreak{}servings\_\allowbreak{}check} & Il numero delle porzioni di base è maggiore di 0.\\\addlinespace[3pt]
{\ttfamily\small recipe\_\allowbreak{}description\_\allowbreak{}check} & La descrizione non è vuota dopo la rimozione degli spazi iniziali e finali.\\\addlinespace[3pt]
{\ttfamily\small recipe\_\allowbreak{}name\_\allowbreak{}check} & Il nome non è vuoto dopo la rimozione degli spazi iniziali e finali.\\\addlinespace[3pt]
{\ttfamily\small recipe\_\allowbreak{}single\_\allowbreak{}portion} & Il numero delle porzioni di base è uguale a 1.\\*
\end{xltabular}}

\Needspace{9\baselineskip}
\subsubsection{recipe\_ingredient}
{\small\renewcommand{\arraystretch}{1.25}
\begin{xltabular}{\linewidth}{>{\hsize=0.680\hsize\linewidth=\hsize\raggedright\arraybackslash}X>{\hsize=1.320\hsize\linewidth=\hsize\raggedright\arraybackslash}X}
\caption{Vincoli di recipe\_ingredient}\label{tab:dictionary-41}\\
\toprule
\textbf{Vincolo} & \textbf{Condizione e significato}\\\midrule\endfirsthead\toprule\textbf{Vincolo} & \textbf{Condizione e significato}\\\midrule\endhead
\midrule\multicolumn{2}{r}{\footnotesize Continua nella pagina successiva}\\\endfoot\bottomrule\endlastfoot

{\ttfamily\small recipe\_\allowbreak{}ingredient\_\allowbreak{}quantity\_\allowbreak{}check} & La quantità dell'ingrediente è maggiore di 0.\\\addlinespace[3pt]
{\ttfamily\small recipe\_\allowbreak{}quantity\_\allowbreak{}precision} & La quantità è maggiore di 0 e minore di 1000000000000; il valore coincide con il suo arrotondamento a 3 decimali, quindi non sono ammesse cifre decimali significative oltre la terza.\\*
\end{xltabular}}

\Needspace{9\baselineskip}
\subsubsection{session\_menu\_snapshot}
{\small\renewcommand{\arraystretch}{1.25}
\begin{xltabular}{\linewidth}{>{\hsize=0.680\hsize\linewidth=\hsize\raggedright\arraybackslash}X>{\hsize=1.320\hsize\linewidth=\hsize\raggedright\arraybackslash}X}
\caption{Vincoli di session\_menu\_snapshot}\label{tab:dictionary-42}\\
\toprule
\textbf{Vincolo} & \textbf{Condizione e significato}\\\midrule\endfirsthead\toprule\textbf{Vincolo} & \textbf{Condizione e significato}\\\midrule\endhead
\midrule\multicolumn{2}{r}{\footnotesize Continua nella pagina successiva}\\\endfoot\bottomrule\endlastfoot

{\ttfamily\small session\_\allowbreak{}menu\_\allowbreak{}snapshot\_\allowbreak{}provenance\_\allowbreak{}check} & La provenienza della copia ammette soltanto CAPTURED o LEGACY\_RECONSTRUCTED.\\*
\end{xltabular}}

\Needspace{9\baselineskip}
\subsubsection{session\_outcome}
{\small\renewcommand{\arraystretch}{1.25}
\begin{xltabular}{\linewidth}{>{\hsize=0.680\hsize\linewidth=\hsize\raggedright\arraybackslash}X>{\hsize=1.320\hsize\linewidth=\hsize\raggedright\arraybackslash}X}
\caption{Vincoli di session\_outcome}\label{tab:dictionary-43}\\
\toprule
\textbf{Vincolo} & \textbf{Condizione e significato}\\\midrule\endfirsthead\toprule\textbf{Vincolo} & \textbf{Condizione e significato}\\\midrule\endhead
\midrule\multicolumn{2}{r}{\footnotesize Continua nella pagina successiva}\\\endfoot\bottomrule\endlastfoot

{\ttfamily\small session\_\allowbreak{}outcome\_\allowbreak{}note\_\allowbreak{}check} & La nota, esclusi gli spazi iniziali e finali, contiene al massimo 2000 caratteri; è ammessa una nota vuota.\\*
\end{xltabular}}

\Needspace{9\baselineskip}
\subsubsection{session\_preparation\_revision}
{\small\renewcommand{\arraystretch}{1.25}
\begin{xltabular}{\linewidth}{>{\hsize=0.680\hsize\linewidth=\hsize\raggedright\arraybackslash}X>{\hsize=1.320\hsize\linewidth=\hsize\raggedright\arraybackslash}X}
\caption{Vincoli di session\_preparation\_revision}\label{tab:dictionary-44}\\
\toprule
\textbf{Vincolo} & \textbf{Condizione e significato}\\\midrule\endfirsthead\toprule\textbf{Vincolo} & \textbf{Condizione e significato}\\\midrule\endhead
\midrule\multicolumn{2}{r}{\footnotesize Continua nella pagina successiva}\\\endfoot\bottomrule\endlastfoot

{\ttfamily\small session\_\allowbreak{}preparation\_\allowbreak{}revision\_\allowbreak{}reason\_\allowbreak{}check} & La motivazione, esclusi gli spazi iniziali e finali, contiene da 1 a 2000 caratteri.\\\addlinespace[3pt]
{\ttfamily\small session\_\allowbreak{}preparation\_\allowbreak{}revision\_\allowbreak{}revision\_\allowbreak{}check} & Il numero della revisione è maggiore di 0.\\*
\end{xltabular}}

\Needspace{9\baselineskip}
\subsubsection{session\_recipe}
{\small\renewcommand{\arraystretch}{1.25}
\begin{xltabular}{\linewidth}{>{\hsize=0.680\hsize\linewidth=\hsize\raggedright\arraybackslash}X>{\hsize=1.320\hsize\linewidth=\hsize\raggedright\arraybackslash}X}
\caption{Vincoli di session\_recipe}\label{tab:dictionary-45}\\
\toprule
\textbf{Vincolo} & \textbf{Condizione e significato}\\\midrule\endfirsthead\toprule\textbf{Vincolo} & \textbf{Condizione e significato}\\\midrule\endhead
\midrule\multicolumn{2}{r}{\footnotesize Continua nella pagina successiva}\\\endfoot\bottomrule\endlastfoot

{\ttfamily\small session\_\allowbreak{}recipe\_\allowbreak{}portions\_\allowbreak{}valid} & Le porzioni per partecipante sono maggiori di 0.\\\addlinespace[3pt]
{\ttfamily\small session\_\allowbreak{}recipe\_\allowbreak{}single\_\allowbreak{}portion} & Le porzioni per partecipante sono uguali a 1.\\*
\end{xltabular}}

Le dosi positive con massimo tre decimali sono verificate senza arrotondamento implicito. Il limite del conteggio salvato nel corso non impone da solo il numero di righe delle sessioni: il caso d’uso di creazione richiede inizialmente 2--30 sessioni.

\clearpage\section{Componenti applicativi e classi di design}

Il progetto usa Java 25, JavaFX 26.0.2, PostgreSQL 18.4, JDBC, HikariCP e Flyway. L’architettura comprende quattro livelli: presentation, application, domain e infrastructure. I controller raccolgono gli input e mostrano gli esiti; i casi d'uso coordinano le operazioni tramite interfacce; il dominio contiene valori e regole indipendenti da GUI e JDBC; l'infrastruttura implementa l'accesso ai dati.

ApplicationRuntime costruisce manualmente i componenti. Il package application.bootstrap dipende dall'infrastruttura perché ne collega le implementazioni: è un'eccezione esplicita, controllata dai test di architettura. Non vengono usati un ORM, un contenitore di dependency injection o un backend HTTP.

Il diagramma di design è organizzato in tavole complementari. Le \cref{fig:diagram-8,fig:diagram-9,fig:diagram-10,fig:diagram-11} dettagliano creazione del corso, valori del calendario, adesioni e pianificazione; le tavole successive collegano le classi principali delle altre aree operative. Sono classi, record e interfacce presenti nel codice. La suddivisione rende leggibili membri e relazioni; uno stesso nome ripetuto indica lo stesso tipo. Il concetto Corso non corrisponde a una classe Java persistente: è rappresentato da comandi, modelli di lettura e tabelle SQL.

La notazione dei membri usa + per pubblico, - per privato e tilde per visibilità di package. Nei record, i componenti sono valori immutabili con accessor; non sono campi pubblici mutabili. La freccia triangolare tratteggiata rappresenta la realizzazione di un'interfaccia; le altre frecce indicano riferimenti o dipendenze.


Le tavole presentano le classi principali delle aree operative. Sono omesse le classi di supporto grafico e i dettagli interni non necessari per comprendere i collegamenti. Le ellissi nei metodi omettono soltanto l’elenco dei parametri. Le frecce continue indicano riferimenti conservati negli oggetti; quelle tratteggiate con punta aperta indicano l’uso di un tipo; quelle con triangolo vuoto indicano la realizzazione di un’interfaccia. Le ripetizioni di uno stesso nome indicano la stessa classe nelle diverse tavole.

\clearpage\subsection{Creazione del corso e coordinamento}\begin{center}

\begin{tikzpicture}[x=1cm,y=1cm,box/.style={draw=ink,rounded corners=1pt,fill=white,text width=6.95cm,minimum height=2.7cm,inner sep=3mm,align=left,font=\small},assoc/.style={-{Stealth[length=2mm]},draw=ink},real/.style={dashed,-{Triangle[open,length=2.4mm]},draw=ink}]

\node[box] (n0) at (0.0,0) {\textbf{\mbox{CourseWorkspaceController}}\\\rule{6.65cm}{0.3pt}\\- creationForm:\newline\hspace*{2mm}CourseCreationForm\\- coordinator:\newline\hspace*{2mm}\mbox{CourseWorkspaceCoordinator}\\\rule{6.65cm}{0.3pt}\\- createCourse(): void\\+ showCourseCreation(): void};

\node[box] (n1) at (8.2,0) {\textbf{\mbox{CourseWorkspaceCoordinator}}\\\rule{6.65cm}{0.3pt}\\- busy: boolean\\- workspaceData:\newline\hspace*{2mm}CourseWorkspaceData\\\rule{6.65cm}{0.3pt}\\+ runAsync(...): void\\+ loadCourses(UUID): void};

\node[box] (n2) at (0.0,-4.54) {\textbf{CourseCreationForm}\\\rule{6.65cm}{0.3pt}\\- formModel:\newline\hspace*{2mm}CourseFormModel\\- createCourse: CreateCourse\\\rule{6.65cm}{0.3pt}\\\textasciitilde{} createCourse(): void\\\textasciitilde{} hasDraft(): boolean};

\node[box] (n3) at (8.2,-4.54) {\textbf{\mbox{CourseWorkspaceActions}}\\\rule{6.65cm}{0.3pt}\\(nessun attributo)\\\rule{6.65cm}{0.3pt}\\+ runAsync(...): void\\+ loadCourses(UUID): void};

\node[box] (n4) at (0.0,-8.72) {\textbf{CreateCourse}\\\rule{6.65cm}{0.3pt}\\- courseRepository:\newline\hspace*{2mm}CourseRepository\\\rule{6.65cm}{0.3pt}\\+ execute(\newline\hspace*{2mm}AuthenticatedUser, NewCourseInput): UUID};

\node[box] (n5) at (8.2,-8.72) {\textbf{CourseFormModel}\\\rule{6.65cm}{0.3pt}\\- sessions: List<SessionDraft>\\- activeSessionCount: int\\\rule{6.65cm}{0.3pt}\\+ setSessionCount(int): void\\+ toInput(...): NewCourseInput};

\node[box] (n6) at (0.0,-12.89) {\textbf{CourseRepository}\\\rule{6.65cm}{0.3pt}\\(nessun attributo)\\\rule{6.65cm}{0.3pt}\\+ createAndPublish(\newline\hspace*{2mm}CourseCreation): UUID\\+ assignRecipe(UUID, UUID, UUID): void};

\node[box] (n7) at (8.2,-12.89) {\textbf{\mbox{JdbcCourseRepository}}\\\rule{6.65cm}{0.3pt}\\- dataSource: DataSource\\\rule{6.65cm}{0.3pt}\\+ createAndPublish(\newline\hspace*{2mm}CourseCreation): UUID\\- insertCourse(\newline\hspace*{2mm}Connection, CourseCreation): UUID};

\draw[assoc] (n0.south) -- (n2.north);

\draw[assoc] (n0.east) -- (n1.west);

\draw[assoc] (n2.south east) -- node[midway,fill=white,inner sep=1pt,font=\footnotesize]{1} (n5.north west);

\draw[assoc] (n2.east) -- (n3.west);

\draw[assoc] (n2.south) -- (n4.north);

\draw[real] (n1.south) -- (n3.north);

\draw[assoc] (n4.south) -- (n6.north);

\draw[real] (n7.west) -- (n6.east);

\end{tikzpicture}
\captionof{figure}{Creazione del corso e coordinamento}\label{fig:diagram-8}
\end{center}

La tavola mostra i collegamenti principali del modulo; gli attributi e i metodi mostrati sono una selezione dei membri reali.

\clearpage\subsection{Valori del corso e calendario}\begin{center}

\begin{tikzpicture}[x=1cm,y=1cm,box/.style={draw=ink,rounded corners=1pt,fill=white,text width=6.95cm,minimum height=2.7cm,inner sep=3mm,align=left,font=\small},assoc/.style={-{Stealth[length=2mm]},draw=ink},real/.style={dashed,-{Triangle[open,length=2.4mm]},draw=ink}]

\node[box] (n0) at (0.0,0) {\textbf{NewCourseInput}\\\rule{6.65cm}{0.3pt}\\categoryId: UUID\\title: String};

\node[box] (n1) at (8.2,0) {\textbf{CourseCreation}\\\rule{6.65cm}{0.3pt}\\sessionToken: UUID\\title: CourseTitle};

\node[box] (n2) at (0.0,-3.88) {\textbf{SessionDetails}\\\rule{6.65cm}{0.3pt}\\mode: SessionMode\\durationMinutes: int\\\rule{6.65cm}{0.3pt}\\+ online(int, URI): SessionDetails\\+ practical(\newline\hspace*{2mm}int, String, int): SessionDetails};

\node[box] (n3) at (8.2,-3.88) {\textbf{CourseSchedule}\\\rule{6.65cm}{0.3pt}\\- startDate: LocalDate\\- recurrenceDays: int\\\rule{6.65cm}{0.3pt}\\+ of(...): CourseSchedule\\+ sessions(): List<SessionPlan>};

\node[box] (n4) at (0.0,-7.87) {\textbf{CourseTitle}\\\rule{6.65cm}{0.3pt}\\value: String\\\rule{6.65cm}{0.3pt}\\+ of(String): CourseTitle\\+ isSingleLine(String): boolean};

\node[box] (n5) at (8.2,-7.87) {\textbf{SessionPlan}\\\rule{6.65cm}{0.3pt}\\sequenceNumber: int\\scheduledAt: LocalDateTime\\\rule{6.65cm}{0.3pt}\\+ mode(): SessionMode};

\draw[assoc] (n1.south) -- node[midway,fill=white,inner sep=1pt,font=\footnotesize]{1} (n3.north);

\draw[assoc] (n0.south) -- node[midway,fill=white,inner sep=1pt,font=\footnotesize]{0..*} (n2.north);

\draw[assoc] (n3.south) -- node[midway,fill=white,inner sep=1pt,font=\footnotesize]{2..30} (n5.north);

\draw[assoc] (n5.south west) -- node[midway,fill=white,inner sep=1pt,font=\footnotesize]{1} (n2.north east);

\draw[assoc] (n1.west) -- (4.12,0) -- node[near end,fill=white,inner sep=1pt,font=\footnotesize]{1} (4.12,-6.1) -- (4.12,-7.87) -- (n4.east);

\end{tikzpicture}
\captionof{figure}{Valori del corso e calendario}\label{fig:diagram-9}
\end{center}

La tavola mostra i collegamenti principali del modulo; gli attributi e i metodi mostrati sono una selezione dei membri reali.

Il caso d'uso CreateCourse della \cref{fig:diagram-8} riceve NewCourseInput, costruisce il comando CourseCreation della \cref{fig:diagram-9} e lo inoltra alla porta CourseRepository. Il comando contiene il titolo validato e il calendario, evitando di passare alla persistenza valori non controllati.

\clearpage\subsection{Gestione delle adesioni}\begin{center}

\begin{tikzpicture}[x=1cm,y=1cm,box/.style={draw=ink,rounded corners=1pt,fill=white,text width=6.95cm,minimum height=2.7cm,inner sep=3mm,align=left,font=\small},assoc/.style={-{Stealth[length=2mm]},draw=ink},real/.style={dashed,-{Triangle[open,length=2.4mm]},draw=ink}]

\node[box] (n0) at (0.0,0) {\textbf{\mbox{ParticipantWorkspaceController}}\\\rule{6.65cm}{0.3pt}\\- data:\newline\hspace*{2mm}\mbox{ParticipantWorkspaceData}\\- async:\newline\hspace*{2mm}\mbox{ParticipantWorkspaceAsync}\\\rule{6.65cm}{0.3pt}\\- changeAttendance(\newline\hspace*{2mm}\mbox{PracticalAttendanceSummary},\newline\hspace*{2mm}AttendanceStatus): void\\- \mbox{invalidateAttendance}(): void};

\node[box] (n1) at (8.2,0) {\textbf{\mbox{ParticipantWorkspaceAsync}}\\\rule{6.65cm}{0.3pt}\\- hooks: Hooks\\- retryRead: Runnable\\\rule{6.65cm}{0.3pt}\\\textasciitilde{} read(...): void\\\textasciitilde{} write(...): void};

\node[box] (n2) at (0.0,-5.1) {\textbf{\mbox{ParticipantWorkspaceData}}\\\rule{6.65cm}{0.3pt}\\- user:\newline\hspace*{2mm}AuthenticatedUser\\- \mbox{attendanceManagement}:\newline\hspace*{2mm}\mbox{ManagePracticalAttendance}\\\rule{6.65cm}{0.3pt}\\\textasciitilde{} confirm(UUID): void\\\textasciitilde{} decline(UUID): void};

\node[box] (n3) at (8.2,-5.1) {\textbf{\mbox{ManagePracticalAttendance}}\\\rule{6.65cm}{0.3pt}\\- \mbox{enrollmentRepository}:\newline\hspace*{2mm}\mbox{EnrollmentRepository}\\\rule{6.65cm}{0.3pt}\\+ confirm(\newline\hspace*{2mm}AuthenticatedUser, UUID): void\\+ decline(\newline\hspace*{2mm}AuthenticatedUser, UUID): void};

\node[box] (n4) at (0.0,-9.64) {\textbf{\mbox{EnrollmentRepository}}\\\rule{6.65cm}{0.3pt}\\(nessun attributo)\\\rule{6.65cm}{0.3pt}\\+ setAttendance(\newline\hspace*{2mm}UUID, UUID, AttendanceStatus): void\\+ practicalSessions(\newline\hspace*{2mm}UUID, UUID): List<\mbox{PracticalAttendanceSummary}>};

\node[box] (n5) at (8.2,-9.64) {\textbf{\mbox{JdbcEnrollmentRepository}}\\\rule{6.65cm}{0.3pt}\\- dataSource: DataSource\\\rule{6.65cm}{0.3pt}\\+ setAttendance(\newline\hspace*{2mm}UUID, UUID, AttendanceStatus): void};

\draw[assoc] (n0.south) -- (n2.north);

\draw[assoc] (n0.east) -- (n1.west);

\draw[assoc] (n2.east) -- (n3.west);

\draw[assoc] (n3.south west) -- (n4.north east);

\draw[real] (n5.west) -- (n4.east);

\end{tikzpicture}
\captionof{figure}{Gestione delle adesioni}\label{fig:diagram-10}
\end{center}

La tavola mostra i collegamenti principali del modulo; gli attributi e i metodi mostrati sono una selezione dei membri reali.

\clearpage\subsection{Lettura della pianificazione}\begin{center}

\begin{tikzpicture}[x=1cm,y=1cm,box/.style={draw=ink,rounded corners=1pt,fill=white,text width=6.95cm,minimum height=2.7cm,inner sep=3mm,align=left,font=\small},assoc/.style={-{Stealth[length=2mm]},draw=ink},real/.style={dashed,-{Triangle[open,length=2.4mm]},draw=ink}]

\node[box] (n0) at (0.0,0) {\textbf{BrowseChefPlanning}\\\rule{6.65cm}{0.3pt}\\- repository:\newline\hspace*{2mm}\mbox{ChefPlanningRepository}\\\rule{6.65cm}{0.3pt}\\+ sessionPreparation(\newline\hspace*{2mm}AuthenticatedUser, UUID):\newline\hspace*{2mm}SessionPreparation\\+ requirements(\newline\hspace*{2mm}AuthenticatedUser, UUID): List<\mbox{IngredientRequirement}>};

\node[box] (n1) at (8.2,0) {\textbf{SessionPreparation}\\\rule{6.65cm}{0.3pt}\\sessionId: UUID\\\mbox{confirmedParticipants}: int\\totals: List<\mbox{IngredientRequirement}>};

\node[box] (n2) at (0.0,-6.0) {\textbf{\mbox{ChefPlanningRepository}}\\\rule{6.65cm}{0.3pt}\\(nessun attributo)\\\rule{6.65cm}{0.3pt}\\+ sessionPreparation(\newline\hspace*{2mm}UUID, UUID):\newline\hspace*{2mm}SessionPreparation\\+ requirements(\newline\hspace*{2mm}UUID, UUID): List<\mbox{IngredientRequirement}>};

\node[box] (n3) at (8.2,-6.0) {\textbf{\mbox{IngredientRequirement}}\\\rule{6.65cm}{0.3pt}\\ingredientName: String\\unit:\newline\hspace*{2mm}MeasurementUnit};

\node[box] (n4) at (0.0,-12.0) {\textbf{\mbox{JdbcChefPlanningRepository}}\\\rule{6.65cm}{0.3pt}\\- dataSource: DataSource\\\rule{6.65cm}{0.3pt}\\+ sessionPreparation(\newline\hspace*{2mm}UUID, UUID):\newline\hspace*{2mm}SessionPreparation\\+ requirements(\newline\hspace*{2mm}UUID, UUID): List<\mbox{IngredientRequirement}>};

\draw[assoc] (n0.south) -- (n2.north);

\draw[real] (n4.north) -- (n2.south);

\draw[dashed,-{Stealth[open,length=2mm]},draw=ink] (n0.east) -- (n1.west);

\draw[assoc] (n1.south) -- node[midway,fill=white,inner sep=1pt,font=\footnotesize]{0..* totals} (n3.north);

\draw[dashed,-{Stealth[open,length=2mm]},draw=ink] (n2.east) -- (n3.west);
\draw[dashed,-{Stealth[open,length=2mm]},draw=ink] (n2.north east) -- (n1.south west);
\end{tikzpicture}
\captionof{figure}{Lettura della pianificazione}\label{fig:diagram-11}
\end{center}

La tavola mostra i collegamenti principali del modulo; gli attributi e i metodi mostrati sono una selezione dei membri reali.

Le transazioni sono nei repository JDBC. Il pool disattiva la conferma automatica e ammette 10 connessioni con attesa massima di 10 secondi. Creazione e pubblicazione sono atomiche. L'esportazione dei fabbisogni usa REPEATABLE\_READ per mantenere letture coerenti.

\clearpage\begin{landscape}
\subsection{Accesso, registrazione e gestione dell’account}
\begin{center}
\begin{tikzpicture}[x=1cm,y=1cm,
dcbox/.style={draw=ink,fill=white,text width=6.8cm,minimum height=2.4cm,inner sep=2.5mm,align=left,font=\fontsize{9}{10.4}\selectfont},
dcassoc/.style={-{Stealth[length=1.8mm]},draw=ink,line width=.4pt},
dcdep/.style={dashed,-{Stealth[open,length=1.8mm]},draw=ink,line width=.4pt},
dcreal/.style={dashed,-{Triangle[open,length=2mm]},draw=ink,line width=.4pt}]

\node[dcbox] (n0) at (0.0,0.0) {\textit{\guillemotleft classe\guillemotright}\\\textbf{Login\allowbreak{}Controller}\\\rule{6.65cm}{.3pt}\\- authenticate\allowbreak{}User: Authenticate\allowbreak{}User\\- account\allowbreak{}Lifecycle: Account\allowbreak{}Lifecycle\\- submit(): void};
\node[dcbox] (n1) at (8.1,0.0) {\textit{\guillemotleft classe\guillemotright}\\\textbf{Authenticate\allowbreak{}User}\\\rule{6.65cm}{.3pt}\\- account\allowbreak{}Repository: Account\allowbreak{}Repository\\+ execute(...): Authenticated\allowbreak{}User};
\node[dcbox] (n2) at (16.2,0.0) {\textit{\guillemotleft interfaccia\guillemotright}\\\textbf{Account\allowbreak{}Repository}\\\rule{6.65cm}{.3pt}\\+ authenticate(...): Optional<Authenticated\allowbreak{}User>};
\node[dcbox] (n3) at (0.0,-3.25) {\textit{\guillemotleft classe\guillemotright}\\\textbf{Account\allowbreak{}Dialogs}\\\rule{6.65cm}{.3pt}\\+ show\allowbreak{}Registration(...): Optional<Registration>\\+ show\allowbreak{}Password\allowbreak{}Change(...): Optional<Password\allowbreak{}Change>};
\node[dcbox] (n4) at (8.1,-3.25) {\textit{\guillemotleft classe\guillemotright}\\\textbf{Account\allowbreak{}Lifecycle}\\\rule{6.65cm}{.3pt}\\- password\allowbreak{}Hasher: Password\allowbreak{}Hasher\\+ register\allowbreak{}Chef(...): Authenticated\allowbreak{}User\\+ change\allowbreak{}Password(...): void\\+ logout(...): void};
\node[dcbox] (n5) at (16.2,-3.25) {\textit{\guillemotleft interfaccia\guillemotright}\\\textbf{Account\allowbreak{}Lifecycle\allowbreak{}Repository}\\\rule{6.65cm}{.3pt}\\+ register\allowbreak{}Participant(...): Authenticated\allowbreak{}User\\+ register\allowbreak{}Chef(...): Authenticated\allowbreak{}User\\+ change\allowbreak{}Password(...): void\\+ logout(...): void};
\node[dcbox] (n6) at (0.0,-6.5) {\textit{\guillemotleft classe\guillemotright}\\\textbf{Scene\allowbreak{}Navigator}\\\rule{6.65cm}{.3pt}\\- runtime: Application\allowbreak{}Runtime\\- change\allowbreak{}Password(...): void\\- logout(...): void};
\node[dcbox] (n7) at (8.1,-6.5) {\textit{\guillemotleft interfaccia\guillemotright}\\\textbf{Password\allowbreak{}Hasher}\\\rule{6.65cm}{.3pt}\\+ hash(char[]): String\\+ matches(...): boolean};
\node[dcbox] (n8) at (16.2,-6.5) {\textit{\guillemotleft classe\guillemotright}\\\textbf{Jdbc\allowbreak{}Account\allowbreak{}Lifecycle\allowbreak{}Repository}\\\rule{6.65cm}{.3pt}\\- data\allowbreak{}Source: Data\allowbreak{}Source\\- password\allowbreak{}Hasher: Pbkdf2Password\allowbreak{}Hasher};
\node[dcbox] (n9) at (0.0,-9.75) {\textit{\guillemotleft record\guillemotright}\\\textbf{Authenticated\allowbreak{}User}\\\rule{6.65cm}{.3pt}\\session\allowbreak{}Token: UUID\\username: Username\\role: Role};
\node[dcbox] (n10) at (8.1,-9.75) {\textit{\guillemotleft classe\guillemotright}\\\textbf{Pbkdf2Password\allowbreak{}Hasher}\\\rule{6.65cm}{.3pt}\\+ hash(char[]): String\\+ authentication\allowbreak{}Proof(...): byte[]};
\node[dcbox] (n11) at (16.2,-9.75) {\textit{\guillemotleft classe\guillemotright}\\\textbf{Jdbc\allowbreak{}Account\allowbreak{}Repository}\\\rule{6.65cm}{.3pt}\\- data\allowbreak{}Source: Data\allowbreak{}Source\\- password\allowbreak{}Hasher: Pbkdf2Password\allowbreak{}Hasher};
\draw[dcassoc] (n0.east) -- (n1.west);
\draw[dcassoc] (n1.east) -- (n2.west);
\draw[dcdep] (n0.south) -- (n3.north);
\draw[dcassoc] (n0.south) -- (0.0,-1.625) -| (n4.north);
\draw[dcdep] (n6.north) -- (n3.south);
\draw[dcdep] (n6.east) -- (4.05,-6.5) -- (4.05,-3.25) -- (n4.west);
\draw[dcassoc] (n4.east) -- (n5.west);
\draw[dcassoc] (n4.south) -- (n7.north);
\draw[dcreal] (n8.north) -- (n5.south);
\draw[dcreal] (n10.north) -- (n7.south);
\draw[dcreal] (n11.east) -- ++(.55,0) |- (n2.east);
\draw[dcassoc] (n8.south) -- (16.2,-8.125) -- (12.05,-8.125) -- (12.05,-9.75) -- (n10.east);
\draw[dcdep] (n6.south) -- (n9.north);
\end{tikzpicture}
\captionof{figure}{Accesso, registrazione e gestione dell’account}\label{fig:design-identity}
\end{center}
{\small La sessione autenticata viene passata ai casi d’uso. SceneNavigator richiama AccountLifecycle attraverso ApplicationRuntime; i dialoghi raccolgono i dati, senza accedere direttamente al database.}
\end{landscape}

\clearpage\begin{landscape}
\subsection{Ricettario, ingredienti e categorie}
\begin{center}
\begin{tikzpicture}[x=1cm,y=1cm,
dcbox/.style={draw=ink,fill=white,text width=6.8cm,minimum height=2.4cm,inner sep=2.5mm,align=left,font=\fontsize{9}{10.4}\selectfont},
dcassoc/.style={-{Stealth[length=1.8mm]},draw=ink,line width=.4pt},
dcdep/.style={dashed,-{Stealth[open,length=1.8mm]},draw=ink,line width=.4pt},
dcreal/.style={dashed,-{Triangle[open,length=2mm]},draw=ink,line width=.4pt}]

\node[dcbox] (n0) at (0.0,0.0) {\textit{\guillemotleft classe\guillemotright}\\\textbf{Catalog\allowbreak{}Workspace\allowbreak{}Controller}\\\rule{6.65cm}{.3pt}\\- browse: Browse\allowbreak{}Catalog\\- manage: Manage\allowbreak{}Recipe\allowbreak{}Catalog\\+ activate(): void};
\node[dcbox] (n1) at (8.1,0.0) {\textit{\guillemotleft classe\guillemotright}\\\textbf{Browse\allowbreak{}Catalog}\\\rule{6.65cm}{.3pt}\\- catalog\allowbreak{}Repository: Catalog\allowbreak{}Repository\\+ recipes(...): Page<Recipe\allowbreak{}Summary>};
\node[dcbox] (n2) at (16.2,0.0) {\textit{\guillemotleft interfaccia\guillemotright}\\\textbf{Catalog\allowbreak{}Repository}\\\rule{6.65cm}{.3pt}\\+ categories(...): List<Category\allowbreak{}Summary>\\+ recipes(...): Page<Recipe\allowbreak{}Summary>\\+ create\allowbreak{}Category(...): Category\allowbreak{}Summary};
\node[dcbox] (n3) at (0.0,-3.25) {\textit{\guillemotleft classe\guillemotright}\\\textbf{Manage\allowbreak{}Recipe\allowbreak{}Catalog}\\\rule{6.65cm}{.3pt}\\- repository: Recipe\allowbreak{}Catalog\allowbreak{}Management\allowbreak{}Repository\\+ save\allowbreak{}Recipe(...): UUID\\+ save\allowbreak{}Ingredient(...): UUID};
\node[dcbox] (n4) at (8.1,-3.25) {\textit{\guillemotleft classe\guillemotright}\\\textbf{Create\allowbreak{}Category}\\\rule{6.65cm}{.3pt}\\- catalog\allowbreak{}Repository: Catalog\allowbreak{}Repository\\+ create(...): Category\allowbreak{}Summary};
\node[dcbox] (n5) at (16.2,-3.25) {\textit{\guillemotleft classe\guillemotright}\\\textbf{Jdbc\allowbreak{}Catalog\allowbreak{}Repository}\\\rule{6.65cm}{.3pt}\\- data\allowbreak{}Source: Data\allowbreak{}Source\\+ recipes(...): Page<Recipe\allowbreak{}Summary>};
\node[dcbox] (n6) at (0.0,-6.5) {\textit{\guillemotleft interfaccia\guillemotright}\\\textbf{Recipe\allowbreak{}Catalog\allowbreak{}Management\allowbreak{}Repository}\\\rule{6.65cm}{.3pt}\\+ save\allowbreak{}Recipe(...): UUID\\+ save\allowbreak{}Ingredient(...): UUID};
\node[dcbox] (n7) at (8.1,-6.5) {\textit{\guillemotleft record\guillemotright}\\\textbf{Recipe\allowbreak{}Input}\\\rule{6.65cm}{.3pt}\\name: String\\ingredients: List<Recipe\allowbreak{}Ingredient\allowbreak{}Input>};
\node[dcbox] (n8) at (16.2,-6.5) {\textit{\guillemotleft record\guillemotright}\\\textbf{Ingredient\allowbreak{}Input}\\\rule{6.65cm}{.3pt}\\name: String\\unit: Measurement\allowbreak{}Unit};
\node[dcbox] (n9) at (0.0,-9.75) {\textit{\guillemotleft classe\guillemotright}\\\textbf{Jdbc\allowbreak{}Recipe\allowbreak{}Catalog\allowbreak{}Management\allowbreak{}Repository}\\\rule{6.65cm}{.3pt}\\- data\allowbreak{}Source: Data\allowbreak{}Source\\+ save\allowbreak{}Recipe(...): UUID};
\node[dcbox] (n10) at (8.1,-9.75) {\textit{\guillemotleft record\guillemotright}\\\textbf{Recipe\allowbreak{}Ingredient\allowbreak{}Input}\\\rule{6.65cm}{.3pt}\\ingredient\allowbreak{}Id: UUID\\quantity: Big\allowbreak{}Decimal};
\node[dcbox] (n11) at (16.2,-9.75) {\textit{\guillemotleft enumerazione\guillemotright}\\\textbf{Measurement\allowbreak{}Unit}\\\rule{6.65cm}{.3pt}\\GRAM\\MILLILITER\\PIECE\\TEASPOON};
\draw[dcassoc] (n0.east) -- (n1.west);
\draw[dcassoc] (n1.east) -- (n2.west);
\draw[dcassoc] (n0.south) -- (n3.north);
\draw[dcassoc] (n0.south) -- (0.0,-1.625) -| (n4.north);
\draw[dcassoc] (n4.north) -- (8.1,-1.625) -| (n2.south);
\draw[dcreal] (n5.north) -- (n2.south);
\draw[dcassoc] (n3.south) -- (n6.north);
\draw[dcreal] (n9.north) -- (n6.south);
\draw[dcdep] (n6.east) -- (n7.west);
\draw[dcdep] (n6.south east) -- (4.05,-8.125) -- (12.15,-8.125) -- (12.15,-6.5) -- (n8.west);
\draw[dcassoc] (n7.south) -- (n10.north);
\draw[dcassoc] (n8.south) -- (n11.north);
\end{tikzpicture}
\captionof{figure}{Ricettario, ingredienti e categorie}\label{fig:design-catalog}
\end{center}
{\small La lettura del catalogo e la sua modifica usano contratti separati. RecipeInput contiene le dosi per una persona; CreateCategory condivide CatalogRepository con BrowseCatalog.}
\end{landscape}

\clearpage\begin{landscape}
\subsection{Calendario dei corsi e menu delle sessioni}
\begin{center}
\begin{tikzpicture}[x=1cm,y=1cm,
dcbox/.style={draw=ink,fill=white,text width=6.8cm,minimum height=2.4cm,inner sep=2.5mm,align=left,font=\fontsize{9}{10.4}\selectfont},
dcassoc/.style={-{Stealth[length=1.8mm]},draw=ink,line width=.4pt},
dcdep/.style={dashed,-{Stealth[open,length=1.8mm]},draw=ink,line width=.4pt},
dcreal/.style={dashed,-{Triangle[open,length=2mm]},draw=ink,line width=.4pt}]

\node[dcbox] (n0) at (0.0,0.0) {\textit{\guillemotleft classe\guillemotright}\\\textbf{Course\allowbreak{}Workspace\allowbreak{}Controller}\\\rule{6.65cm}{.3pt}\\- schedule\allowbreak{}Commands: Course\allowbreak{}Schedule\allowbreak{}Commands\\- recipe\allowbreak{}Menu: Course\allowbreak{}Recipe\allowbreak{}Menu};
\node[dcbox] (n1) at (8.1,0.0) {\textit{\guillemotleft classe\guillemotright}\\\textbf{Course\allowbreak{}Schedule\allowbreak{}Commands}\\\rule{6.65cm}{.3pt}\\- manage\allowbreak{}Course\allowbreak{}Schedule: Manage\allowbreak{}Course\allowbreak{}Schedule\\- schedule\allowbreak{}Editor: Course\allowbreak{}Schedule\allowbreak{}Editor\\\textasciitilde{} cancel\allowbreak{}Session(): void};
\node[dcbox] (n2) at (16.2,0.0) {\textit{\guillemotleft classe\guillemotright}\\\textbf{Course\allowbreak{}Schedule\allowbreak{}Editor}\\\rule{6.65cm}{.3pt}\\\textasciitilde{} configure\allowbreak{}Schedule\allowbreak{}Editor(): void\\\textasciitilde{} populate\allowbreak{}Schedule\allowbreak{}Editor(...): void};
\node[dcbox] (n3) at (0.0,-3.25) {\textit{\guillemotleft classe\guillemotright}\\\textbf{Course\allowbreak{}Recipe\allowbreak{}Menu}\\\rule{6.65cm}{.3pt}\\- assign\allowbreak{}Recipe\allowbreak{}To\allowbreak{}Session: Assign\allowbreak{}Recipe\allowbreak{}To\allowbreak{}Session\\- lifecycle: Manage\allowbreak{}Session\allowbreak{}Lifecycle\\\textasciitilde{} assign\allowbreak{}Recipe(...): void};
\node[dcbox] (n4) at (8.1,-3.25) {\textit{\guillemotleft classe\guillemotright}\\\textbf{Manage\allowbreak{}Course\allowbreak{}Schedule}\\\rule{6.65cm}{.3pt}\\- repository: Schedule\allowbreak{}Repository\\+ add(...): UUID\\+ reschedule(...): void\\+ cancel(...): void};
\node[dcbox] (n5) at (16.2,-3.25) {\textit{\guillemotleft interfaccia\guillemotright}\\\textbf{Schedule\allowbreak{}Repository}\\\rule{6.65cm}{.3pt}\\+ add(...): UUID\\+ reschedule(...): void\\+ cancel(...): void};
\node[dcbox] (n6) at (0.0,-6.5) {\textit{\guillemotleft classe\guillemotright}\\\textbf{Assign\allowbreak{}Recipe\allowbreak{}To\allowbreak{}Session}\\\rule{6.65cm}{.3pt}\\- course\allowbreak{}Repository: Course\allowbreak{}Repository\\+ execute(...): void\\+ remove(...): void};
\node[dcbox] (n7) at (8.1,-6.5) {\textit{\guillemotleft classe\guillemotright}\\\textbf{Manage\allowbreak{}Session\allowbreak{}Lifecycle}\\\rule{6.65cm}{.3pt}\\- repository: Session\allowbreak{}Lifecycle\allowbreak{}Repository\\+ menu(...): List<Saved\allowbreak{}Recipe>\\+ refresh(...): void};
\node[dcbox] (n8) at (16.2,-6.5) {\textit{\guillemotleft classe\guillemotright}\\\textbf{Jdbc\allowbreak{}Schedule\allowbreak{}Repository}\\\rule{6.65cm}{.3pt}\\- data\allowbreak{}Source: Data\allowbreak{}Source\\+ reschedule(...): void};
\node[dcbox] (n9) at (0.0,-9.75) {\textit{\guillemotleft interfaccia\guillemotright}\\\textbf{Course\allowbreak{}Repository}\\\rule{6.65cm}{.3pt}\\+ assign\allowbreak{}Recipe(...): void\\+ remove\allowbreak{}Recipe(...): void};
\node[dcbox] (n10) at (8.1,-9.75) {\textit{\guillemotleft interfaccia\guillemotright}\\\textbf{Session\allowbreak{}Lifecycle\allowbreak{}Repository}\\\rule{6.65cm}{.3pt}\\+ menu(...): List<Saved\allowbreak{}Recipe>\\+ refresh(...): void};
\node[dcbox] (n11) at (16.2,-9.75) {\textit{\guillemotleft classe\guillemotright}\\\textbf{Jdbc\allowbreak{}Session\allowbreak{}Lifecycle\allowbreak{}Repository}\\\rule{6.65cm}{.3pt}\\- data\allowbreak{}Source: Data\allowbreak{}Source\\+ menu(...): List<Saved\allowbreak{}Recipe>};
\draw[dcassoc] (n0.east) -- (n1.west);
\draw[dcassoc] (n1.east) -- (n2.west);
\draw[dcassoc] (n0.south) -- (n3.north);
\draw[dcassoc] (n1.south) -- (n4.north);
\draw[dcassoc] (n4.east) -- (n5.west);
\draw[dcreal] (n8.north) -- (n5.south);
\draw[dcassoc] (n3.south) -- (n6.north);
\draw[dcassoc] (n3.south) -- (0.0,-4.875) -| (n7.north);
\draw[dcassoc] (n6.south) -- (n9.north);
\draw[dcassoc] (n7.south) -- (n10.north);
\draw[dcreal] (n11.west) -- (n10.east);
\end{tikzpicture}
\captionof{figure}{Calendario dei corsi e menu delle sessioni}\label{fig:design-calendar-menu}
\end{center}
{\small I comandi del calendario includono la bozza della comunicazione. L’associazione di una ricetta e la lettura o l’aggiornamento delle copie del menu seguono contratti distinti. CourseRepository e JdbcCourseRepository sono collegati nella tavola sulla creazione del corso.}
\end{landscape}

\clearpage\begin{landscape}
\subsection{Consultazione e iscrizioni del Partecipante}
\begin{center}
\begin{tikzpicture}[x=1cm,y=1cm,
dcbox/.style={draw=ink,fill=white,text width=6.8cm,minimum height=2.4cm,inner sep=2.5mm,align=left,font=\fontsize{9}{10.4}\selectfont},
dcassoc/.style={-{Stealth[length=1.8mm]},draw=ink,line width=.4pt},
dcdep/.style={dashed,-{Stealth[open,length=1.8mm]},draw=ink,line width=.4pt},
dcreal/.style={dashed,-{Triangle[open,length=2mm]},draw=ink,line width=.4pt}]

\node[dcbox] (n0) at (0.0,0.0) {\textit{\guillemotleft classe\guillemotright}\\\textbf{Participant\allowbreak{}Workspace\allowbreak{}Controller}\\\rule{6.65cm}{.3pt}\\- data: Participant\allowbreak{}Workspace\allowbreak{}Data\\- async: Participant\allowbreak{}Workspace\allowbreak{}Async};
\node[dcbox] (n1) at (8.1,0.0) {\textit{\guillemotleft classe\guillemotright}\\\textbf{Participant\allowbreak{}Workspace\allowbreak{}Data}\\\rule{6.65cm}{.3pt}\\- published\allowbreak{}Courses: Browse\allowbreak{}Published\allowbreak{}Courses\\\textasciitilde{} enroll(UUID): void\\\textasciitilde{} withdraw(UUID): void};
\node[dcbox] (n2) at (16.2,0.0) {\textit{\guillemotleft classe\guillemotright}\\\textbf{Browse\allowbreak{}Published\allowbreak{}Courses}\\\rule{6.65cm}{.3pt}\\- course\allowbreak{}Repository: Course\allowbreak{}Repository\\+ sessions(...): List<Published\allowbreak{}Course\allowbreak{}Session\allowbreak{}Summary>\\+ menu(...): List<Published\allowbreak{}Session\allowbreak{}Recipe\allowbreak{}Summary>};
\node[dcbox] (n3) at (0.0,-3.25) {\textit{\guillemotleft classe\guillemotright}\\\textbf{Browse\allowbreak{}Catalog}\\\rule{6.65cm}{.3pt}\\+ categories(...): List<Category\allowbreak{}Summary>};
\node[dcbox] (n4) at (8.1,-3.25) {\textit{\guillemotleft classe\guillemotright}\\\textbf{Browse\allowbreak{}Enrollments}\\\rule{6.65cm}{.3pt}\\- enrollment\allowbreak{}Repository: Enrollment\allowbreak{}Repository\\+ enrollments(...): Page<Enrollment\allowbreak{}Summary>};
\node[dcbox] (n5) at (16.2,-3.25) {\textit{\guillemotleft interfaccia\guillemotright}\\\textbf{Course\allowbreak{}Repository}\\\rule{6.65cm}{.3pt}\\+ published\allowbreak{}Sessions(...): List<Published\allowbreak{}Course\allowbreak{}Session\allowbreak{}Summary>\\+ published\allowbreak{}Menu(...): List<Published\allowbreak{}Session\allowbreak{}Recipe\allowbreak{}Summary>};
\node[dcbox] (n6) at (0.0,-6.5) {\textit{\guillemotleft classe\guillemotright}\\\textbf{Manage\allowbreak{}Enrollment}\\\rule{6.65cm}{.3pt}\\- enrollment\allowbreak{}Repository: Enrollment\allowbreak{}Repository\\+ enroll(...): void\\+ withdraw(...): void};
\node[dcbox] (n7) at (8.1,-6.5) {\textit{\guillemotleft interfaccia\guillemotright}\\\textbf{Enrollment\allowbreak{}Repository}\\\rule{6.65cm}{.3pt}\\+ enroll(...): void\\+ withdraw(...): void};
\node[dcbox] (n8) at (16.2,-6.5) {\textit{\guillemotleft classe\guillemotright}\\\textbf{Jdbc\allowbreak{}Course\allowbreak{}Repository}\\\rule{6.65cm}{.3pt}\\- data\allowbreak{}Source: Data\allowbreak{}Source};
\node[dcbox] (n9) at (0.0,-9.75) {\textit{\guillemotleft classe\guillemotright}\\\textbf{Browse\allowbreak{}Notifications}\\\rule{6.65cm}{.3pt}\\+ for\allowbreak{}Participant(...): List<Participant\allowbreak{}Notification>};
\node[dcbox] (n10) at (8.1,-9.75) {\textit{\guillemotleft classe\guillemotright}\\\textbf{Manage\allowbreak{}Notifications}\\\rule{6.65cm}{.3pt}\\+ mark\allowbreak{}Read(...): void};
\node[dcbox] (n11) at (16.2,-9.75) {\textit{\guillemotleft classe\guillemotright}\\\textbf{Jdbc\allowbreak{}Enrollment\allowbreak{}Repository}\\\rule{6.65cm}{.3pt}\\- data\allowbreak{}Source: Data\allowbreak{}Source\\+ enroll(...): void\\+ withdraw(...): void};
\draw[dcassoc] (n0.east) -- (n1.west);
\draw[dcassoc] (n1.east) -- (n2.west);
\draw[dcassoc] (n2.south) -- (n5.north);
\draw[dcassoc] (n1.south) -- (8.1,-1.625) -| (n3.north);
\draw[dcassoc] (n1.south) -- (n4.north);
\draw[dcassoc] (n1.south west) -- ++(-.5,-.2) -- (-4.1,-1.5) -- (-4.1,-6.5) -- (n6.west);
\draw[dcassoc] (n4.south) -- (n7.north);
\draw[dcassoc] (n6.east) -- (n7.west);
\draw[dcreal] (n8.north) -- (n5.south);
\draw[dcreal] (n11.north) -- (16.2,-8.125) -| (n7.south);
\draw[dcassoc] (n1.south west) -- ++(-.7,-.4) -- (-4.4,-1.7) -- (-4.4,-9.75) -- (n9.west);
\draw[dcassoc] (n1.south east) -- ++(.5,-.2) -- (12.1,-1.5) -- (12.1,-9.75) -- (n10.east);
\end{tikzpicture}
\captionof{figure}{Consultazione e iscrizioni del Partecipante}\label{fig:design-participant}
\end{center}
{\small ParticipantWorkspaceData unisce letture e operazioni della stessa area. La consultazione usa modelli di lettura; EnrollmentRepository resta separato da CourseRepository. Le adesioni sono dettagliate nella tavola dedicata, mentre BrowseNotifications e ManageNotifications sono ripresi nella tavola delle comunicazioni.}
\end{landscape}

\clearpage\begin{landscape}
\subsection{Comunicazioni e report mensile}
\begin{center}
\begin{tikzpicture}[x=1cm,y=1cm,
dcbox/.style={draw=ink,fill=white,text width=6.8cm,minimum height=2.4cm,inner sep=2.5mm,align=left,font=\fontsize{9}{10.4}\selectfont},
dcassoc/.style={-{Stealth[length=1.8mm]},draw=ink,line width=.4pt},
dcdep/.style={dashed,-{Stealth[open,length=1.8mm]},draw=ink,line width=.4pt},
dcreal/.style={dashed,-{Triangle[open,length=2mm]},draw=ink,line width=.4pt}]

\node[dcbox] (n0) at (0.0,0.0) {\textit{\guillemotleft classe\guillemotright}\\\textbf{Notification\allowbreak{}Workspace\allowbreak{}Controller}\\\rule{6.65cm}{.3pt}\\- browse\allowbreak{}Notifications: Browse\allowbreak{}Notifications\\- manage\allowbreak{}Notifications: Manage\allowbreak{}Notifications\\+ compose\allowbreak{}For\allowbreak{}Course(UUID): void};
\node[dcbox] (n1) at (8.1,0.0) {\textit{\guillemotleft classe\guillemotright}\\\textbf{Manage\allowbreak{}Notifications}\\\rule{6.65cm}{.3pt}\\- repository: Notification\allowbreak{}Repository\\+ send(...): UUID\\+ mark\allowbreak{}Read(...): void};
\node[dcbox] (n2) at (16.2,0.0) {\textit{\guillemotleft classe\guillemotright}\\\textbf{Browse\allowbreak{}Notifications}\\\rule{6.65cm}{.3pt}\\- repository: Notification\allowbreak{}Repository\\+ for\allowbreak{}Chef(...): List<Chef\allowbreak{}Notification>\\+ for\allowbreak{}Participant(...): List<Participant\allowbreak{}Notification>};
\node[dcbox] (n3) at (0.0,-3.25) {\textit{\guillemotleft record\guillemotright}\\\textbf{Notification\allowbreak{}Draft}\\\rule{6.65cm}{.3pt}\\title: String\\message: String};
\node[dcbox] (n4) at (8.1,-3.25) {\textit{\guillemotleft interfaccia\guillemotright}\\\textbf{Notification\allowbreak{}Repository}\\\rule{6.65cm}{.3pt}\\+ send(...): UUID\\+ find\allowbreak{}For\allowbreak{}Chef(...): List<Chef\allowbreak{}Notification>\\+ mark\allowbreak{}Read(...): void};
\node[dcbox] (n5) at (16.2,-3.25) {\textit{\guillemotleft classe\guillemotright}\\\textbf{Jdbc\allowbreak{}Notification\allowbreak{}Repository}\\\rule{6.65cm}{.3pt}\\- data\allowbreak{}Source: Data\allowbreak{}Source\\+ send(...): UUID};
\node[dcbox] (n6) at (0.0,-6.5) {\textit{\guillemotleft classe\guillemotright}\\\textbf{Planning\allowbreak{}Workspace\allowbreak{}Controller}\\\rule{6.65cm}{.3pt}\\- report\allowbreak{}Presenter: Monthly\allowbreak{}Report\allowbreak{}Presenter\\+ show\allowbreak{}Monthly\allowbreak{}Report(): void};
\node[dcbox] (n7) at (8.1,-6.5) {\textit{\guillemotleft classe\guillemotright}\\\textbf{Monthly\allowbreak{}Report\allowbreak{}Presenter}\\\rule{6.65cm}{.3pt}\\- browse\allowbreak{}Chef\allowbreak{}Planning: Browse\allowbreak{}Chef\allowbreak{}Planning\\- monthly\allowbreak{}Report\allowbreak{}View: Monthly\allowbreak{}Report\allowbreak{}View\\\textasciitilde{} load(): void};
\node[dcbox] (n8) at (16.2,-6.5) {\textit{\guillemotleft classe\guillemotright}\\\textbf{Browse\allowbreak{}Chef\allowbreak{}Planning}\\\rule{6.65cm}{.3pt}\\- repository: Chef\allowbreak{}Planning\allowbreak{}Repository\\+ monthly\allowbreak{}Report(...): Monthly\allowbreak{}Report};
\node[dcbox] (n9) at (0.0,-9.75) {\textit{\guillemotleft record\guillemotright}\\\textbf{Monthly\allowbreak{}Report}\\\rule{6.65cm}{.3pt}\\month: Year\allowbreak{}Month\\total\allowbreak{}Courses: long\\average\allowbreak{}Recipes: Optional<Big\allowbreak{}Decimal>};
\node[dcbox] (n10) at (8.1,-9.75) {\textit{\guillemotleft classe\guillemotright}\\\textbf{Monthly\allowbreak{}Report\allowbreak{}View}\\\rule{6.65cm}{.3pt}\\- session\allowbreak{}Chart: Bar\allowbreak{}Chart<Number, String>\\- recipe\allowbreak{}Chart: Bar\allowbreak{}Chart<Number, String>\\\textasciitilde{} render(Monthly\allowbreak{}Report): void};
\node[dcbox] (n11) at (16.2,-9.75) {\textit{\guillemotleft interfaccia\guillemotright}\\\textbf{Chef\allowbreak{}Planning\allowbreak{}Repository}\\\rule{6.65cm}{.3pt}\\+ monthly\allowbreak{}Report(...): Monthly\allowbreak{}Report};
\draw[dcassoc] (n0.east) -- (n1.west);
\draw[dcassoc] (n0.north) -- ++(0,.5) -| (n2.north);
\draw[dcassoc] (n1.south) -- (n4.north);
\draw[dcassoc] (n2.south) -- (16.2,-1.625) -| (n4.north);
\draw[dcdep] (n1.south) -- (8.1,-1.745) -| (n3.north);
\draw[dcreal] (n5.west) -- (n4.east);
\draw[dcassoc] (n6.east) -- (n7.west);
\draw[dcassoc] (n7.east) -- (n8.west);
\draw[dcassoc] (n7.south) -- (n10.north);
\draw[dcassoc] (n8.south) -- (n11.north);
\draw[dcdep] (n10.west) -- (n9.east);
\end{tikzpicture}
\captionof{figure}{Comunicazioni e report mensile}\label{fig:design-communications-report}
\end{center}
{\small Le comunicazioni vengono registrate e consultate nell’applicazione. Il report usa il medesimo BrowseChefPlanning della pianificazione, con ChefPlanningRepository e la sua implementazione JDBC già collegati nella tavola di lettura della pianificazione. MonthlyReportView costruisce i grafici a partire dal risultato immutabile.}
\end{landscape}

\clearpage\begin{landscape}
\subsection{Consultazione dei corsi Chef e agenda}
\begin{center}
\begin{tikzpicture}[x=1cm,y=1cm,
dcbox/.style={draw=ink,fill=white,text width=6.8cm,minimum height=2.4cm,inner sep=2.5mm,align=left,font=\fontsize{9}{10.4}\selectfont},
dcassoc/.style={-{Stealth[length=1.8mm]},draw=ink,line width=.4pt},
dcdep/.style={dashed,-{Stealth[open,length=1.8mm]},draw=ink,line width=.4pt},
dcreal/.style={dashed,-{Triangle[open,length=2mm]},draw=ink,line width=.4pt}]

\node[dcbox] (n0) at (0.0,0.0) {\textit{\guillemotleft classe\guillemotright}\\\textbf{Course\allowbreak{}Workspace\allowbreak{}Controller}\\\rule{6.65cm}{.3pt}\\- workspace\allowbreak{}Data: Course\allowbreak{}Workspace\allowbreak{}Data\\- browse\allowbreak{}Courses: Browse\allowbreak{}Courses};
\node[dcbox] (n1) at (8.1,0.0) {\textit{\guillemotleft classe\guillemotright}\\\textbf{Course\allowbreak{}Workspace\allowbreak{}Data}\\\rule{6.65cm}{.3pt}\\- browse\allowbreak{}Courses: Browse\allowbreak{}Courses\\\textasciitilde{} load\allowbreak{}Courses(UUID): void\\\textasciitilde{} load\allowbreak{}Sessions(...): void};
\node[dcbox] (n2) at (16.2,0.0) {\textit{\guillemotleft classe\guillemotright}\\\textbf{Browse\allowbreak{}Courses}\\\rule{6.65cm}{.3pt}\\- course\allowbreak{}Repository: Course\allowbreak{}Repository\\+ courses(...): Page<Course\allowbreak{}Summary>\\+ sessions(...): List<Course\allowbreak{}Session\allowbreak{}Summary>};
\node[dcbox] (n3) at (0.0,-3.25) {\textit{\guillemotleft classe\guillemotright}\\\textbf{Dashboard\allowbreak{}Workspace\allowbreak{}Controller}\\\rule{6.65cm}{.3pt}\\- overview\allowbreak{}Loader: Dashboard\allowbreak{}Overview\allowbreak{}Loader\\+ activate(): void};
\node[dcbox] (n4) at (8.1,-3.25) {\textit{\guillemotleft classe\guillemotright}\\\textbf{Dashboard\allowbreak{}Overview\allowbreak{}Loader}\\\rule{6.65cm}{.3pt}\\- browse\allowbreak{}Courses: Browse\allowbreak{}Courses\\\textasciitilde{} load(): Dashboard\allowbreak{}Snapshot};
\node[dcbox] (n5) at (16.2,-3.25) {\textit{\guillemotleft interfaccia\guillemotright}\\\textbf{Course\allowbreak{}Repository}\\\rule{6.65cm}{.3pt}\\+ find\allowbreak{}By\allowbreak{}Chef(...): Page<Course\allowbreak{}Summary>\\+ sessions(...): List<Course\allowbreak{}Session\allowbreak{}Summary>};
\node[dcbox] (n6) at (0.0,-6.5) {\textit{\guillemotleft classe\guillemotright}\\\textbf{Course\allowbreak{}Page\allowbreak{}Queries}\\\rule{6.65cm}{.3pt}\\- data\allowbreak{}Source: Data\allowbreak{}Source\\\textasciitilde{} by\allowbreak{}Chef(...): Page<Course\allowbreak{}Summary>};
\node[dcbox] (n7) at (8.1,-6.5) {\textit{\guillemotleft classe\guillemotright}\\\textbf{Jdbc\allowbreak{}Course\allowbreak{}Repository}\\\rule{6.65cm}{.3pt}\\- chef\allowbreak{}Courses: Chef\allowbreak{}Course\allowbreak{}Queries\\- course\allowbreak{}Pages: Course\allowbreak{}Page\allowbreak{}Queries};
\node[dcbox] (n8) at (16.2,-6.5) {\textit{\guillemotleft classe\guillemotright}\\\textbf{Chef\allowbreak{}Course\allowbreak{}Queries}\\\rule{6.65cm}{.3pt}\\- data\allowbreak{}Source: Data\allowbreak{}Source\\\textasciitilde{} find\allowbreak{}By\allowbreak{}Chef(...): List<Course\allowbreak{}Summary>\\\textasciitilde{} sessions(...): List<Course\allowbreak{}Session\allowbreak{}Summary>};
\draw[dcassoc] (n0.east) -- (n1.west);
\draw[dcassoc] (n1.east) -- (n2.west);
\draw[dcassoc] (n0.north) -- ++(0,.5) -| (n2.north);
\draw[dcassoc] (n3.east) -- (n4.west);
\draw[dcassoc] (n4.north) -- (8.1,-1.625) -- (12.15,-1.625) -- (12.15,0) -- (n2.west);
\draw[dcassoc] (n2.south) -- (n5.north);
\draw[dcreal] (n7.north) -- (8.1,-4.875) -- (12.15,-4.875) -- (12.15,-3.25) -- (n5.west);
\draw[dcassoc] (n7.west) -- (n6.east);
\draw[dcassoc] (n7.east) -- (n8.west);
\end{tikzpicture}
\captionof{figure}{Consultazione dei corsi Chef e agenda}\label{fig:design-chef-consultation}
\end{center}
{\small L’elenco filtrato e l’agenda riutilizzano BrowseCourses. Il repository delega le letture paginabili a CoursePageQueries e le letture delle sessioni a ChefCourseQueries: entrambe interrogano le funzioni SQL e restituiscono modelli di lettura, senza ricostruire una classe Corso persistente.}
\end{landscape}

\clearpage\begin{landscape}

\section{Comportamento di due funzionalità}

\subsection{Creazione e pubblicazione del corso}\begin{center}

% generated by Plantuml 1.2026.8
\definecolor{plantucolor0000}{RGB}{159,67,43}%
\definecolor{plantucolor0001}{RGB}{255,255,255}%
\definecolor{plantucolor0002}{RGB}{128,144,149}%
\definecolor{plantucolor0003}{RGB}{0,0,0}%
\definecolor{plantucolor0004}{RGB}{226,226,240}%
\definecolor{plantucolor0005}{RGB}{24,24,24}%
\definecolor{plantucolor0006}{RGB}{255,255,255}%
\definecolor{plantucolor0007}{RGB}{23,57,70}%
\begin{tikzpicture}[xscale=0.90,yscale=-0.50,font=\fontsize{9}{10.2}\selectfont,every node/.style={yshift=3.5pt}
,pstyle0/.style={color=plantucolor0000,line width=1.5pt}
,pstyle1/.style={color=plantucolor0001,line width=0.0pt}
,pstyle2/.style={color=plantucolor0002,line width=0.5pt,dash pattern=on 5.0pt off 5.0pt}
,pstyle3/.style={color=plantucolor0005,fill=plantucolor0004,line width=0.5pt}
,pstyle4/.style={color=plantucolor0005,line width=0.5pt}
,pstyle5/.style={color=plantucolor0007,fill=white,line width=0.5pt}
,pstyle6/.style={color=plantucolor0007,fill=plantucolor0007,line width=1.0pt}
,pstyle7/.style={color=plantucolor0007,line width=1.0pt}
,pstyle8/.style={color=plantucolor0000,fill=white,line width=1.5pt}
,pstyle9/.style={color=plantucolor0007,line width=1.0pt,dash pattern=on 2.0pt off 2.0pt}
]
\draw[pstyle0] (97.24pt,262pt) rectangle (691.29pt,604pt);
\draw[pstyle0] (116.24pt,360pt) rectangle (666.29pt,590pt);
\draw[pstyle1] (20pt,80pt) rectangle (28pt,628pt);
\draw[pstyle2] (23.5pt,80pt) -- (23.5pt,628pt);
\draw[pstyle1] (128.74pt,80pt) rectangle (136.74pt,628pt);
\draw[pstyle2] (132.24pt,80pt) -- (132.24pt,628pt);
\draw[pstyle1] (287.39pt,80pt) rectangle (295.39pt,628pt);
\draw[pstyle2] (290.89pt,80pt) -- (290.89pt,628pt);
\draw[pstyle1] (403.37pt,80pt) rectangle (411.37pt,628pt);
\draw[pstyle2] (406.87pt,80pt) -- (406.87pt,628pt);
\draw[pstyle1] (521.23pt,80pt) rectangle (529.23pt,628pt);
\draw[pstyle2] (524.73pt,80pt) -- (524.73pt,628pt);
\draw[pstyle1] (646.79pt,80pt) rectangle (654.79pt,628pt);
\draw[pstyle2] (650.29pt,80pt) -- (650.29pt,628pt);
\node at (10.36pt,71.475pt)[below right,color=black,inner sep=0,yshift=-3.5pt]{Chef};
\draw[pstyle3] (23.5pt,18.5pt) ellipse (8pt and 8pt);
\draw[pstyle4] (23.5pt,26.5pt) -- (23.5pt,53.5pt)(10.5pt,34.5pt) -- (36.5pt,34.5pt)(23.5pt,53.5pt) -- (10.5pt,68.5pt)(23.5pt,53.5pt) -- (36.5pt,68.5pt);
\draw[pstyle5] (113.765pt,55pt) rectangle (150.715pt,79pt);
\node at (120.765pt,63.6pt)[below right,color=plantucolor0007,inner sep=0]{Form};
\draw[pstyle5] (255.105pt,55pt) rectangle (326.675pt,79pt);
\node at (262.105pt,63.475pt)[below right,color=plantucolor0007,inner sep=0]{Coordinatore};
\draw[pstyle5] (370.375pt,55pt) rectangle (443.365pt,79pt);
\node at (377.375pt,63.475pt)[below right,color=plantucolor0007,inner sep=0]{CreateCourse};
\draw[pstyle5] (494.005pt,55pt) rectangle (555.455pt,79pt);
\node at (501.005pt,63.585pt)[below right,color=plantucolor0007,inner sep=0]{Repository};
\node at (621.135pt,71.475pt)[below right,color=black,inner sep=0,yshift=-3.5pt]{PostgreSQL};
\draw[pstyle3] (632.29pt,34pt) ..controls (632.29pt,24pt) and (650.29pt,24pt) .. (650.29pt,24pt) ..controls (650.29pt,24pt) and (668.29pt,24pt) .. (668.29pt,34pt) -- (668.29pt,60pt) ..controls (668.29pt,70pt) and (650.29pt,70pt) .. (650.29pt,70pt) ..controls (650.29pt,70pt) and (632.29pt,70pt) .. (632.29pt,60pt) -- (632.29pt,34pt);
\draw[pstyle4] (632.29pt,34pt) ..controls (632.29pt,44pt) and (650.29pt,44pt) .. (650.29pt,44pt) ..controls (650.29pt,44pt) and (668.29pt,44pt) .. (668.29pt,34pt);
\draw[pstyle6] (120.24pt,100pt) -- (130.24pt,104pt) -- (120.24pt,108pt) -- (124.24pt,104pt) -- cycle;
\draw[pstyle7] (23.5pt,104pt) -- (126.24pt,104pt);
\node at (35.5pt,93.475pt)[below right,color=black,inner sep=0]{compila e conferma};
\draw[pstyle7] (132.24pt,128pt) -- (174.24pt,128pt);
\draw[pstyle7] (174.24pt,128pt) -- (174.24pt,141pt);
\draw[pstyle7] (133.24pt,141pt) -- (174.24pt,141pt);
\draw[pstyle6] (143.24pt,137pt) -- (133.24pt,141pt) -- (143.24pt,145pt) -- (139.24pt,141pt) -- cycle;
\node at (139.24pt,117.53pt)[below right,color=black,inner sep=0]{valida l'input};
\draw[pstyle6] (278.89pt,161pt) -- (288.89pt,165pt) -- (278.89pt,169pt) -- (282.89pt,165pt) -- cycle;
\draw[pstyle7] (132.24pt,165pt) -- (284.89pt,165pt);
\node at (182.625pt,154.25pt)[below right,color=black,inner sep=0]{runAsync(...)};
\draw[pstyle7] (404.87pt,189pt) -- (394.87pt,185pt);
\draw[pstyle7] (404.87pt,189pt) -- (394.87pt,193pt);
\draw[pstyle7] (290.89pt,189pt) -- (405.87pt,189pt);
\node at (302.89pt,178.25pt)[below right,color=black,inner sep=0]{execute(actor, input)};
\draw[pstyle7] (406.87pt,213pt) -- (448.87pt,213pt);
\draw[pstyle7] (448.87pt,213pt) -- (448.87pt,226pt);
\draw[pstyle7] (407.87pt,226pt) -- (448.87pt,226pt);
\draw[pstyle6] (417.87pt,222pt) -- (407.87pt,226pt) -- (417.87pt,230pt) -- (413.87pt,226pt) -- cycle;
\node at (413.87pt,202.53pt)[below right,color=black,inner sep=0]{valida il dominio};
\draw[pstyle6] (512.73pt,246pt) -- (522.73pt,250pt) -- (512.73pt,254pt) -- (516.73pt,250pt) -- cycle;
\draw[pstyle7] (406.87pt,250pt) -- (518.73pt,250pt);
\node at (418.87pt,239.25pt)[below right,color=black,inner sep=0]{createAndPublish(...)};
\draw[pstyle8] (97.24pt,262pt) -- (199.27pt,262pt) -- (199.27pt,264pt) -- (189.27pt,282pt) -- (97.24pt,282pt) -- (97.24pt,262pt);
\draw[pstyle0] (97.24pt,262pt) rectangle (691.29pt,604pt);
\node at (112.24pt,264.575pt)[below right,color=plantucolor0000,inner sep=0,yshift=-3.5pt]{\textbf{transazione}};
\draw[pstyle6] (638.29pt,296pt) -- (648.29pt,300pt) -- (638.29pt,304pt) -- (642.29pt,300pt) -- cycle;
\draw[pstyle7] (524.73pt,300pt) -- (644.29pt,300pt);
\node at (539.535pt,289.25pt)[below right,color=black,inner sep=0]{app.create\_course(...)};
\draw[pstyle6] (535.73pt,320pt) -- (525.73pt,324pt) -- (535.73pt,328pt) -- (531.73pt,324pt) -- cycle;
\draw[pstyle9] (529.73pt,324pt) -- (649.29pt,324pt);
\node at (565.94pt,313.53pt)[below right,color=black,inner sep=0]{course\_id};
\draw[pstyle6] (638.29pt,344pt) -- (648.29pt,348pt) -- (638.29pt,352pt) -- (642.29pt,348pt) -- cycle;
\draw[pstyle7] (524.73pt,348pt) -- (644.29pt,348pt);
\node at (536.73pt,337.25pt)[below right,color=black,inner sep=0]{app.publish\_course(...)};
\draw[pstyle8] (116.24pt,360pt) -- (174.49pt,360pt) -- (174.49pt,362pt) -- (164.49pt,380pt) -- (116.24pt,380pt) -- (116.24pt,360pt);
\draw[pstyle0] (116.24pt,360pt) rectangle (666.29pt,590pt);
\node at (131.24pt,362.53pt)[below right,color=plantucolor0000,inner sep=0,yshift=-3.5pt]{\textbf{alt}};
\node at (189.49pt,363.25pt)[below right,color=black,inner sep=0,yshift=-3.5pt]{\textbf{[successo]}};
\draw[pstyle6] (638.29pt,394pt) -- (648.29pt,398pt) -- (638.29pt,402pt) -- (642.29pt,398pt) -- cycle;
\draw[pstyle7] (524.73pt,398pt) -- (644.29pt,398pt);
\node at (571.125pt,387.715pt)[below right,color=black,inner sep=0]{commit};
\draw[pstyle6] (417.87pt,418pt) -- (407.87pt,422pt) -- (417.87pt,426pt) -- (413.87pt,422pt) -- cycle;
\draw[pstyle9] (411.87pt,422pt) -- (523.73pt,422pt);
\node at (447.42pt,411.53pt)[below right,color=black,inner sep=0]{risultato};
\draw[pstyle6] (301.89pt,442pt) -- (291.89pt,446pt) -- (301.89pt,450pt) -- (297.89pt,446pt) -- cycle;
\draw[pstyle9] (295.89pt,446pt) -- (405.87pt,446pt);
\node at (318.83pt,435.53pt)[below right,color=black,inner sep=0]{esito del Task};
\draw[pstyle6] (143.24pt,466pt) -- (133.24pt,470pt) -- (143.24pt,474pt) -- (139.24pt,470pt) -- cycle;
\draw[pstyle9] (137.24pt,470pt) -- (289.89pt,470pt);
\node at (148.665pt,459.53pt)[below right,color=black,inner sep=0]{callback FX: esito e rilettura};
\draw[color=plantucolor0000,line width=1.0pt,dash pattern=on 2.0pt off 2.0pt] (116.24pt,479pt) -- (666.29pt,479pt);
\node at (121.24pt,482.25pt)[below right,color=black,inner sep=0,yshift=-3.5pt]{\textbf{[errore SQL]}};
\draw[pstyle6] (638.29pt,506pt) -- (648.29pt,510pt) -- (638.29pt,514pt) -- (642.29pt,510pt) -- cycle;
\draw[pstyle7] (524.73pt,510pt) -- (644.29pt,510pt);
\node at (570.27pt,499.53pt)[below right,color=black,inner sep=0]{rollback};
\draw[pstyle6] (417.87pt,530pt) -- (407.87pt,534pt) -- (417.87pt,538pt) -- (413.87pt,534pt) -- cycle;
\draw[pstyle9] (411.87pt,534pt) -- (523.73pt,534pt);
\node at (445.81pt,523.715pt)[below right,color=black,inner sep=0]{eccezione};
\draw[pstyle6] (301.89pt,554pt) -- (291.89pt,558pt) -- (301.89pt,562pt) -- (297.89pt,558pt) -- cycle;
\draw[pstyle9] (295.89pt,558pt) -- (405.87pt,558pt);
\node at (316.035pt,547.53pt)[below right,color=black,inner sep=0]{errore del Task};
\draw[pstyle6] (143.24pt,578pt) -- (133.24pt,582pt) -- (143.24pt,586pt) -- (139.24pt,582pt) -- cycle;
\draw[pstyle9] (137.24pt,582pt) -- (289.89pt,582pt);
\node at (144.24pt,571.53pt)[below right,color=black,inner sep=0]{callback FX: conserva la bozza};
\end{tikzpicture}%

\captionof{figure}{Creazione e pubblicazione del corso}\label{fig:diagram-12}
\end{center}

{\small Nella \cref{fig:diagram-12}, Form indica CourseCreationForm, Coordinatore indica CourseWorkspaceCoordinator con FxAsyncExecutor e Task, Repository indica JdbcCourseRepository. Il ramo di rollback è alternativo a quello di successo. Un input non valido termina prima della chiamata al repository. Il form chiama CourseWorkspaceCoordinator attraverso CourseWorkspaceActions; il coordinatore usa FxAsyncExecutor. La bozza di input non equivale a un corso DRAFT già salvato.}

\end{landscape}

\clearpage\begin{landscape}

\subsection{Conferma di adesione alla pratica}\begin{center}

% generated by Plantuml 1.2026.8
\definecolor{plantucolor0000}{RGB}{159,67,43}%
\definecolor{plantucolor0001}{RGB}{255,255,255}%
\definecolor{plantucolor0002}{RGB}{128,144,149}%
\definecolor{plantucolor0003}{RGB}{0,0,0}%
\definecolor{plantucolor0004}{RGB}{226,226,240}%
\definecolor{plantucolor0005}{RGB}{24,24,24}%
\definecolor{plantucolor0006}{RGB}{255,255,255}%
\definecolor{plantucolor0007}{RGB}{23,57,70}%
\begin{tikzpicture}[xscale=0.82,yscale=-0.54,font=\fontsize{9}{10.2}\selectfont,every node/.style={yshift=3.5pt}
,pstyle0/.style={color=plantucolor0000,line width=1.5pt}
,pstyle1/.style={color=plantucolor0001,line width=0.0pt}
,pstyle2/.style={color=plantucolor0002,line width=0.5pt,dash pattern=on 5.0pt off 5.0pt}
,pstyle3/.style={color=plantucolor0005,fill=plantucolor0004,line width=0.5pt}
,pstyle4/.style={color=plantucolor0005,line width=0.5pt}
,pstyle5/.style={color=plantucolor0007,fill=white,line width=0.5pt}
,pstyle6/.style={color=plantucolor0007,fill=plantucolor0007,line width=1.0pt}
,pstyle7/.style={color=plantucolor0007,line width=1.0pt}
,pstyle8/.style={color=plantucolor0000,fill=white,line width=1.5pt}
,pstyle9/.style={color=plantucolor0007,line width=1.0pt,dash pattern=on 2.0pt off 2.0pt}
]
\draw[pstyle0] (100.485pt,225pt) rectangle (849.255pt,556pt);
\draw[pstyle0] (119.485pt,312pt) rectangle (766.105pt,542pt);
\draw[pstyle1] (37.085pt,80pt) rectangle (45.085pt,580pt);
\draw[pstyle2] (40.585pt,80pt) -- (40.585pt,580pt);
\draw[pstyle1] (131.985pt,80pt) rectangle (139.985pt,580pt);
\draw[pstyle2] (135.485pt,80pt) -- (135.485pt,580pt);
\draw[pstyle1] (285.335pt,80pt) rectangle (293.335pt,580pt);
\draw[pstyle2] (288.835pt,80pt) -- (288.835pt,580pt);
\draw[pstyle1] (417.055pt,80pt) rectangle (425.055pt,580pt);
\draw[pstyle2] (420.555pt,80pt) -- (420.555pt,580pt);
\draw[pstyle1] (574.725pt,80pt) rectangle (582.725pt,580pt);
\draw[pstyle2] (578.225pt,80pt) -- (578.225pt,580pt);
\draw[pstyle1] (746.605pt,80pt) rectangle (754.605pt,580pt);
\draw[pstyle2] (750.105pt,80pt) -- (750.105pt,580pt);
\node at (10pt,71.585pt)[below right,color=black,inner sep=0,yshift=-3.5pt]{Partecipante};
\draw[pstyle3] (40.585pt,18.5pt) ellipse (8pt and 8pt);
\draw[pstyle4] (40.585pt,26.5pt) -- (40.585pt,53.5pt)(27.585pt,34.5pt) -- (53.585pt,34.5pt)(40.585pt,53.5pt) -- (27.585pt,68.5pt)(40.585pt,53.5pt) -- (53.585pt,68.5pt);
\draw[pstyle5] (105.555pt,55pt) rectangle (165.415pt,79pt);
\node at (112.555pt,63.475pt)[below right,color=plantucolor0007,inner sep=0]{Interfaccia};
\draw[pstyle5] (272.18pt,55pt) rectangle (305.49pt,79pt);
\node at (279.18pt,63.585pt)[below right,color=plantucolor0007,inner sep=0]{Dati};
\draw[pstyle5] (394.915pt,55pt) rectangle (446.195pt,79pt);
\node at (401.915pt,63.42pt)[below right,color=plantucolor0007,inner sep=0]{Adesioni};
\draw[pstyle5] (547.5pt,55pt) rectangle (608.95pt,79pt);
\node at (554.5pt,63.585pt)[below right,color=plantucolor0007,inner sep=0]{Repository};
\node at (720.95pt,71.475pt)[below right,color=black,inner sep=0,yshift=-3.5pt]{PostgreSQL};
\draw[pstyle3] (732.105pt,34pt) ..controls (732.105pt,24pt) and (750.105pt,24pt) .. (750.105pt,24pt) ..controls (750.105pt,24pt) and (768.105pt,24pt) .. (768.105pt,34pt) -- (768.105pt,60pt) ..controls (768.105pt,70pt) and (750.105pt,70pt) .. (750.105pt,70pt) ..controls (750.105pt,70pt) and (732.105pt,70pt) .. (732.105pt,60pt) -- (732.105pt,34pt);
\draw[pstyle4] (732.105pt,34pt) ..controls (732.105pt,44pt) and (750.105pt,44pt) .. (750.105pt,44pt) ..controls (750.105pt,44pt) and (768.105pt,44pt) .. (768.105pt,34pt);
\draw[pstyle6] (123.485pt,100pt) -- (133.485pt,104pt) -- (123.485pt,108pt) -- (127.485pt,104pt) -- cycle;
\draw[pstyle7] (40.585pt,104pt) -- (129.485pt,104pt);
\node at (52.585pt,93.475pt)[below right,color=black,inner sep=0]{sceglie conferma};
\draw[pstyle7] (135.485pt,128pt) -- (177.485pt,128pt);
\draw[pstyle7] (177.485pt,128pt) -- (177.485pt,141pt);
\draw[pstyle7] (136.485pt,141pt) -- (177.485pt,141pt);
\draw[pstyle6] (146.485pt,137pt) -- (136.485pt,141pt) -- (146.485pt,145pt) -- (142.485pt,141pt) -- cycle;
\node at (142.485pt,117.475pt)[below right,color=black,inner sep=0]{verifica disponibilità e dialogo};
\draw[pstyle7] (286.835pt,165pt) -- (276.835pt,161pt);
\draw[pstyle7] (286.835pt,165pt) -- (276.835pt,169pt);
\draw[pstyle7] (135.485pt,165pt) -- (287.835pt,165pt);
\node at (172.96pt,154.25pt)[below right,color=black,inner sep=0]{write: confirm(...)};
\draw[pstyle6] (408.555pt,185pt) -- (418.555pt,189pt) -- (408.555pt,193pt) -- (412.555pt,189pt) -- cycle;
\draw[pstyle7] (288.835pt,189pt) -- (414.555pt,189pt);
\node at (300.835pt,178.25pt)[below right,color=black,inner sep=0]{confirm(actor, sessionId)};
\draw[pstyle6] (566.225pt,209pt) -- (576.225pt,213pt) -- (566.225pt,217pt) -- (570.225pt,213pt) -- cycle;
\draw[pstyle7] (420.555pt,213pt) -- (572.225pt,213pt);
\node at (432.555pt,202.25pt)[below right,color=black,inner sep=0]{setAttendance(CONFIRMED)};
\draw[pstyle8] (100.485pt,225pt) -- (202.515pt,225pt) -- (202.515pt,227pt) -- (192.515pt,245pt) -- (100.485pt,245pt) -- (100.485pt,225pt);
\draw[pstyle0] (100.485pt,225pt) rectangle (849.255pt,556pt);
\node at (115.485pt,227.575pt)[below right,color=plantucolor0000,inner sep=0,yshift=-3.5pt]{\textbf{transazione}};
\draw[pstyle6] (738.105pt,259pt) -- (748.105pt,263pt) -- (738.105pt,267pt) -- (742.105pt,263pt) -- cycle;
\draw[pstyle7] (578.225pt,263pt) -- (744.105pt,263pt);
\node at (590.225pt,252.25pt)[below right,color=black,inner sep=0]{app.set\_practical\_attendance(...)};
\draw[pstyle7] (750.105pt,287pt) -- (792.105pt,287pt);
\draw[pstyle7] (792.105pt,287pt) -- (792.105pt,300pt);
\draw[pstyle7] (751.105pt,300pt) -- (792.105pt,300pt);
\draw[pstyle6] (761.105pt,296pt) -- (751.105pt,300pt) -- (761.105pt,304pt) -- (757.105pt,300pt) -- cycle;
\node at (757.105pt,276.53pt)[below right,color=black,inner sep=0]{controlla i vincoli};
\draw[pstyle8] (119.485pt,312pt) -- (177.735pt,312pt) -- (177.735pt,314pt) -- (167.735pt,332pt) -- (119.485pt,332pt) -- (119.485pt,312pt);
\draw[pstyle0] (119.485pt,312pt) rectangle (766.105pt,542pt);
\node at (134.485pt,314.53pt)[below right,color=plantucolor0000,inner sep=0,yshift=-3.5pt]{\textbf{alt}};
\node at (192.735pt,315.25pt)[below right,color=black,inner sep=0,yshift=-3.5pt]{\textbf{[successo]}};
\draw[pstyle6] (738.105pt,346pt) -- (748.105pt,350pt) -- (738.105pt,354pt) -- (742.105pt,350pt) -- cycle;
\draw[pstyle7] (578.225pt,350pt) -- (744.105pt,350pt);
\node at (647.78pt,339.715pt)[below right,color=black,inner sep=0]{commit};
\draw[pstyle6] (431.555pt,370pt) -- (421.555pt,374pt) -- (431.555pt,378pt) -- (427.555pt,374pt) -- cycle;
\draw[pstyle9] (425.555pt,374pt) -- (577.225pt,374pt);
\node at (489.365pt,363.715pt)[below right,color=black,inner sep=0]{esito};
\draw[pstyle6] (299.835pt,394pt) -- (289.835pt,398pt) -- (299.835pt,402pt) -- (295.835pt,398pt) -- cycle;
\draw[pstyle9] (293.835pt,398pt) -- (419.555pt,398pt);
\node at (344.67pt,387.715pt)[below right,color=black,inner sep=0]{esito};
\draw[pstyle6] (146.485pt,418pt) -- (136.485pt,422pt) -- (146.485pt,426pt) -- (142.485pt,422pt) -- cycle;
\draw[pstyle9] (140.485pt,422pt) -- (287.835pt,422pt);
\node at (147.485pt,411.53pt)[below right,color=black,inner sep=0]{callback FX: cache e rilettura};
\draw[color=plantucolor0000,line width=1.0pt,dash pattern=on 2.0pt off 2.0pt] (119.485pt,431pt) -- (766.105pt,431pt);
\node at (124.485pt,434.25pt)[below right,color=black,inner sep=0,yshift=-3.5pt]{\textbf{[errore SQL]}};
\draw[pstyle6] (738.105pt,458pt) -- (748.105pt,462pt) -- (738.105pt,466pt) -- (742.105pt,462pt) -- cycle;
\draw[pstyle7] (578.225pt,462pt) -- (744.105pt,462pt);
\node at (646.925pt,451.53pt)[below right,color=black,inner sep=0]{rollback};
\draw[pstyle6] (431.555pt,482pt) -- (421.555pt,486pt) -- (431.555pt,490pt) -- (427.555pt,486pt) -- cycle;
\draw[pstyle9] (425.555pt,486pt) -- (577.225pt,486pt);
\node at (479.4pt,475.715pt)[below right,color=black,inner sep=0]{eccezione};
\draw[pstyle6] (299.835pt,506pt) -- (289.835pt,510pt) -- (299.835pt,514pt) -- (295.835pt,510pt) -- cycle;
\draw[pstyle9] (293.835pt,510pt) -- (419.555pt,510pt);
\node at (341.875pt,500.76pt)[below right,color=black,inner sep=0]{errore};
\draw[pstyle6] (146.485pt,530pt) -- (136.485pt,534pt) -- (146.485pt,538pt) -- (142.485pt,534pt) -- cycle;
\draw[pstyle9] (140.485pt,534pt) -- (287.835pt,534pt);
\node at (149.82pt,523.53pt)[below right,color=black,inner sep=0]{callback FX: messaggio utile};
\end{tikzpicture}%

\captionof{figure}{Conferma di adesione alla pratica}\label{fig:diagram-13}
\end{center}

{\small Nella \cref{fig:diagram-13}, Interfaccia raggruppa ParticipantWorkspaceController e ParticipantWorkspaceAsync; Dati indica ParticipantWorkspaceData, Adesioni indica ManagePracticalAttendance e Repository indica JdbcEnrollmentRepository. Il raggruppamento mantiene visibile il passaggio asincrono senza introdurre una nuova classe. Se il partecipante annulla il dialogo, il controller termina senza chiamare il caso d'uso. La risposta di lettura obsoleta viene scartata tramite versione del contesto. Una scrittura riuscita e una successiva rilettura fallita sono esiti distinti: il secondo non prova che il salvataggio sia fallito. Le scritture non vengono ripetute automaticamente dall'helper asincrono.}

\end{landscape}

\clearpage\section{Realizzazione della persistenza}

Lo schema app utilizza ruoli distinti. Il migratore può modificare lo schema; il ruolo applicativo dispone dei privilegi necessari alle operazioni esposte. La migrazione iniziale V001 definisce tabelle, vincoli, indici, funzioni, trigger e privilegi; V002 aggiorna il controllo temporale dello spostamento delle sessioni. Flyway applica le migrazioni in ordine e controlla versione e checksum di ciascun file. Una correzione viene aggiunta senza modificare la migrazione già applicata, così anche i database esistenti ricevono lo stesso aggiornamento. I dati sintetici della demo e delle prove sono separati dalla struttura.

Il catalogo effettivo comprende 81 routine, tutte funzioni PostgreSQL, e 11 trigger applicativi. Non sono presenti procedure con prokind=p nel catalogo corrente: le operazioni sono implementate come funzioni SQL/PLpgSQL e trigger. Molte funzioni autorizzate usano SECURITY DEFINER con search\_path controllato. Il JDBC usa PreparedStatement e parametri per gli input.

Le operazioni di dominio passano normalmente attraverso l'API SQL. Alcune letture tecniche autorizzate, come quelle di schema\_metadata, accedono direttamente alle tabelle di supporto. Il percorso di accesso dipende quindi dalla natura dell'operazione.


\subsection{Struttura SQL delle tabelle: casi significativi}\label{sec:sql-tables}

La definizione completa dello schema è nel file del progetto
\sqlsource.
Il file contiene le 22 tabelle, i vincoli introdotti con ALTER TABLE, gli indici, le funzioni, i trigger e i privilegi. Lo schema logico e il dizionario degli attributi sono descritti nella \cref{sec:logical-schema}; qui sono approfondite tre scelte che richiedono una spiegazione ulteriore.

\subsubsection{Sessioni online e pratiche}

Una sola tabella rappresenta le due modalità perché data, durata e collegamento al corso sono comuni. I CHECK rendono invece obbligatori i dati specifici: una sessione online richiede un link HTTPS e non usa sede o capienza; una pratica richiede sede, capienza positiva e chiusura delle adesioni precedente all'inizio. Il \cref{lst:course-session} mostra questi vincoli nella definizione fisica. Le chiavi, aggiunte separatamente nella migrazione, conservano il collegamento al corso e l'unicità del numero di sessione al suo interno.

\par\noindent\begin{minipage}{\linewidth}
\begin{lstlisting}[language=SQL,caption={Dati comuni e vincoli delle due modalità di sessione},label={lst:course-session}]
CREATE TABLE app.course_session (
    id uuid DEFAULT gen_random_uuid() NOT NULL,
    course_id uuid NOT NULL,
    sequence_number integer NOT NULL,
    scheduled_at timestamp without time zone NOT NULL,
    mode character varying(20) NOT NULL,
    status character varying(20) DEFAULT 'SCHEDULED'::character varying NOT NULL,
    duration_minutes integer DEFAULT 120 NOT NULL,
    online_url text,
    venue character varying(180),
    capacity integer,
    attendance_closes_at timestamp without time zone,
    completed_at timestamp with time zone,
    CONSTRAINT course_session_attendance_cutoff_valid CHECK (((((mode)::text = 'ONLINE'::text) AND (attendance_closes_at IS NULL)) OR (((mode)::text = 'PRACTICAL'::text) AND (attendance_closes_at IS NOT NULL) AND (attendance_closes_at < scheduled_at)))),
    CONSTRAINT course_session_duration_valid CHECK (((duration_minutes >= 1) AND (duration_minutes <= 720))),
    CONSTRAINT course_session_mode_check CHECK (((mode)::text = ANY ((ARRAY['ONLINE'::character varying, 'PRACTICAL'::character varying])::text[]))),
    CONSTRAINT course_session_mode_fields_valid CHECK (((((mode)::text = 'ONLINE'::text) AND (online_url IS NOT NULL) AND (online_url ~* '^https://[[:alnum:]][[:alnum:].-]*(:[0-9]+)?([/?#][^[:space:]]*)?$'::text) AND (venue IS NULL) AND (capacity IS NULL)) OR (((mode)::text = 'PRACTICAL'::text) AND (online_url IS NULL) AND (venue IS NOT NULL) AND ((venue)::text = btrim((venue)::text)) AND (btrim((venue)::text) <> ''::text) AND (capacity IS NOT NULL) AND (capacity > 0)))),
    CONSTRAINT course_session_sequence_number_check CHECK ((sequence_number > 0)),
    CONSTRAINT course_session_status_check CHECK (((status)::text = ANY ((ARRAY['SCHEDULED'::character varying, 'COMPLETED'::character varying, 'CANCELLED'::character varying])::text[])))
);
\end{lstlisting}
\end{minipage}\par

\subsubsection{Dosi esatte per una persona}

La tabella associativa \code{recipe_ingredient} conserva la quantità della coppia ricetta/ingrediente. La chiave composta impedisce una dose duplicata per lo stesso ingrediente. La colonna NUMERIC conserva il valore inserito; il controllo nel \cref{lst:recipe-dose} ammette quantità positive inferiori a $10^{12}$ e con massimo tre decimali. Il confronto con round verifica la precisione e rifiuta il valore non valido, senza correggerlo mediante arrotondamento.

\par\noindent\begin{minipage}{\linewidth}
\begin{lstlisting}[language=SQL,caption={Associazione tra ricetta e ingrediente con controllo della precisione},label={lst:recipe-dose}]
CREATE TABLE app.recipe_ingredient (
    recipe_id uuid NOT NULL,
    ingredient_id uuid NOT NULL,
    quantity numeric NOT NULL,
    CONSTRAINT recipe_ingredient_quantity_check CHECK ((quantity > (0)::numeric)),
    CONSTRAINT recipe_quantity_precision CHECK (((quantity > (0)::numeric) AND (quantity < ('1000000000000'::bigint)::numeric) AND (quantity = round(quantity, 3))))
);
\end{lstlisting}
\end{minipage}\par

\subsubsection{Copie storiche del menu}

La copia della ricetta conserva insieme nome, descrizione, unità e dosi. JSONB permette di acquisire questo insieme senza collegarlo ai valori successivi del ricettario. È una ridondanza intenzionale: il menu della sessione deve restare interpretabile anche dopo la modifica di una ricetta. Il \cref{lst:menu-copy} mostra struttura e provenienza della copia; la FK composta, definita nel file completo, la collega alla coppia di \code{session_recipe}.

La provenienza CAPTURED distingue una copia acquisita da una ricostruzione LEGACY\_RECONSTRUCTED. Una ricostruzione non prova quali fossero le dosi originarie. Prima della chiusura delle adesioni è consentito l'aggiornamento esplicito dal ricettario; le normali modifiche del ricettario non riscrivono automaticamente le copie.

\par\noindent\begin{minipage}{\linewidth}
\begin{lstlisting}[language=SQL,caption={Copia della ricetta associata alla sessione},label={lst:menu-copy}]
CREATE TABLE app.session_menu_snapshot (
    session_id uuid NOT NULL,
    recipe_id uuid NOT NULL,
    recipe_name text NOT NULL,
    description text NOT NULL,
    ingredients jsonb NOT NULL,
    captured_at timestamp with time zone DEFAULT clock_timestamp() NOT NULL,
    provenance text NOT NULL,
    CONSTRAINT session_menu_snapshot_provenance_check CHECK ((provenance = ANY (ARRAY['CAPTURED'::text, 'LEGACY_RECONSTRUCTED'::text])))
);
\end{lstlisting}
\end{minipage}\par

\subsection{Trigger e funzioni SQL: operazioni significative}\label{sec:sql-operations}

PK, FK, UNIQUE e CHECK proteggono identità, riferimenti e condizioni sulla singola riga. Le regole che coinvolgono più righe richiedono invece controlli coordinati. Le operazioni applicative passano attraverso funzioni che verificano anche token, ruolo e proprietà dei dati; i trigger mantengono i vincoli sulle modifiche delle tabelle.

La realizzazione corrente comprende funzioni SQL/PLpgSQL e trigger, senza procedure SQL. Le dichiarazioni e i corpi completi della struttura iniziale sono disponibili in \sqlsource; la correzione successiva è in \sqlcorrection. Entrambi i file appartengono alla directory delle migrazioni e definiscono insieme il comportamento corrente. Sono riportati soltanto i casi in cui la copia storica, la concorrenza o il calcolo dei risultati meritano attenzione.

La funzione \code{reschedule_course_session} blocca una modifica quando l'inizio originale è già stato raggiunto, confrontando l'appuntamento con l'ora dell'operazione nel fuso Europe/Rome. V002 applica questo controllo anche quando il parametro della sede o del link è assente: conservarne il valore non deve consentire di spostare una sessione già iniziata. La funzione esistente viene sostituita con \code{CREATE OR REPLACE}, mantenendo firma, autorizzazioni e transazione; tabelle, vincoli e numero delle routine restano invariati.

\subsubsection{Acquisizione automatica della ricetta}

Il trigger del \cref{lst:capture-trigger} interviene dopo l'associazione della ricetta alla pratica. La funzione del \cref{lst:capture-menu} acquisisce nome e descrizione, costruisce l'elenco JSONB degli ingredienti e registra la provenienza. L'ordinamento per identificatore rende stabile la rappresentazione. La copia e l'associazione appartengono alla stessa transazione: un errore nell'acquisizione impedisce di salvare un'associazione ordinaria priva della propria copia.

\par\noindent\begin{minipage}{\linewidth}
\begin{lstlisting}[language=SQL,caption={Trigger di acquisizione della copia dopo l’associazione},label={lst:capture-trigger}]
CREATE TRIGGER capture_session_recipe AFTER INSERT ON app.session_recipe FOR EACH ROW EXECUTE FUNCTION app.capture_session_recipe();
\end{lstlisting}
\end{minipage}\par

\par\noindent\begin{minipage}{\linewidth}
\begin{lstlisting}[language=SQL,caption={Acquisizione della copia del menu},label={lst:capture-menu}]
CREATE FUNCTION app.capture_session_recipe() RETURNS trigger
    LANGUAGE plpgsql
    SET search_path TO 'pg_catalog', 'app', 'pg_temp'
    AS $$
BEGIN
    INSERT INTO app.session_menu_snapshot(session_id, recipe_id, recipe_name, description, ingredients, provenance)
    SELECT NEW.session_id, r.id, r.name, r.description,
        coalesce((SELECT jsonb_agg(jsonb_build_object('ingredient_id', i.id,
            'ingredient_name', i.name, 'unit', i.unit, 'quantity', ri.quantity) ORDER BY i.id)
            FROM app.recipe_ingredient ri JOIN app.ingredient i ON i.id=ri.ingredient_id
            WHERE ri.recipe_id=r.id), '[]'::jsonb), 'CAPTURED'
    FROM app.recipe r WHERE r.id=NEW.recipe_id;
    RETURN NEW;
END; $$;
\end{lstlisting}
\end{minipage}\par

\subsubsection{Capienza e validità delle adesioni}

Due conferme simultanee non devono occupare lo stesso ultimo posto. La funzione \code{enforce_practical_attendance}, richiamata dal trigger del \cref{lst:attendance-trigger}, blocca le righe di iscrizione e sessione con FOR UPDATE prima di verificare coerenza e capienza. Il conteggio considera soltanto iscrizioni attive e risposte del periodo corrente: una reiscrizione non riattiva una conferma precedente.

Il \cref{lst:attendance-capacity} è un estratto del ramo per le conferme, successivo all'acquisizione dei blocchi. La funzione completa controlla anche che la pratica appartenga al corso dell'iscrizione e aggiorna il momento della risposta. La funzione di accesso \code{set_practical_attendance} verifica inoltre scadenza delle adesioni e presenza del menu; questi controlli non sono attribuiti al solo trigger. Entrambe le funzioni sono nel file SQL completo.

\par\noindent\begin{minipage}{\linewidth}
\begin{lstlisting}[language=SQL,caption={Trigger di protezione delle adesioni alla pratica},label={lst:attendance-trigger}]
CREATE TRIGGER practical_attendance_guard BEFORE INSERT OR UPDATE OF enrollment_id, session_id, status ON app.practical_attendance FOR EACH ROW EXECUTE FUNCTION app.enforce_practical_attendance();
\end{lstlisting}
\end{minipage}\par

\par\noindent\begin{minipage}{\linewidth}
\begin{lstlisting}[language=SQL,caption={Estratto: conferme valide e controllo della capienza dopo il blocco delle righe},label={lst:attendance-capacity}]
    IF NEW.status = 'CONFIRMED' THEN
        IF enrollment_status <> 'ACTIVE'
                OR session_status <> 'SCHEDULED'
                OR session_time <= (current_timestamp AT TIME ZONE 'Europe/Rome') THEN
            RAISE EXCEPTION 'Attendance requires an active enrollment and a future session'
                USING ERRCODE = '23514';
        END IF;

        SELECT count(*)::INTEGER
          INTO confirmed_count
          FROM app.practical_attendance
          JOIN app.enrollment current_enrollment ON current_enrollment.id = practical_attendance.enrollment_id
         WHERE practical_attendance.session_id = NEW.session_id
           AND practical_attendance.status = 'CONFIRMED'
           AND current_enrollment.status = 'ACTIVE'
           AND practical_attendance.updated_at >= current_enrollment.enrolled_at
           AND practical_attendance.enrollment_id <> NEW.enrollment_id;
        IF confirmed_count >= session_capacity THEN
            RAISE EXCEPTION 'Practical session capacity has been reached'
                USING ERRCODE = '23514';
        END IF;
    END IF;
\end{lstlisting}
\end{minipage}\par

\subsubsection{Appuntamenti dello stesso chef}

Un lock consultivo di transazione, calcolato sullo chef, coordina le modifiche concorrenti al suo calendario. La funzione del \cref{lst:chef-overlap} confronta gli intervalli delle sessioni programmate ed esclude la sessione in modifica. Le disuguaglianze strette consentono appuntamenti consecutivi: una sessione può iniziare esattamente quando termina l'altra. Il trigger \code{chef_session_overlap_guard}, definito nel file completo, attiva il controllo all'inserimento e alla modifica dei dati dell'appuntamento.

\par\noindent\begin{minipage}{\linewidth}
\begin{lstlisting}[language=SQL,caption={Blocco per chef e verifica degli intervalli del calendario},label={lst:chef-overlap}]
CREATE FUNCTION app.prevent_chef_session_overlap() RETURNS trigger
    LANGUAGE plpgsql SECURITY DEFINER
    SET search_path TO 'pg_catalog', 'app', 'pg_temp'
    AS $$
DECLARE owner_id UUID;
BEGIN
    IF NEW.status <> 'SCHEDULED' THEN
        RETURN NEW;
    END IF;

    SELECT chef_profile_id INTO owner_id FROM app.course WHERE id = NEW.course_id;
    IF owner_id IS NULL THEN
        RAISE EXCEPTION 'Course owner is unavailable' USING ERRCODE = '23503';
    END IF;

    PERFORM pg_advisory_xact_lock(hashtextextended(owner_id::TEXT, 0));
    IF EXISTS (
        SELECT 1
          FROM app.course_session existing_session
          JOIN app.course existing_course ON existing_course.id = existing_session.course_id
         WHERE existing_course.chef_profile_id = owner_id
           AND existing_session.status = 'SCHEDULED'
           AND existing_session.id IS DISTINCT FROM NEW.id
           AND existing_session.scheduled_at
                 < NEW.scheduled_at + make_interval(mins => NEW.duration_minutes)
           AND NEW.scheduled_at
                 < existing_session.scheduled_at + make_interval(mins => existing_session.duration_minutes)
    ) THEN
        RAISE EXCEPTION 'Chef already has another session in this time interval'
            USING ERRCODE = '23P01';
    END IF;
    RETURN NEW;
END;
$$;
\end{lstlisting}
\end{minipage}\par

\subsubsection{Fabbisogni dalle copie e dalle risposte correnti}

La query del \cref{lst:requirements} separa le risposte degli iscritti dalle dosi del menu. Il collegamento temporale esclude le risposte anteriori all'iscrizione corrente. JSONB viene trasformato in righe; la somma per identificatore di ingrediente è poi moltiplicata per le conferme valide. Le quantità e le unità provengono dalla copia, mentre il nome mostrato è quello corrente del catalogo.

\code{preliminary_session_preparation} è una funzione interna di calcolo. La funzione applicativa \code{session_preparation}, presente nel file completo, verifica l'identità dello chef e la proprietà della pratica prima di richiamarla. La separazione permette di riusare il calcolo senza rendere pubblica una lettura priva di autorizzazione.

Per un ingrediente $i$ della pratica $s$, il fabbisogno è $Q_{s,i}=n_s\sum_{r\in M_s}q_{r,i}$: $n_s$ indica le conferme valide, $M_s$ il menu copiato e $q_{r,i}$ la dose per una persona. Con 6 conferme e dosi di 100 g e 50 g dello stesso ingrediente, il totale è 900 g.

\par\noindent\begin{minipage}{\linewidth}
\begin{lstlisting}[language=SQL,caption={Calcolo dei fabbisogni per ingrediente dalle copie del menu},label={lst:requirements}]
CREATE FUNCTION app.preliminary_session_preparation(p_session uuid) RETURNS TABLE(confirmed_participants integer, pending_participants integer, declined_participants integer, ingredient_margin numeric, recipe_id uuid, recipe_name text, base_servings integer, portions_per_participant numeric, ingredient_id uuid, ingredient_name text, unit text, contribution_quantity numeric, required_quantity numeric)
    LANGUAGE sql STABLE
    SET search_path TO 'pg_catalog', 'app', 'pg_temp'
    AS $$
    WITH responses AS (
        SELECT count(*) FILTER(WHERE pa.status='CONFIRMED')::INTEGER confirmed,
            count(*) FILTER(WHERE pa.status IS NULL)::INTEGER pending,
            count(*) FILTER(WHERE pa.status='DECLINED')::INTEGER declined
        FROM app.enrollment e JOIN app.course_session s ON s.course_id=e.course_id
        LEFT JOIN app.practical_attendance pa ON pa.enrollment_id=e.id AND pa.session_id=s.id AND pa.updated_at>=e.enrolled_at
        WHERE s.id=p_session AND e.status='ACTIVE'
    ), doses AS (
        SELECT m.recipe_id,m.recipe_name,i.* FROM app.session_menu_snapshot m
        CROSS JOIN LATERAL jsonb_to_recordset(m.ingredients)
            AS i(ingredient_id UUID,ingredient_name TEXT,unit TEXT,quantity NUMERIC)
        WHERE m.session_id=p_session
    )
    SELECT r.confirmed,r.pending,r.declined,0::NUMERIC,d.recipe_id,d.recipe_name,1,1::NUMERIC,
        d.ingredient_id,i.name::TEXT,d.unit,d.quantity*r.confirmed,
        sum(d.quantity) OVER(PARTITION BY d.ingredient_id)*r.confirmed
    FROM responses r LEFT JOIN doses d ON TRUE
    LEFT JOIN app.ingredient i ON i.id=d.ingredient_id;
$$;
\end{lstlisting}
\end{minipage}\par

\subsubsection{Svolgimento derivato e report mensile}

Lo svolgimento è ricavato dalla fine prevista e mantiene la precedenza dell'annullamento. Il \cref{lst:effective-status} rende il calcolo indipendente dall'apertura dell'applicazione o da un processo pianificato. L'orario di confronto è interpretato nel fuso Europe/Rome dalle funzioni chiamanti.

\par\noindent\begin{minipage}{\linewidth}
\begin{lstlisting}[language=SQL,caption={Stato effettivo derivato dalla fine prevista della sessione},label={lst:effective-status}]
CREATE FUNCTION app.effective_session_status(p_status text, p_scheduled_at timestamp without time zone, p_duration_minutes integer, p_as_of timestamp without time zone) RETURNS text
    LANGUAGE sql IMMUTABLE STRICT
    SET search_path TO 'pg_catalog', 'app', 'pg_temp'
    AS $$
    SELECT CASE WHEN p_status='CANCELLED' THEN 'CANCELLED'
        WHEN p_scheduled_at+make_interval(mins=>p_duration_minutes)<=p_as_of THEN 'COMPLETED'
        ELSE 'SCHEDULED' END;
$$;
\end{lstlisting}
\end{minipage}\par

Il \cref{lst:monthly-report} aggrega soltanto le sessioni svolte. Il numero di corsi è distinto; media, minimo e massimo delle ricette riguardano le sole pratiche. La funzione \code{chef_monthly_report_details}, disponibile nel file completo, controlla ruolo e mese richiesto, attribuisce le sessioni al mese della data programmata e conta le copie delle ricette associate. Si evita così di duplicare nel riepilogo la regola temporale e i controlli di accesso.

\par\noindent\begin{minipage}{\linewidth}
\begin{lstlisting}[language=SQL,caption={Indicatori del report mensile dello chef},label={lst:monthly-report}]
CREATE FUNCTION app.chef_monthly_report(p_session_token uuid, p_year integer, p_month integer) RETURNS TABLE(total_courses bigint, online_sessions bigint, practical_sessions bigint, average_recipes numeric, minimum_recipes integer, maximum_recipes integer)
    LANGUAGE sql STABLE SECURITY DEFINER
    SET search_path TO 'pg_catalog', 'app', 'pg_temp'
    AS $$
    SELECT count(DISTINCT s.course_id),count(*) FILTER(WHERE d.mode='ONLINE'),count(*) FILTER(WHERE d.mode='PRACTICAL'),
        avg(d.recipes) FILTER(WHERE d.mode='PRACTICAL'),min(d.recipes) FILTER(WHERE d.mode='PRACTICAL'),max(d.recipes) FILTER(WHERE d.mode='PRACTICAL')
    FROM app.chef_monthly_report_details(p_session_token,p_year,p_month) d JOIN app.course_session s ON s.id=d.session_id
    WHERE d.state='COMPLETED';
$$;
\end{lstlisting}
\end{minipage}\par

\subsection{Autenticazione e gestione delle connessioni}

Le password sono protette con PBKDF2 e sale casuale; \code{auth_session} conserva l'hash del token di accesso. Le routine verificano scadenza, account attivo e ruolo, oltre alla proprietà o visibilità dei dati richiesti. Questi controlli proteggono le operazioni anche oltre la validazione della GUI.

Migrazione e accesso applicativo usano credenziali separate per limitare i privilegi del client. Le query JDBC sono parametrizzate; le operazioni lente vengono eseguite fuori dal thread JavaFX e gli errori sono tradotti in messaggi comprensibili.

\end{document}
